\chapter{The classification of nice central-charge-one vertex operator algebras (Gannon--Rayhaun)}\label{ch:L22}
\begingroup\setlength{\emergencystretch}{3em}\sloppy

\providecommand{\ltwotwoS}{S_{00}}
\providecommand{\ltwotwoch}{\operatorname{ch}}
\providecommand{\ltwotwoH}{\mathbb{H}}
\providecommand{\ltwotwoSL}{\mathrm{SL}_2(\Z)}

This chapter was first written against an earlier state of the LeanVOA library \emph{together with the
working tree of 2026-10-13} (the book's snapshot). The files that prove the \emph{converse}
(Sections~\ref{l22:sec:B3} and~\ref{l22:sec:final}: the logarithmic tensor category, the McRae orbifold argument and
\leanpath{MainConverseFinal.lean}) were added \emph{after} the snapshot, in the commits that followed it
(2026-10-13, 23:45, through 2026-10-14); this revision of the chapter was
checked against the last of those commits. Checks that were run at an earlier state of the library and not re-run
for this edition are marked as such. Every Lean quotation was copied mechanically from the source files and is
machine-checked by \texttt{scripts/textbook/sync.py} to occur verbatim (up to whitespace and comments) in the source files of the release documented by this edition. The chapter covers
the directory \leanpath{VOA/\allowbreak Classification/} (about 1750 files) and the umbrella
\leanpath{VOA/\allowbreak Classification.lean}. The deep representation-theoretic inputs ---
Huang's modular tensor category, the Verlinde formula, the Huang--Kirillov--Lepowsky extension theory,
the Dong--Lin--Ng congruence theorem and Ng--Schauenburg --- live in \leanpath{VOA/\allowbreak Huang/} and
are the subject of Chapter~\ref{ch:L21}; here they appear as named theorems whose statements we quote
and whose proofs we only sketch.

We rely on earlier chapters and do not repeat them: Lean idioms (structures, \lean{Prop}-valued
definitions, universes, \lean{∃}/\lean{∀}, \lean{Nonempty}, \lean{≃*}, \lean{HasSum}, \lean{finrank})
in Chapter~\ref{ch:01}; vertex operator algebras, the conformal vector, \lean{centralCharge} and the
grading \lean{V.grading n} in Chapters~\ref{ch:L03} and~\ref{ch:L04}; automorphism groups \lean{Aut W}
and fixed-point subalgebras \lean{orbifoldVOA G} in Chapters~\ref{ch:03} and~\ref{ch:04}; lattice vertex
algebras \lean{latticeVOAGen} and the lift \lean{thetaAut} of $-1$ in Chapters~\ref{ch:L05}, \ref{ch:03}
and~\ref{ch:04}; ordinary modules \lean{OrdVModule}, irreducibility, \lean{ModuleEquiv},
\lean{adjointModule}, weight spaces, $C_2$-cofiniteness, admissible rationality, self-duality and
\lean{StronglyRational} in Chapter~\ref{ch:L12} (Sections~\ref{l12:sec:categories},
\ref{l12:sec:rationality} and~\ref{l12:sec:strong}); one-point functions, square-bracket weights and Zhu's
modularity theorem in Chapters~\ref{ch:L14} and~\ref{ch:L15}; the Virasoro Kac determinant in
Chapter~\ref{ch:L06}.

\paragraph{What this chapter is about, in one paragraph.} A two-dimensional conformal field theory
with central charge $c=1$ is, in the simplest cases, a single free boson. Its chiral algebra is a vertex
operator algebra of central charge $1$. Gannon and Rayhaun (arXiv:2609.00122, Theorem~3.1) proved that a
vertex operator algebra which is \emph{nice} --- strongly rational and pseudo-unitary --- and has $c=1$ is
isomorphic to exactly one of the following: the lattice VOA $V_{2m}$ of the rank-one lattice $\sqrt{2m}\,\Z$
($m\ge1$); its $\Z_2$-orbifold $V_{2m}^+$ ($m\ge2$); or one of the three \emph{exceptional orbifolds}
$V_2^{A_4}$, $V_2^{S_4}$, $V_2^{A_5}$ of $V_2 \cong L_{\mathfrak{sl}_2}(1,0)$ by the tetrahedral, octahedral and
icosahedral subgroups of $\Aut V_2\cong \mathrm{PSL}_2(\C)$. Conversely every VOA in this list is nice.
The library formalizes the three parts of this theorem: the \emph{forward classification} (nice and
$c=1$ implies in the list), \emph{pairwise distinctness} (no VOA carries two labels), and the
\emph{converse} (everything in the list is nice).

\paragraph{The findings of this chapter, in brief.} Read these first; each is explained where it occurs.
\begin{enumerate}
\item \textbf{The forward classification is proved with no hypothesis}
  (\lean{VOA.Classification.COne.gannonRayhaun_forward}, Theorem~\ref{l22:thm:forward}). Its
  \lean{#print axioms} is Lean's three standard axioms (re-probed at an earlier commit, not re-run for this edition; at the snapshot it also
  listed \lean{VOA.pbw}, then a project axiom). It is unconditional. It uses the PBW theorem \lean{VOA.pbw}
  (Chapter~\ref{ch:L02}) through the modular-data package \lean{iModData_holds} (Chapter~\ref{ch:L21}).
\item \textbf{Pairwise distinctness is proved with no hypothesis and without} \lean{VOA.pbw}
  (\lean{gannonRayhaunDistinct}, Theorem~\ref{l22:thm:distinct}), by computing the graded dimensions
  of the exceptional orbifolds in weights $4$ and $16$ exactly (Molien's formula).
\item \textbf{The converse is proved with no hypothesis, and so is the whole classification}
  (\lean{gannonRayhaunConverse_holds} and \lean{gannonRayhaun_classification}, Theorem~\ref{l22:thm:final}). The last
  input of the converse was \lean{FiniteV2OrbifoldRationalOfC2Fact}, the \emph{rationality of the exceptional orbifolds}: every $C_2$-cofinite fixed-point algebra
  $V_2^G$ of a finite group is rational. It is now a theorem, \lean{finiteV2OrbifoldRationalOfC2}, proved by formalizing
  McRae's orbifold theorem (arXiv:1910.13226) for $V=V_2$ in seven steps (Section~\ref{l22:sec:B3}): the braided
  tensor category of logarithmic $V_2^G$-modules, $V_2$ as a rigid commutative algebra in it, McRae's projections onto
  $g$-twisted parts, the monodromy argument making each twisted part an untwisted $V_2^{\langle g\rangle}$-module, the
  rationality of $V_2^{\langle g\rangle}\cong V_{2n^2}$, McRae's Lemma~4.21, and the passage from split surjections to
  rationality. The historical reduction \lean{gannonRayhaunConverse_of_B3} is kept (the \texttt{B3} in this and other
  Lean names is the project's internal label for the rationality of the exceptional orbifolds); its hypothesis is now proved.
  \lean{#print axioms gannonRayhaun_classification} is Lean's three axioms (re-probed at an earlier commit, not re-run for this edition; at the
  snapshot also \lean{VOA.pbw}). The proof of the converse uses \lean{VOA.pbw} through the logarithmic tensor category
  (every categorical step uses it, through the finiteness theorem proved in Step~1 of that argument); since
  \lean{VOA.pbw} is now a theorem, \textbf{the whole classification is unconditional}. The three parts are stated together in universe \lean{0} (Remark~\ref{l22:rem:universe}).
\item \textbf{``Rational'' is the library's \lean{AdmissiblyRational}, and ``universe'' matters.} The rationality
  conjunct of \lean{Nice} is \lean{AdmissiblyRational} (weak submodules of admissible modules have weak complements). Its
  relation to Dong--Li--Mason rationality, which Gannon--Rayhaun use, is \textbf{not settled in Lean}; only DLM-rational
  $\Rightarrow$ \lean{AdmissiblyRational} is known, on paper. So the forward theorem covers every Lean-nice VOA, a class
  that contains every DLM-nice VOA by the paper direction, and the converse is proved in the library's sense
  (Remark~\ref{l22:rem:rational}). The classification and the converse are stated for carriers in \lean{Type 0} only; the
  forward theorem and distinctness hold in every universe (Remark~\ref{l22:rem:universe}).
\item \textbf{The modular-data input had to be repaired once.} An earlier form of the modular-data
  package constrained the $S$-matrix only through characters; it was \emph{satisfied by a continuum of
  wrong $S$-matrices} (machine-checked counterexample \lean{isModularDataWeak_perturbS}). The current
  definition pins $S$ on all one-point functions (\lean{SPinned}); see Warning~\ref{l22:warn:spin}.
\item \textbf{Lean defects found while writing} (Section~\ref{l22:sec:defects}): stale docstrings that
  still say ``not proved here'' for statements now proved (\lean{PlusExceptionalResidual} and
  \lean{ExceptionalResidual} in \leanpath{Distinct.lean}, \lean{GannonRayhaunDistinct} and \lean{GannonRayhaunConverse} in
  \leanpath{Statement.lean}), and docstrings of the orbifold route that still call \lean{TwistedJacobiBridge} and
  \lean{TwistedJacobiBridgeAll} ``open'' although both are proved; no single
  declaration packages each input as a theorem of the form \lean{IVirasoro} or \lean{ISerreStark24}
  (they are discharged inside the assembly chain). None of these affects any statement.
\end{enumerate}

%======================================================================
\section{The theorem and its physics}\label{l22:sec:theorem}
%======================================================================

\subsection{The free boson and its rational points}

Let $\varphi(z)$ be a chiral free boson: a field with operator product expansion
$\partial\varphi(z)\,\partial\varphi(w)\sim -(z-w)^{-2}$. The current $J=i\partial\varphi$ generates the
Heisenberg vertex algebra $M(1)$ of Chapter~\ref{ch:L05}, of central charge $1$. Its irreducible modules
are Fock spaces $M(1,p)$ labelled by a momentum $p\in\C$, with lowest conformal weight $p^2/2$. A
\emph{compact} boson takes values in a circle; its momenta are then quantized, and if the radius is
rational in suitable units, finitely many momentum sectors close under fusion. Algebraically: for an even
lattice $L=\sqrt{2m}\,\Z$ (so that $\alpha^2/2 = m$ for the generator $\alpha$), the lattice VOA
\begin{equation*}
V_{2m} = V_{\sqrt{2m}\Z} = \bigoplus_{n\in\Z} M(1, n\sqrt{2m})
\end{equation*}
is a vertex operator algebra of central charge $1$ containing the vertex operators $e^{\pm\alpha}$ of
weight $m$. Its character is
\begin{equation}\label{l22:eq:latchar}
\ltwotwoch_{V_{2m}}(\tau) = \frac{\theta_{0,1,m}(\tau)}{\eta(\tau)}, \qquad
\theta_{0,1,m}(\tau)=\sum_{n\in\Z} q^{m n^2},\quad \eta(\tau) = q^{1/24}\prod_{k\ge1}(1-q^k),
\end{equation}
with $q=e^{2\pi i\tau}$. The automorphism $\theta$ of Chapter~\ref{ch:04}, the lift of $-1$ on $L$
($\varphi\mapsto-\varphi$), has fixed-point subalgebra $V_{2m}^+$. At $m=1$ the weight-one space of $V_2$ is
three-dimensional, spanned by $\alpha(-1)\mathbf 1$ and $e^{\pm\alpha}$, and forms $\mathfrak{sl}_2$; in fact
$V_2\cong L_{\mathfrak{sl}_2}(1,0)$, the $SU(2)$ level-one WZW chiral algebra, and $\Aut V_2 \cong
\mathrm{PSL}_2(\C)$ acts through the adjoint group. Its finite subgroups are those of $SO(3)$ (Klein):
cyclic, dihedral, tetrahedral $A_4$, octahedral $S_4$, icosahedral $A_5$.

\begin{physics}
The moduli space of $c=1$ unitary conformal field theories (Ginsparg 1988, Dijkgraaf--Verlinde--Verlinde
1988) consists of the \emph{circle line} (a compact boson of radius $R$), the \emph{orbifold line}
($S^1/\Z_2$), meeting at a point, and three \emph{isolated} theories obtained by orbifolding the $SU(2)_1$
point by the tetrahedral, octahedral and icosahedral groups. The Gannon--Rayhaun theorem is the
\emph{chiral} shadow of this picture: the rational points of the circle line have chiral algebra $V_{2m}$,
those of the orbifold line $V_{2m}^+$, and the three isolated theories have chiral algebras
$V_2^{A_4},V_2^{S_4},V_2^{A_5}$. The theorem says that, among chiral algebras with good representation theory
and positive energies, there is nothing else. This is a heuristic correspondence between full CFTs and
chiral algebras; the theorem itself is a statement about vertex operator algebras only.
\end{physics}

\subsection{Why the hypotheses}

\emph{Strongly rational} (Chapter~\ref{ch:L12}, Definition~\ref{l12:def:strong}) means: simple, of CFT type
(built into the library's \lean{VertexOperatorAlgebra}), self-dual, $C_2$-cofinite and rational. These are
exactly Huang's hypotheses for $\operatorname{Rep}(V)$ to be a modular tensor category, which is what makes
modular-form arguments possible. \emph{Pseudo-unitary} means that every irreducible module other than $V$
itself has strictly positive lowest conformal weight; in a unitary theory this is automatic (energies are
positive and only the vacuum has energy zero). Without it there are non-unitary rational $c=1$ VOAs that the
theorem would have to list. The Heisenberg VOA $M(1)$ has $c=1$ but is not $C_2$-cofinite, so it is excluded by
the hypothesis, as it should be: it has a continuum of irreducible modules.

%======================================================================
\section{The statement layer}\label{l22:sec:statement}
%======================================================================

All definitions of this section are in \leanpath{VOA/\allowbreak Classification/\allowbreak COne/\allowbreak Statement.lean},
which proves nothing about the classification; it only fixes the vocabulary. Throughout, $\mathcal V$ is the
carrier type of a vertex operator algebra \lean{V : VertexOperatorAlgebra 𝒱} (Chapter~\ref{ch:L03}).

\begin{definition}[Isomorphism of VOAs]\label{l22:def:iso}
Two VOAs $V$ and $T$ are \emph{isomorphic} if there is a linear bijection $e:V\to T$ with
$e(\mathbf 1)=\mathbf 1$, $e(u_{(n)}v) = e(u)_{(n)}e(v)$ for all $u,v,n$, and $e(\omega_V)=\omega_T$.
\end{definition}
\noindent File \leanpath{VOA/\allowbreak Classification/\allowbreak COne/\allowbreak Statement.lean}.
\begin{leancode}
def VOAIsomorphic (V : VertexOperatorAlgebra 𝒱) (T : VertexOperatorAlgebra 𝒲) : Prop :=
  ∃ e : 𝒱 ≃ₗ[ℂ] 𝒲, IsVOAHom V T e.toLinearMap ∧ e V.conformalVector = T.conformalVector
\end{leancode}

\begin{dictionary}
\lean{𝒱 ≃ₗ[ℂ] 𝒲} is a $\C$-linear equivalence (Chapter~\ref{ch:01}). \lean{IsVOAHom V T e.toLinearMap} is
the library's VOA homomorphism predicate (Chapter~\ref{ch:L04}): it preserves the vacuum and every
$n$-th product. The extra clause \lean{e V.conformalVector = T.conformalVector} makes the isomorphism
respect conformal vectors, hence gradings and central charges; without it two VOA structures with the
same vertex algebra but different conformal vectors (which exist, e.g. a Heisenberg VOA with shifted
$\omega$) would count as isomorphic. The two VOAs may live in \emph{different} universes (\lean{u} and
\lean{v}); this matters for the converse (Remark~\ref{l22:rem:universe}). Same as the textbook notion.
\end{dictionary}

\begin{physics}
Two chiral algebras are ``the same theory'' when there is a dictionary of states preserving the vacuum,
all operator product expansions and the stress tensor. Preserving $T(z)$ is essential: the same algebra of
fields can carry several stress tensors, which describe physically different theories (different spins).
\end{physics}

\begin{status}
Definition. It is reflexive, symmetric and transitive (\lean{VOAIsomorphic.symm}, \lean{.trans} in
\leanpath{Conditional.lean}).
\end{status}

\begin{definition}[Pseudo-unitarity; Gannon--Rayhaun Def.~2.4]\label{l22:def:pu}
A VOA $V$ is \emph{pseudo-unitary} if for every irreducible ordinary $V$-module $M\not\cong V$ and every
$\alpha\in\C$ with $M_\alpha\neq0$, $\operatorname{Re}\alpha>0$.
\end{definition}
\noindent File \leanpath{VOA/\allowbreak Classification/\allowbreak COne/\allowbreak Statement.lean}.
\begin{leancode}
def PseudoUnitary (V : VertexOperatorAlgebra 𝒱) : Prop :=
  ∀ (M : Type u) [AddCommGroup M] [Module ℂ M] (N : OrdVModule V M),
    N.IsIrreducible → IsEmpty (ModuleEquiv V N (adjointModule V)) →
      ∀ α : ℂ, N.weight α ≠ ⊥ → 0 < α.re
\end{leancode}

\begin{dictionary}
The quantifier \lean{∀ (M : Type u) …} ranges over all ordinary modules (\lean{OrdVModule}, Chapter~\ref{ch:L12},
Definition~\ref{l12:def:ordmodule}) whose carrier lives in the universe of $V$. \lean{N.IsIrreducible}: no
proper nonzero submodule. \lean{IsEmpty (ModuleEquiv V N (adjointModule V))}: $N$ is \emph{not} isomorphic to
the adjoint module $V$. \lean{N.weight α ≠ ⊥}: the $L_0$-eigenspace of eigenvalue $\alpha$ is nonzero. The
conclusion asks for $\operatorname{Re}\alpha>0$ for \emph{every} weight, not only the lowest; since the
weights of an irreducible ordinary module lie in $h+\N$ with $h$ the lowest weight, this is equivalent to
$\operatorname{Re} h>0$. GR write $h>0$; for a rational VOA all conformal weights are rational (Dong--Li--Mason),
so $\operatorname{Re}h>0$ and $h>0$ agree. \textbf{Could it be vacuous?} If $V$ had no irreducible modules other
than itself (a holomorphic VOA) the condition holds trivially --- correctly so. It is not satisfied by every
VOA: it is false for the Lee--Yang model (Chapter~\ref{ch:L19}),
whose nonvacuum module has $h=-1/5$.
\end{dictionary}

\begin{physics}
In a unitary CFT the Hamiltonian $L_0+\bar L_0$ is positive and the vacuum is the unique state of zero
energy, so every nonvacuum primary has positive weight. Pseudo-unitarity keeps this consequence without
demanding a positive inner product. It is equivalent (for strongly rational $V$) to the quantum dimensions
$d_a = S_{a0}/S_{00}$ being $\ge1$, which is how it is used below.
\end{physics}

\begin{status}
Definition. Satisfied by every $V_{2m}$ (\lean{V2m_pseudoUnitary}, Theorem~\ref{l22:thm:V2mnice}).
\end{status}

\begin{definition}[Nice; Gannon--Rayhaun Def.~2.6]\label{l22:def:nice}
$V$ is \emph{nice} if it is strongly rational and pseudo-unitary.
\end{definition}
\noindent File \leanpath{VOA/\allowbreak Classification/\allowbreak COne/\allowbreak Statement.lean}.
\begin{leancode}
def Nice (V : VertexOperatorAlgebra 𝒱) : Prop :=
  StronglyRational V ∧ PseudoUnitary V
\end{leancode}

\begin{dictionary}
\lean{StronglyRational V} is Definition~\ref{l12:def:strong} of Chapter~\ref{ch:L12}: simple, self-dual
(nondegenerate invariant bilinear form), $C_2$-cofinite and admissibly rational, with CFT type built into
\lean{VertexOperatorAlgebra}. Simplicity, self-duality, $C_2$-cofiniteness and CFT type match GR Def.~2.1.
\textbf{Rationality does not literally match.} GR's ``rational'' is Dong--Li--Mason's: every admissible module is a direct
sum of simple admissible modules. The library's \lean{AdmissiblyRational} (Chapter~\ref{ch:L12}) asks that every
\emph{weak} submodule (a mode-stable subspace, not necessarily level-graded) of an admissible module have a weak
complement. The relation between the two readings is \textbf{not settled in Lean}: the library's own note
(\leanpath{VOA/\allowbreak Modules/\allowbreak AdmissibleRationality.lean}, ll.~80--91) says that neither implication is
proved. Only one direction is known on paper: DLM-rational $\Rightarrow$ \lean{AdmissiblyRational}, by DLM's Theorem~8.1
(every admissible module of a DLM-rational VOA is a direct sum of simple ordinary modules, so every weak submodule is
complemented). The converse direction is unchecked. Note also that \lean{StronglyRational} uses
\lean{SelfDual}, not the vacuous \lean{SelfContragredient} recorded in Chapter~\ref{ch:L12}.

\textbf{The structure theorems behind niceness are proved.} The standard consequences of strong rationality that
Gannon--Rayhaun take from the literature --- regularity (Abe--Buhl--Dong, Thm.~4.5), finitely many irreducible admissible
modules, all of them ordinary (Dong--Li--Mason, Thm.~8.1(b),(c)) and hence Gannon-rationality --- are theorems of the library
for every \lean{StronglyRational} $V$, with no literature hypothesis: \lean{StronglyRational.regular},
\lean{StronglyRational.finite_Φadm} and \lean{StronglyRational.gannonRational'} (Chapter~\ref{ch:L12},
Theorem~\ref{l12:thm:srlit}; \leanpath{VOA/\allowbreak Classification/\allowbreak COne/\allowbreak LiteratureInputsProved.lean},
\texttt{\#print axioms} standard three at an earlier commit, not re-run for this edition). They are proved from the library's \lean{AdmissiblyRational} and
\lean{C2Cofinite}, so they do not settle the comparison with Dong--Li--Mason rationality described above. The one named
literature predicate of Chapter~\ref{ch:L12} that remains unproved, H.~Li's ``regular $\Rightarrow C_2$-cofinite''
(\lean{LiC2CofiniteOfRegular}), is not used anywhere in this chapter.
\end{dictionary}

\begin{remark}[What the rationality convention means for the classification]\label{l22:rem:rational}
Write ``Lean-nice'' for \lean{Nice} (with \lean{AdmissiblyRational}) and ``DLM-nice'' for Gannon--Rayhaun's niceness (with
Dong--Li--Mason rationality). Then:
\begin{itemize}
\item \textbf{The forward theorem} (Theorem~\ref{l22:thm:forward}) is proved for every Lean-nice VOA with $c=1$. By the
  paper direction DLM-rational $\Rightarrow$ \lean{AdmissiblyRational}, every DLM-nice VOA is Lean-nice, so the forward
  theorem covers every DLM-nice VOA, i.e.\ it implies Gannon--Rayhaun's Theorem~3.1 as stated --- \emph{but that implication
  rests on an on-paper argument (DLM Theorem~8.1), not on Lean}.
\item \textbf{The converse} (Theorem~\ref{l22:thm:final}) proves that every listed VOA is Lean-nice, i.e.\ nice \emph{in the
  library's sense}. It does not prove DLM complete reducibility for the listed VOAs; that would need the direction
  \lean{AdmissiblyRational} $\Rightarrow$ DLM-rational, which is not known in Lean.
\item \textbf{The classification} is therefore an ``if and only if'' for Lean-niceness. Read as a statement about GR's
  niceness, the forward half holds by the paper argument above and the converse half is the library's notion.
\end{itemize}
The release audit (\texttt{Notes/\allowbreak Textbook/\allowbreak referee/\allowbreak RELEASE\_AUDIT\_L22.md}, finding M1)
sketches a graded-projection argument for the missing direction (take the degree-$0$ part, for the level grading, of the
projection onto a level-graded submodule along a weak complement); it is offered as a pointer and is not checked.
\end{remark}

\begin{physics}
``Nice'' collects what a physicist assumes of a rational chiral algebra of a unitary CFT: finitely many
primaries, closed fusion, a modular-invariant set of characters, and positive energies.
\end{physics}

\begin{status}
Definition. Satisfiable: $V_{2m}$ and $V_{2m}^+$ are nice (Theorems~\ref{l22:thm:V2mnice}
and~\ref{l22:thm:plusnice}).
\end{status}

\begin{definition}[The rank-one lattice VOAs and their $\Z_2$-orbifolds]\label{l22:def:V2m}
For $m\ge1$, $V_{2m}$ is the lattice VOA of $\sqrt{2m}\,\Z$; $\theta$ is the lift of $-1$; $V_{2m}^+$ is the
fixed-point subalgebra of the cyclic group generated by $\theta$.
\end{definition}
\noindent File \leanpath{VOA/\allowbreak Classification/\allowbreak COne/\allowbreak Statement.lean}.
\begin{leancode}
noncomputable def V2m (m : ℕ) (hm : 1 ≤ m) : VertexOperatorAlgebra (LatModuleGen (latP m)) :=
  latticeVOAGen (latP m) (latP_even m) (latP_posDef m hm)
\end{leancode}
\noindent File \leanpath{VOA/\allowbreak Classification/\allowbreak COne/\allowbreak Statement.lean}.
\begin{leancode}
noncomputable def thetaV2m (m : ℕ) (hm : 1 ≤ m) : Aut (V2m m hm) :=
  thetaAut (latP m) (latP_even m) (latP_posDef m hm)
\end{leancode}
\noindent File \leanpath{VOA/\allowbreak Classification/\allowbreak COne/\allowbreak Statement.lean}.
\begin{leancode}
noncomputable def V2mPlus (m : ℕ) (hm : 1 ≤ m) :=
  orbifoldVOA (Subgroup.zpowers (thetaV2m m hm))
\end{leancode}

\begin{dictionary}
\lean{latP m} is the rank-one lattice $\Z$ with bilinear form $\langle x,y\rangle = 2m\,xy$, i.e.\
$\sqrt{2m}\,\Z$; \lean{latP_even m} and \lean{latP_posDef m hm} certify that it is even and positive definite
(the latter needs $m\ge1$). \lean{latticeVOAGen} is the general lattice VOA construction of
Chapter~\ref{ch:L05}; \lean{thetaAut} the lift of $-1$ (Chapter~\ref{ch:04}); \lean{Subgroup.zpowers g} the
cyclic subgroup $\langle g\rangle$; \lean{orbifoldVOA G} the fixed-point VOA $W^G$ (Chapter~\ref{ch:04}).
The hypothesis \lean{hm : 1 ≤ m} is part of the data, so $V_0$ is never formed.
\end{dictionary}

\begin{physics}
$V_{2m}$ is the chiral algebra of the compact boson at the rational radius where the momentum lattice is
$\sqrt{2m}\Z$; $V_{2m}^+$ is the chiral algebra of the corresponding point on the orbifold line, obtained by
gauging $\varphi\mapsto-\varphi$.
\end{physics}

\begin{status}
Definitions. Simplicity, self-duality, $C_2$-cofiniteness and rationality of $V_{2m}$ are proved
(\lean{V2m_stronglyRational}); so is niceness of $V_{2m}^+$ (Section~\ref{l22:sec:converse}).
\end{status}

\begin{definition}[Exceptional orbifolds and the list]\label{l22:def:list}
$V$ is a \emph{$V_2$-orbifold of type $H$} if $V\cong V_2^G$ for some subgroup $G\le\Aut V_2$ isomorphic to
the group $H$. $V$ is \emph{in the $c=1$ list} if it is isomorphic to some $V_{2m}$ ($m\ge1$), or to some
$V_{2m}^+$ ($m\ge2$), or is a $V_2$-orbifold of type $A_4$, $S_4$ or $A_5$.
\end{definition}
\noindent File \leanpath{VOA/\allowbreak Classification/\allowbreak COne/\allowbreak Statement.lean}.
\begin{leancode}
def IsV2Orbifold (V : VertexOperatorAlgebra 𝒱) (H : Type) [Group H] : Prop :=
  ∃ G : Subgroup (Aut (V2m 1 le_rfl)), Nonempty (G ≃* H) ∧ VOAIsomorphic V (orbifoldVOA G)
\end{leancode}
\noindent File \leanpath{VOA/\allowbreak Classification/\allowbreak COne/\allowbreak Statement.lean}.
\begin{leancode}
def InCOneList (V : VertexOperatorAlgebra 𝒱) : Prop :=
  (∃ (m : ℕ) (hm : 1 ≤ m), VOAIsomorphic V (V2m m hm)) ∨
  (∃ (m : ℕ) (hm : 1 ≤ m), 2 ≤ m ∧ VOAIsomorphic V (V2mPlus m hm)) ∨
  IsV2Orbifold V (alternatingGroup (Fin 4)) ∨
  IsV2Orbifold V (Equiv.Perm (Fin 4)) ∨
  IsV2Orbifold V (alternatingGroup (Fin 5))
\end{leancode}

\begin{dictionary}
\lean{V2m 1 le_rfl} is $V_2$ (\lean{le_rfl : 1 ≤ 1}). \lean{Nonempty (G ≃* H)}: there is a group isomorphism
from the subgroup $G$ (coerced to a type) to $H$. \lean{alternatingGroup (Fin 4)} is $A_4$,
\lean{Equiv.Perm (Fin 4)} is $S_4$, \lean{alternatingGroup (Fin 5)} is $A_5$. The orbifolds are stated
\emph{existentially}: \emph{some} subgroup isomorphic to $A_4$. All such subgroups of $\mathrm{PSL}_2(\C)$ are
conjugate, so the orbifold does not depend on the choice. The forward theorem and distinctness do not need this
conjugacy (distinctness works for \emph{every} such subgroup); the library does prove it for $A_4$ and $S_4$
(\lean{OrbC2d.tetrahedralConjugacy}, \lean{OrbC2d.octahedralConjugacy}), where it is used in the converse.
The side condition $m\ge2$ for $V_{2m}^+$ is GR's: $V_2^+\cong V_8$ is already in the first family.
\textbf{Could the exceptional clauses be vacuous?} If no subgroup of $\Aut V_2$ were isomorphic to $A_5$, the
clause would be empty. It is not: \lean{exists_A4_subgroup}, \lean{exists_S4_subgroup} and
\lean{exists_A5_subgroup} (Theorem~\ref{l22:thm:subgroups}) construct them.
\end{dictionary}

\begin{physics}
The three exceptional theories are gaugings of the $SU(2)_1$ WZW model by the finite subgroups of the
global $SO(3)$ symmetry that are \emph{not} contained in a $U(1)$ or its normalizer --- the binary polyhedral
groups modulo the centre. Gauging a $U(1)$ subgroup's finite cyclic or dihedral subgroup moves along the
circle or orbifold line instead.
\end{physics}

\begin{status}
Definitions; all five clauses are inhabited (Theorem~\ref{l22:thm:subgroups}).
\end{status}

\begin{theorem}[The exceptional subgroups exist]\label{l22:thm:subgroups}
$\Aut V_2$ contains subgroups isomorphic to $A_4$, $S_4$ and $A_5$.
\end{theorem}
\noindent File \leanpath{VOA/\allowbreak Classification/\allowbreak COne/\allowbreak V2Automorphisms.lean}.
\begin{leancode}
theorem exists_S4_subgroup :
    ∃ G : Subgroup (Aut (V2m 1 le_rfl)), Nonempty (G ≃* Equiv.Perm (Fin 4)) :=
\end{leancode}
\noindent File \leanpath{VOA/\allowbreak Classification/\allowbreak COne/\allowbreak V2Automorphisms.lean}.
\begin{leancode}
theorem exists_A4_subgroup :
    ∃ G : Subgroup (Aut (V2m 1 le_rfl)), Nonempty (G ≃* alternatingGroup (Fin 4)) :=
\end{leancode}
\noindent File \leanpath{VOA/\allowbreak Classification/\allowbreak COne/\allowbreak V2Icosahedral.lean}.
\begin{leancode}
theorem exists_A5_subgroup :
    ∃ G : Subgroup (Aut (V2m 1 le_rfl)), Nonempty (G ≃* alternatingGroup (Fin 5)) :=
\end{leancode}

\begin{dictionary}
The statements say exactly ``there is a subgroup $G$ with $G\cong H$''. The witnesses are images of explicit
injective homomorphisms \lean{octHom : Perm (Fin 4) →* Aut V₂} and \lean{icoHom}, built from the
$\mathfrak{sl}_2$ action on $V_2$: $S_4$ acts through rotations of the octahedron, $A_5$ through the icosahedral
rotations written with $\sqrt5$.
\end{dictionary}

\begin{proof}[Proof on paper]
$V_2$ contains the currents $h=\alpha(-1)\mathbf 1$, $e=e^{\alpha}$, $f=e^{-\alpha}$ (up to a cocycle sign), and
$\exp(\operatorname{ad}x_{(0)})$ for $x\in (V_2)_1$ is an automorphism. These generate a copy of $\mathrm{PSL}_2(\C)$
acting on $(V_2)_1\cong\mathfrak{sl}_2$ by the adjoint action. A finite rotation group $H\le SO(3)$ lifts to
$\mathrm{PSL}_2(\C)$ through $SO(3)\cong \mathrm{PSU}(2)$; writing generators as explicit matrices (the
icosahedral ones involve the golden ratio) gives a homomorphism $H\to\Aut V_2$, and injectivity follows
because its action on the weight-one space is the faithful adjoint action.
\end{proof}

\begin{physics}
These are the global symmetries of the $SU(2)_1$ theory that one can gauge to reach the isolated points of
the $c=1$ moduli space.
\end{physics}

\begin{status}
Proved, \lean{#print axioms}-clean (Lean's three).
\end{status}

\begin{definition}[The forward statement]\label{l22:def:GRS}
\emph{Gannon--Rayhaun Theorem~3.1, forward direction}: every nice VOA of central charge $1$ is in the list.
\end{definition}
\noindent File \leanpath{VOA/\allowbreak Classification/\allowbreak COne/\allowbreak Statement.lean}.
\begin{leancode}
def GannonRayhaunStatement : Prop :=
  ∀ (𝒱 : Type u) [AddCommGroup 𝒱] [Module ℂ 𝒱] (V : VertexOperatorAlgebra 𝒱),
    Nice V → V.centralCharge = 1 → InCOneList V
\end{leancode}

\begin{dictionary}
A \lean{Prop}: the claim is quantified over every carrier type \lean{𝒱} in a fixed universe \lean{u}, every
additive group and $\C$-module structure on it, and every VOA structure. The theorem
\lean{gannonRayhaun_forward : GannonRayhaunStatement.{u}} is universe-polymorphic, so it holds in every
universe. Note that \lean{V.centralCharge = 1} is a hypothesis \emph{in addition} to niceness. Same as GR, with
``rational'' read as \lean{AdmissiblyRational} (Remark~\ref{l22:rem:rational}).
\end{dictionary}

\begin{physics}
``Every rational unitary chiral algebra with $c=1$ is a free boson at a rational point, its $\Z_2$ orbifold, or
one of three exceptional gaugings of $SU(2)_1$.''
\end{physics}

\begin{status}
Proved without hypotheses, unconditional (it uses the theorem \lean{VOA.pbw}; Theorem~\ref{l22:thm:forward}).
\end{status}

\begin{definition}[Labels, distinctness and converse]\label{l22:def:labels}
The labels are $\mathrm{lat}(m)$, $\mathrm{plus}(m)$, $\mathrm{tetra}$, $\mathrm{octa}$, $\mathrm{icosa}$;
$V$ has label $l$ if it is isomorphic to the corresponding VOA (with the side conditions $m\ge1$, resp.\ $m\ge2$).
\emph{Distinctness}: a VOA has at most one label. \emph{Converse}: every VOA in the list is nice and has
central charge $1$.
\end{definition}
\noindent File \leanpath{VOA/\allowbreak Classification/\allowbreak COne/\allowbreak Statement.lean}.
\begin{leancode}
inductive COneLabel
  | lat (m : ℕ)
  | plus (m : ℕ)
  | tetra
  | octa
  | icosa
  deriving DecidableEq
\end{leancode}
\noindent File \leanpath{VOA/\allowbreak Classification/\allowbreak COne/\allowbreak Statement.lean}.
\begin{leancode}
def HasCOneLabel (V : VertexOperatorAlgebra 𝒱) : COneLabel → Prop
  | .lat m => ∃ hm : 1 ≤ m, VOAIsomorphic V (V2m m hm)
  | .plus m => ∃ hm : 1 ≤ m, 2 ≤ m ∧ VOAIsomorphic V (V2mPlus m hm)
  | .tetra => IsV2Orbifold V (alternatingGroup (Fin 4))
  | .octa => IsV2Orbifold V (Equiv.Perm (Fin 4))
  | .icosa => IsV2Orbifold V (alternatingGroup (Fin 5))
\end{leancode}
\noindent File \leanpath{VOA/\allowbreak Classification/\allowbreak COne/\allowbreak Statement.lean}.
\begin{leancode}
def GannonRayhaunDistinct : Prop :=
  ∀ (𝒱 : Type u) [AddCommGroup 𝒱] [Module ℂ 𝒱] (V : VertexOperatorAlgebra 𝒱) (l l' : COneLabel),
    HasCOneLabel V l → HasCOneLabel V l' → l = l'
\end{leancode}
\noindent File \leanpath{VOA/\allowbreak Classification/\allowbreak COne/\allowbreak Statement.lean}.
\begin{leancode}
def GannonRayhaunConverse : Prop :=
  ∀ (𝒱 : Type u) [AddCommGroup 𝒱] [Module ℂ 𝒱] (V : VertexOperatorAlgebra 𝒱),
    InCOneList V → Nice V ∧ V.centralCharge = 1
\end{leancode}

\begin{dictionary}
\lean{HasCOneLabel V (.lat m)} unfolds to \lean{∃ hm : 1 ≤ m, VOAIsomorphic V (V2m m hm)}: the side
condition is \emph{inside} the predicate, so an ``invalid'' label such as $\mathrm{lat}(0)$ is never realized.
\lean{inCOneList_iff_exists_label} proves $\text{InCOneList}\,V \iff \exists l,\ \text{HasCOneLabel}\,V\,l$.
\textbf{Distinctness} quantifies over all $V$ and pairs of labels and concludes \lean{l = l'}; in particular it
says $V_{2m}\not\cong V_{2m'}$ for $m\ne m'$, $V_{2m}\not\cong V_{2m'}^+$, and that the three exceptional orbifolds
are pairwise non-isomorphic and not isomorphic to any lattice or plus VOA. Together with the forward theorem
it gives GR's ``exactly one''. \textbf{The converse} includes $c=1$ in its conclusion.
\end{dictionary}

\begin{physics}
Distinctness says the listed theories are genuinely different chiral algebras (different spectra); the
converse says each of them is a legitimate rational unitary chiral algebra.
\end{physics}

\begin{status}
Definitions. Distinctness is proved (Theorem~\ref{l22:thm:distinct}); the converse is proved with no hypothesis, in
universe \lean{0}, using the theorem \lean{VOA.pbw} (Theorem~\ref{l22:thm:final}).
\end{status}

%======================================================================
\section{Modular data and the seven external inputs}\label{l22:sec:inputs}
%======================================================================

The proof of Gannon and Rayhaun is a modular-forms argument resting on several large published theorems.
The library's strategy, fixed early in the project, was: \emph{first} state each published theorem as a named
\lean{Prop} at its natural generality, prove the classification \emph{conditionally} on these
(Section~\ref{l22:sec:conditional}), and \emph{then} discharge each \lean{Prop} by an actual proof
(Section~\ref{l22:sec:discharge}). All inputs are in
\leanpath{VOA/\allowbreak Classification/\allowbreak COne/\allowbreak Inputs.lean}. We present them in full,
because the referee must agree that each is a true theorem; if one were false (or too strong), the
conditional theorem would be vacuous.

\subsection{Shared vocabulary}

\begin{definition}[Möbius action, $\theta$-type forms]\label{l22:def:thetatype}
$\gamma=\begin{pmatrix}a&b\\c&d\end{pmatrix}\in\ltwotwoSL$ acts on $\C$ by $\gamma\tau = (a\tau+b)/(c\tau+d)$.
A function $f$ on the upper half plane $\ltwotwoH$ is a \emph{holomorphic $\theta$-type modular form of
weight $\tfrac12$ for $\Gamma_1(N)$} if it is holomorphic on $\ltwotwoH$, transforms like
$\theta(\tau)=\sum_n q^{n^2}$ under $\Gamma_1(N)$, i.e.\ $f(\gamma\tau)/f(\tau)=\theta(\gamma\tau)/\theta(\tau)$, and
$|(c\tau+d)^{-1/2}f(\gamma\tau)|$ is bounded on $\operatorname{Im}\tau\ge1$ for every $\gamma\in\ltwotwoSL$
(holomorphy at all cusps).
\end{definition}
\noindent File \leanpath{VOA/\allowbreak Classification/\allowbreak COne/\allowbreak Inputs.lean}.
\begin{leancode}
def mobius (γ : Matrix.SpecialLinearGroup (Fin 2) ℤ) (τ : ℂ) : ℂ :=
  ((γ 0 0 : ℂ) * τ + γ 0 1) / ((γ 1 0 : ℂ) * τ + γ 1 1)
\end{leancode}
\noindent File \leanpath{VOA/\allowbreak Classification/\allowbreak COne/\allowbreak Inputs.lean}.
\begin{leancode}
def thetaRef (τ : ℂ) : ℂ := rankOneTheta 0 1 1 τ
\end{leancode}
\noindent File \leanpath{VOA/\allowbreak Classification/\allowbreak COne/\allowbreak Inputs.lean}.
\begin{leancode}
structure IsThetaTypeForm (N : ℕ) (f : ℂ → ℂ) : Prop where
  holo : DifferentiableOn ℂ f {τ : ℂ | 0 < τ.im}
  multiplier : ∀ γ ∈ CongruenceSubgroup.Gamma1 N, ∀ τ : ℂ, 0 < τ.im →
    f (mobius γ τ) * thetaRef τ = thetaRef (mobius γ τ) * f τ
  cusps : ∀ γ : Matrix.SpecialLinearGroup (Fin 2) ℤ, ∃ C : ℝ, ∀ τ : ℂ, 1 ≤ τ.im →
    ‖f (mobius γ τ)‖ ≤ C * ‖(γ 1 0 : ℂ) * τ + γ 1 1‖ ^ (1 / 2 : ℝ)
\end{leancode}

\begin{dictionary}
\lean{Matrix.SpecialLinearGroup (Fin 2) ℤ} is $\ltwotwoSL$, with entries \lean{γ 0 0}$=a$, \lean{γ 0 1}$=b$,
\lean{γ 1 0}$=c$, \lean{γ 1 1}$=d$. Functions on $\ltwotwoH$ are encoded as functions $\C\to\C$ constrained only
where \lean{0 < τ.im}. The multiplier condition is \emph{cross-multiplied} ($f(\gamma\tau)\theta(\tau) =
\theta(\gamma\tau)f(\tau)$), which avoids dividing by $\theta$ (it has no zeros on $\ltwotwoH$, so the two forms are
equivalent). The cusp condition is a growth bound $\|f(\gamma\tau)\|\le C|c\tau+d|^{1/2}$; for a form with
this multiplier, it is exactly the absence of negative powers in the Puiseux $q$-expansion at the cusp
$\gamma(i\infty)$. \lean{CongruenceSubgroup.Gamma1 N} is Mathlib's $\Gamma_1(N)$. \lean{rankOneTheta} is
Definition~\ref{l22:def:theta} below.
\end{dictionary}

\begin{physics}
Weight-$\frac12$ forms are the characters of a single chiral boson times $\eta$: the $\eta$ removes the
oscillator contribution, leaving the theta function that counts momentum states.
\end{physics}

\begin{status}
Definitions.
\end{status}

\begin{definition}[Graded dimension and Euler's function as formal series]\label{l22:def:series}
$\mathrm{vac}_V=\sum_{n\ge0}\dim V_n\, q^n\in\Z[[q]]$ and $\prod_{n\ge1}(1-q^n)\in\Z[[q]]$.
\end{definition}
\noindent File \leanpath{VOA/\allowbreak Classification/\allowbreak COne/\allowbreak Inputs.lean}.
\begin{leancode}
def vacuumSeries {𝒱 : Type u} [AddCommGroup 𝒱] [Module ℂ 𝒱] (V : VertexOperatorAlgebra 𝒱) :
    PowerSeries ℤ :=
  PowerSeries.mk fun n => (finrank ℂ (V.grading n) : ℤ)
\end{leancode}
\noindent File \leanpath{VOA/\allowbreak Classification/\allowbreak COne/\allowbreak Inputs.lean}.
\begin{leancode}
def eulerSeries : PowerSeries ℤ :=
  PowerSeries.mk fun N =>
    PowerSeries.coeff N (∏ n ∈ Finset.range N, (1 - PowerSeries.X ^ (n + 1) : PowerSeries ℤ))
\end{leancode}

\begin{dictionary}
\lean{PowerSeries.mk f} is the series with coefficients $f(n)$; \lean{PowerSeries.coeff N} extracts the $N$-th
coefficient. The infinite product is defined coefficientwise by the finite product over $n<N$ --- correct,
because factors $1-q^{n+1}$ with $n+1>N$ do not affect the coefficient of $q^N$. These are formal series: the
link to analytic characters is always an explicit \lean{HasSum} (next definition).
\end{dictionary}

\begin{physics}
$\mathrm{vac}_V$ counts states of the chiral algebra by energy; $\prod(1-q^n)^{-1}$ counts the oscillator
states of one free boson, so the product $\prod(1-q^n)\cdot\mathrm{vac}_V$ ``removes one boson''.
\end{physics}

\begin{status}
Definitions.
\end{status}

\subsection{Modular data}

\begin{definition}[Candidate modular data]\label{l22:def:moddata}
Candidate modular data for $V$ consist of a finite index set $I$ with a vacuum label $0$, a $V$-module $M_a$ for
each $a\in I$, rational numbers $h_a$, functions $\ltwotwoch_a$ on $\ltwotwoH$, a matrix $S\in\C^{I\times I}$, a real
number $S_{00}$ and real numbers $d_a$. No properties are imposed.
\end{definition}
\noindent File \leanpath{VOA/\allowbreak Classification/\allowbreak COne/\allowbreak Inputs.lean}.
\begin{leancode}
structure ModularData {𝒱 : Type u} [AddCommGroup 𝒱] [Module ℂ 𝒱]
    (V : VertexOperatorAlgebra 𝒱) where
  ι : Type
  [instFintype : Fintype ι]
  [instDecEq : DecidableEq ι]
  vac : ι
  carrier : ι → Type u
  [instACG : ∀ a, AddCommGroup (carrier a)]
  [instMod : ∀ a, Module ℂ (carrier a)]
  module : ∀ a, OrdVModule V (carrier a)
  h : ι → ℚ
  ch : ι → ℂ → ℂ
  S : ι → ι → ℂ
  S00 : ℝ
  d : ι → ℝ
\end{leancode}

\begin{dictionary}
A Lean \lean{structure} bundles data. \lean{ι : Type} with \lean{[Fintype ι] [DecidableEq ι]} is a finite
label set; \lean{carrier a : Type u} is the underlying space of the $a$-th module, in the universe of $V$;
\lean{module a : OrdVModule V (carrier a)} is the module structure. Everything else is plain data;
\lean{IsModularData} below says which data are \emph{the} modular data of $V$.
\end{dictionary}

\begin{physics}
The modular data of a rational CFT: the list of primaries, their conformal weights, their characters, and
the modular $S$-matrix that implements $\tau\mapsto-1/\tau$ (exchange of space and time on the torus).
\end{physics}

\begin{status}
Definition.
\end{status}

\begin{definition}[The trace function and $S$-pinning]\label{l22:def:spin}
For a label $a$ and a state $v$, $T_a(v,\tau)=\operatorname{tr}_{M_a} o(v)\,q^{L_0-c/24}$ (Zhu's one-point
function). $S$ is \emph{pinned} if for every square-bracket-homogeneous $v$ of weight $n$,
$T_a(v,-1/\tau)=\tau^n\sum_b S_{ab}T_b(v,\tau)$.
\end{definition}
\noindent File \leanpath{VOA/\allowbreak Classification/\allowbreak COne/\allowbreak Inputs.lean}.
\begin{leancode}
def trace (D : ModularData V) (a : D.ι) (v : 𝒱) (τ : UpperHalfPlane) : ℂ :=
  qEvalAt (onePoint (D.module a) v) ((D.h a : ℂ) - V.centralCharge / 24) τ
\end{leancode}
\noindent File \leanpath{VOA/\allowbreak Classification/\allowbreak COne/\allowbreak Inputs.lean}.
\begin{leancode}
def SPinned (D : ModularData V) (S : D.ι → D.ι → ℂ) : Prop :=
  ∀ (n : ℕ) (v : 𝒱), v ∈ sqWeightSpace V (n : ℂ) → ∀ (a : D.ι) (τ : UpperHalfPlane),
    D.trace a v (ModularGroup.S • τ) = (τ : ℂ) ^ n * ∑ b, S a b * D.trace b v τ
\end{leancode}

\begin{dictionary}
\lean{onePoint (D.module a) v} is the formal one-point series $\sum_n \operatorname{tr}_{(M_a)_{h_a+n}}o(v)\,q^n$
and \lean{qEvalAt … s τ} evaluates $q^{s}\sum c_n q^n$ at $q=e^{2\pi i\tau}$ (Chapter~\ref{ch:L14}); the shift
$s = h_a-c/24$ restores the $q^{L_0-c/24}$. \lean{sqWeightSpace V n} is the weight-$n$ space for Zhu's
square-bracket $L[0]$ (Chapter~\ref{ch:L15}); \lean{ModularGroup.S • τ} is $-1/\tau$. This is Zhu's theorem
(Zhu 1996, (5.3.18a)) at $\gamma=S$.
\end{dictionary}

\begin{warning}[title={The $S$-matrix is not determined by characters}]\label{l22:warn:spin}
The first version of the modular-data package required only the character law
$\ltwotwoch_a(-1/\tau)=\sum_b S_{ab}\ltwotwoch_b(\tau)$. Characters do not separate a module from its
contragredient ($\ltwotwoch_M=\ltwotwoch_{M'}$), so the matrix $S$ was not determined: the library proved that a
whole continuum of perturbed $S$-matrices satisfied the old package
(\lean{VOA.Huang.HuangModular.isModularDataWeak_perturbS}). Since several later steps read off
properties of $S$, an input satisfied by wrong $S$-matrices is dangerous. The repaired definition adds
\lean{S_pinned}, and \lean{VOA.Huang.SPinning.S_unique_of_pinned} proves that a pinned $S$ is unique, because
one-point functions of inequivalent irreducibles are linearly independent.
\end{warning}

\begin{physics}
Characters are partition functions without insertions; one-point functions are torus amplitudes with one
operator inserted. The $S$-matrix is defined by the modular transformation of all torus amplitudes, and only
the insertions distinguish a sector from its charge conjugate.
\end{physics}

\begin{status}
Definitions. The uniqueness theorem is proved (Chapter~\ref{ch:L21}).
\end{status}

\begin{definition}[Modular data of a nice VOA]\label{l22:def:ismoddata}
The candidate data are \emph{the modular data of $V$} if: (classification) each $M_a$ is irreducible, every
irreducible ordinary module is isomorphic to exactly one $M_a$, and $M_0\cong V$; (grading) the weights of $M_a$
lie in $h_a+\N$, $h_a$ is attained, $h_0=0$; (characters) $\ltwotwoch_a(\tau)=\sum_n \dim(M_a)_{h_a+n}\,
q^{h_a+n-c/24}$, a convergent sum, and the vacuum character is $\sum_n\dim V_n\,q^{n-c/24}$; ($S$) the
character law and the pinned law; (pseudo-unitarity) the real number $S_{00}$ is positive and equals the matrix entry $S_{\mathrm{vac},\mathrm{vac}}$,
$S_{a0}=S_{0a}=S_{00}d_a$ with
$d_a\ge1$; (Gannon's bound) $|S_{ab}|\le S_{00}d_ad_b$.
\end{definition}
\noindent File \leanpath{VOA/\allowbreak Classification/\allowbreak COne/\allowbreak Inputs.lean}.
\begin{leancode}
structure IsModularData (D : ModularData V) : Prop where
  irred : ∀ a, (D.module a).IsIrreducible
  exhaust : ∀ (M : Type u) [AddCommGroup M] [Module ℂ M] (N : OrdVModule V M),
    N.IsIrreducible → ∃ a, Nonempty (ModuleEquiv V N (D.module a))
  distinct : ∀ a b, Nonempty (ModuleEquiv V (D.module a) (D.module b)) → a = b
  vac_iso : Nonempty (ModuleEquiv V (D.module D.vac) (adjointModule V))
  weight_support : ∀ a (α : ℂ), (D.module a).weight α ≠ ⊥ → ∃ n : ℕ, α = (D.h a : ℂ) + n
  weight_lowest : ∀ a, (D.module a).weight (D.h a : ℂ) ≠ ⊥
  h_vac : D.h D.vac = 0
  char : ∀ a (τ : ℂ), 0 < τ.im →
    HasSum (fun n : ℕ => (finrank ℂ ((D.module a).weight ((D.h a : ℂ) + n)) : ℂ) *
      cexp (2 * Real.pi * I * ((D.h a : ℂ) + n - V.centralCharge / 24) * τ)) (D.ch a τ)
  char_vac : ∀ τ : ℂ, 0 < τ.im →
    HasSum (fun n : ℕ => (finrank ℂ (V.grading n) : ℂ) *
      cexp (2 * Real.pi * I * ((n : ℂ) - V.centralCharge / 24) * τ)) (D.ch D.vac τ)
  S_transform : ∀ a (τ : ℂ), 0 < τ.im → D.ch a (-1 / τ) = ∑ b, D.S a b * D.ch b τ
  S_pinned : D.SPinned D.S
  S00_pos : 0 < D.S00
  S_vac_vac : D.S D.vac D.vac = D.S00
  d_ge_one : ∀ a, 1 ≤ D.d a
  S_vac_left : ∀ a, D.S D.vac a = D.S00 * D.d a
  S_vac_right : ∀ a, D.S a D.vac = D.S00 * D.d a
  S_bound : ∀ a b, ‖D.S a b‖ ≤ D.S00 * D.d a * D.d b
\end{leancode}

\begin{dictionary}
Each field is one clause of the mathematical definition. \lean{cexp (2 * Real.pi * I * (…) * τ)} is
$q^{(\dots)}$. \lean{HasSum f x} is unconditional convergence of $\sum_n f(n)$ to $x$ (Chapter~\ref{ch:01}),
so \lean{char} says the character is an honest convergent series of \emph{actual} graded dimensions --- this
is what prevents the package from being satisfied by arbitrary functions. \lean{S_vac_left},
\lean{S_vac_right} and \lean{d_ge_one} encode pseudo-unitarity in its categorical form: $d_a=S_{a0}/S_{00}\ge1$
are the Frobenius--Perron dimensions. \lean{S_bound} is Gannon (2005), (4.2a), quoted as GR (2.35).
\textbf{Stronger or weaker?} This is the standard package of GR §2.1. It contains nothing specific to $c=1$.
\textbf{Vacuity:} every clause constrains the actual modules of $V$; the package is satisfiable (it holds for
every nice $V$ by \lean{iModData_holds}, Chapter~\ref{ch:L21}).
\end{dictionary}

\begin{physics}
A rational CFT's torus partition functions form a finite-dimensional representation of $\ltwotwoSL$. The
entries $S_{a0}/S_{00}$ are the quantum dimensions: asymptotic ratios of the number of states in sector $a$ and
in the vacuum sector at high energy. Pseudo-unitarity makes them $\ge1$, as for the dimensions of
representations of a group.
\end{physics}

\begin{status}
Definition; satisfied by every nice VOA, proved using the theorem \lean{VOA.pbw} (\lean{iModData_holds}).
\end{status}

\begin{definition}[Level and Galois symmetry]\label{l22:def:level}
$N>0$ is a \emph{level} if $N(h_a-c/24)\in\Z$ for all $a$. The data are \emph{Galois-symmetric at level $N$}
if for all $\ell,\ell'$ with $\ell\ell'\equiv1\pmod N$ and every $\gamma\equiv\operatorname{diag}(\ell,\ell')
\pmod N$ there are a permutation $\pi$ of the labels and signs $\varepsilon_a=\pm1$ with
$\ltwotwoch_a(\gamma\tau)=\varepsilon_a\ltwotwoch_{\pi a}(\tau)$.
\end{definition}
\noindent File \leanpath{VOA/\allowbreak Classification/\allowbreak COne/\allowbreak Inputs.lean}.
\begin{leancode}
def IsLevel (D : ModularData V) (N : ℕ) : Prop :=
  0 < N ∧ ∀ a, ∃ k : ℤ, (N : ℂ) * ((D.h a : ℂ) - V.centralCharge / 24) = k
\end{leancode}
\noindent File \leanpath{VOA/\allowbreak Classification/\allowbreak COne/\allowbreak Inputs.lean}.
\begin{leancode}
def GaloisSymmetric (D : ModularData V) (N : ℕ) : Prop :=
  ∀ ℓ ℓ' : ℤ, ℓ * ℓ' ≡ 1 [ZMOD N] → ∀ γ : Matrix.SpecialLinearGroup (Fin 2) ℤ,
    γ 0 0 ≡ ℓ [ZMOD N] → γ 1 1 ≡ ℓ' [ZMOD N] → γ 0 1 ≡ 0 [ZMOD N] → γ 1 0 ≡ 0 [ZMOD N] →
    ∃ (π : D.ι ≃ D.ι) (ε : D.ι → ℂ), (∀ a, ε a = 1 ∨ ε a = -1) ∧
      ∀ a (τ : ℂ), 0 < τ.im → D.ch a (mobius γ τ) = ε a * D.ch (π a) τ
\end{leancode}

\begin{dictionary}
\lean{[ZMOD N]} is congruence of integers modulo $N$. The conductor $C_V$ of GR is the least level; every
level is a multiple of it, so stating things ``for every level'' is equivalent to ``at the conductor''.
\end{dictionary}

\begin{physics}
$q^{h_a-c/24}$ is $N$-periodic in $\tau$ exactly when $N$ is a level, so $T^N$ acts trivially. Galois symmetry
is the statement that the modular representation is defined over a cyclotomic field and the Galois group acts on
the $S$-matrix by signed permutations (Coste--Gannon).
\end{physics}

\begin{status}
Definitions.
\end{status}

\subsection{The seven inputs}

\begin{definition}[The seven external inputs]\label{l22:def:inputs}
\begin{enumerate}
\item \textbf{IModData}: every nice VOA has modular data (Zhu, Dong--Li--Mason, Huang, Gannon).
\item \textbf{ICongruenceGalois}: for every level $N$, the characters are $\Gamma(N)$-invariant and Galois
  symmetric (Dong--Lin--Ng 2015; Coste--Gannon; Ng--Schauenburg).
\item \textbf{ISerreStark}: for $4\mid N$, every holomorphic $\theta$-type weight-$\frac12$ form for
  $\Gamma_1(N)$ is a finite combination of theta series $\theta_{n_0,r,t}$ with $4tr^2\mid N$ (Serre--Stark
  1977). The library uses only levels divisible by $24$ (\textbf{ISerreStark24}).
\item \textbf{IVirasoro}: for $c=1$, $\prod(1-q^n)\cdot\mathrm{vac}_V=\sum_h A_h\,\eta\chi_h$ with
  multiplicities $A_h\in\N$, $A_0=1$ (Kac; Feigin--Fuchs).
\item \textbf{ITannakian}: if $S_{ab}=S_{00}d_ad_b$ on integer-weight labels, then $V\cong W^G$ for a nice $W$ of
  the same central charge and a finite $G\le\Aut W$, with $\dim W_n=\sum_{h_a\in\Z}d_a\dim(M_a)_n$ (GR
  Props.~2.13, 2.17: Müger, Deligne, Kirillov--Ostrik, Huang--Kirillov--Lepowsky).
\item \textbf{IDongMason}: a nice $c=1$ VOA $W$ with $W_1\ne0$ is some $V_{2m}$ (Dong--Mason 2004).
\item \textbf{IEnumeration}: $V_{2m}^G$ is in the list for every finite $G\le\Aut V_{2m}$ (GR Prop.~C.4).
\end{enumerate}
\end{definition}
\noindent File \leanpath{VOA/\allowbreak Classification/\allowbreak COne/\allowbreak Inputs.lean}.
\begin{leancode}
def IModData : Prop :=
  ∀ (𝒱 : Type u) [AddCommGroup 𝒱] [Module ℂ 𝒱] (V : VertexOperatorAlgebra 𝒱), Nice V →
    ∃ D : ModularData V, D.IsModularData
\end{leancode}
\noindent File \leanpath{VOA/\allowbreak Classification/\allowbreak COne/\allowbreak Inputs.lean}.
\begin{leancode}
def ICongruenceGalois : Prop :=
  ∀ (𝒱 : Type u) [AddCommGroup 𝒱] [Module ℂ 𝒱] (V : VertexOperatorAlgebra 𝒱), Nice V →
    ∀ D : ModularData V, D.IsModularData → ∀ N : ℕ, D.IsLevel N →
      (∀ γ ∈ CongruenceSubgroup.Gamma N, ∀ a (τ : ℂ), 0 < τ.im →
        D.ch a (mobius γ τ) = D.ch a τ) ∧ D.GaloisSymmetric N
\end{leancode}
\noindent File \leanpath{VOA/\allowbreak Classification/\allowbreak COne/\allowbreak Inputs.lean}.
\begin{leancode}
def ISerreStark24 : Prop :=
  ∀ (N : ℕ) (f : ℂ → ℂ), 0 < N → 24 ∣ N → IsThetaTypeForm N f →
    ∃ (s : Finset (ℤ × ℕ × ℕ)) (coef : ℤ × ℕ × ℕ → ℂ),
      (∀ p ∈ s, 1 ≤ p.2.1 ∧ 1 ≤ p.2.2 ∧ 4 * p.2.2 * p.2.1 ^ 2 ∣ N) ∧
      ∀ τ : ℂ, 0 < τ.im → f τ = ∑ p ∈ s, coef p * rankOneTheta p.1 p.2.1 p.2.2 τ
\end{leancode}
\noindent File \leanpath{VOA/\allowbreak Classification/\allowbreak COne/\allowbreak Inputs.lean}.
\begin{leancode}
def IVirasoro : Prop :=
  ∀ (𝒱 : Type u) [AddCommGroup 𝒱] [Module ℂ 𝒱] (V : VertexOperatorAlgebra 𝒱),
    V.vacuum ≠ 0 → V.centralCharge = 1 →
      ∃ A : ℕ → ℕ, A 0 = 1 ∧ eulerSeries * vacuumSeries V = etaVirasoroSum A
\end{leancode}
\noindent File \leanpath{VOA/\allowbreak Classification/\allowbreak COne/\allowbreak Inputs.lean}.
\begin{leancode}
def ITannakian : Prop :=
  ∀ (𝒱 : Type u) [AddCommGroup 𝒱] [Module ℂ 𝒱] (V : VertexOperatorAlgebra 𝒱), Nice V →
    ∀ D : ModularData V, D.IsModularData →
    (∀ a b, D.IsIntWeight a → D.IsIntWeight b → D.S a b = D.S00 * D.d a * D.d b) →
    ∃ (𝒲 : Type u) (_ : AddCommGroup 𝒲) (_ : Module ℂ 𝒲) (W : VertexOperatorAlgebra 𝒲)
      (G : Subgroup (Aut W)), Nice W ∧ Finite G ∧ W.centralCharge = V.centralCharge ∧
      VOAIsomorphic V (orbifoldVOA G) ∧
      ∀ n : ℕ, (finrank ℂ (W.grading n) : ℝ) =
        ∑ a ∈ D.intLabels, D.d a * (finrank ℂ ((D.module a).weight (n : ℂ)) : ℝ)
\end{leancode}
\noindent File \leanpath{VOA/\allowbreak Classification/\allowbreak COne/\allowbreak Inputs.lean}.
\begin{leancode}
def IDongMason : Prop :=
  ∀ (𝒲 : Type u) [AddCommGroup 𝒲] [Module ℂ 𝒲] (W : VertexOperatorAlgebra 𝒲), Nice W →
    W.centralCharge = 1 → finrank ℂ (W.grading 1) ≠ 0 →
      ∃ (m : ℕ) (hm : 1 ≤ m), VOAIsomorphic W (V2m m hm)
\end{leancode}
\noindent File \leanpath{VOA/\allowbreak Classification/\allowbreak COne/\allowbreak Inputs.lean}.
\begin{leancode}
def IEnumeration : Prop :=
  ∀ (m : ℕ) (hm : 1 ≤ m) (G : Subgroup (Aut (V2m m hm))), Finite G →
    InCOneList (orbifoldVOA G)
\end{leancode}

\begin{dictionary}
\textbf{Each input is a \lean{def … : Prop}, never an \lean{axiom}.} Nothing is assumed by declaring them; the
conditional theorem takes them as explicit hypotheses. We check each against its source.
\begin{itemize}
\item \lean{IModData}: quantified over every nice $V$ in universe \lean{u}; the conclusion is the existence of
  modular data in the sense of Definition~\ref{l22:def:ismoddata}. Natural generality, no $c=1$.
\item \lean{ICongruenceGalois}: for every nice $V$, every modular data \lean{D} of $V$ (not a chosen one) and
  every level, (1) $\Gamma(N)$-invariance of \emph{all} characters, (2) Galois symmetry. This is DLN Theorem~I plus
  the Coste--Gannon/Ng--Schauenburg Galois action.
\item \lean{ISerreStark24}: a statement about modular forms only; the triple \lean{p = (n₀, r, t)} encodes
  $\theta_{n_0,r,t}$, and \lean{4 * p.2.2 * p.2.1 ^ 2 ∣ N} is $4tr^2\mid N$. Serre--Stark's basis
  $\theta_{\psi,t}=\sum\psi(n)q^{tn^2}$ is a combination of these. The library also states the version for
  $4\mid N$ (\lean{ISerreStark}) and proves \lean{iSerreStark24_of_iSerreStark}.
\item \lean{IVirasoro}: the hypothesis \lean{V.vacuum ≠ 0} excludes the zero space, for which $A_0=1$ would be
  false. \lean{etaVirasoroSum A} is $\sum_hA_h\,\eta\chi_h$ with GR's $c=1$ Virasoro characters
  (Definition~\ref{l22:def:etavir}). It holds for every $c=1$ VOA, rational or not, so it carries no
  classification content by itself.
\item \lean{ITannakian}: the hypothesis is the factorization of $S$ on integer-weight labels; the conclusion
  produces $W$, $G$ with \lean{Nice W}, \lean{Finite G}, equal central charges, \lean{V ≅ W^G}, and the graded
  dimension formula. \textbf{Flag}: this is a GR proposition assembled from several published results, not a single
  published theorem. It is stated at the generality of GR §2.2 (any nice $V$).
\item \lean{IDongMason}: Dong--Mason Theorems~1 and~3 (the weight-one Lie algebra is reductive of rank
  $\le\tilde c$, with equality iff $W$ is a lattice VOA), with the elementary steps $\tilde c=c=1$
  (pseudo-unitarity) and ``rank-one even lattices are $\sqrt{2m}\Z$'' folded in.
\item \lean{IEnumeration}: GR Prop.~C.4 --- finite subgroups of $\Aut V_2\cong\mathrm{PSL}_2(\C)$ and of
  $\Aut V_{2m}\cong\C^\times\rtimes\Z_2$ ($m\ge2$), plus $V_{2m}^{\Z_n}\cong V_{2mn^2}$ and
  $V_{2m}^{D_{2n}}\cong V_{2mn^2}^+$. \textbf{Flag}: assembled.
\end{itemize}
\textbf{Could the conditional theorem be vacuous?} Only if some input were false. Each is a true published theorem
(or a GR proposition assembled from them), and --- more decisively --- \emph{each has since been proved in Lean}
(Section~\ref{l22:sec:discharge}), so the question is settled by the machine.
\end{dictionary}

\begin{physics}
The inputs are the standard toolkit of rational CFT: modular invariance of torus amplitudes (Zhu), the
modular tensor category of anyons (Huang), the congruence property (the representation of $\ltwotwoSL$ factors
through a finite quotient), the theta-series description of one-boson partition functions (Serre--Stark),
the Virasoro decomposition of the vacuum module, the gauging/ungauging correspondence (Tannakian
reconstruction: a symmetric subcategory of anyons with trivial spins is $\operatorname{Rep}G$ and can be
condensed), and the fact that a rank-one current algebra at $c=1$ is a lattice boson.
\end{physics}

\begin{status}
Definitions (\lean{Prop}s). All seven are now theorems of the library, see
Section~\ref{l22:sec:discharge}: \lean{IModData} (using the theorem \lean{VOA.pbw}), \lean{IEnumeration} and the
Serre--Stark and Virasoro inputs with Lean's three axioms only.
\end{status}

%======================================================================
\section{The paper's own lemmas}\label{l22:sec:core}
%======================================================================

The new mathematics of Gannon--Rayhaun §3 is formalized directly, as four lemmas: the
$S$-transformation of the rank-one theta functions, the positivity of the theta weights (GR Lemma~3.11), the
uniqueness of the surviving theta weight (GR Prop.~3.12), and the Tannakian squeeze (GR Theorem~2.18). They are
stated in elementary terms (theta functions, finite sums, power series), so the referee can check them without any
VOA theory.

\subsection{The rank-one theta functions and their \texorpdfstring{$S$}{S}-transformation}

\begin{definition}[Rank-one theta functions; GR (3.9)]\label{l22:def:theta}
For $n_0,r\in\Z$ and $t>0$,
$\theta_{n_0,r,t}(\tau)=\sum_{n\equiv n_0\ (r)} q^{tn^2}=\sum_{m\in\Z}e^{2\pi i t(n_0+rm)^2\tau}$.
\end{definition}
\noindent File \leanpath{VOA/\allowbreak Classification/\allowbreak COne/\allowbreak RankOneTheta.lean}.
\begin{leancode}
def rankOneThetaTerm (n₀ r : ℤ) (t : ℝ) (τ : ℂ) (m : ℤ) : ℂ :=
  cexp (2 * π * I * t * ((n₀ + r * m : ℤ) : ℂ) ^ 2 * τ)
\end{leancode}
\noindent File \leanpath{VOA/\allowbreak Classification/\allowbreak COne/\allowbreak RankOneTheta.lean}.
\begin{leancode}
def rankOneTheta (n₀ r : ℤ) (t : ℝ) (τ : ℂ) : ℂ :=
  ∑' m : ℤ, rankOneThetaTerm n₀ r t τ m
\end{leancode}

\begin{dictionary}
\lean{∑' m : ℤ} is Mathlib's \lean{tsum}: the sum if it converges unconditionally, and $0$ otherwise (the
junk-value convention of Chapter~\ref{ch:01}). For $\operatorname{Im}\tau>0$ and $t>0$ the series converges
absolutely (Gaussian decay), so the junk value never occurs where the theorems use it; the library also proves
\lean{rankOneTheta_eq_jacobiTheta₂}, identifying it with Mathlib's two-variable Jacobi theta function, which
settles convergence. The parameter $t$ is a real number in the definition and a natural number in all uses.
\end{dictionary}

\begin{physics}
$\theta_{n_0,r,t}/\eta$ is the character of the momentum sector $n_0+r\Z$ of a compact boson: each momentum
$n$ contributes $q^{tn^2}$ (its $L_0$ eigenvalue) times the oscillator partition function $1/\eta$.
\end{physics}

\begin{status}
Definition.
\end{status}

\begin{theorem}[$S$-transformation of $\theta_{n_0,r,t}$; GR (3.15)]\label{l22:thm:thetaS}
For $r\ne0$, $t>0$, $\operatorname{Im}\tau>0$:
$\theta_{n_0,r,t}(-1/\tau)=\bigl(-i\tau/(2tr^2)\bigr)^{1/2}\,\vartheta\!\left(n_0/r,\ \tau/(2tr^2)\right)$, with
$\vartheta(z,\tau)=\sum_n e^{2\pi i nz+\pi i n^2\tau}$.
\end{theorem}
\noindent File \leanpath{VOA/\allowbreak Classification/\allowbreak COne/\allowbreak RankOneTheta.lean}.
\begin{leancode}
theorem rankOneTheta_S (n₀ : ℤ) {r : ℤ} (hr : r ≠ 0) {t : ℝ} (ht : 0 < t) {τ : ℂ}
    (hτ : 0 < τ.im) :
    rankOneTheta n₀ r t (-1 / τ) =
      (-I * (τ / (2 * t * (r : ℂ) ^ 2))) ^ (1 / 2 : ℂ) *
        jacobiTheta₂ ((n₀ : ℂ) / r) (τ / (2 * t * (r : ℂ) ^ 2)) :=
\end{leancode}

\begin{dictionary}
\lean{jacobiTheta₂ z τ} is Mathlib's $\sum_n\exp(2\pi i n z+\pi i n^2\tau)$; \lean{(…) ^ (1 / 2 : ℂ)} is the
principal branch of the square root, the correct one since $-i\tau/(2tr^2)$ has positive real part on
$\ltwotwoH$. The library derives the forms used later (\lean{rankOneTheta_S_residues},
\lean{rankOneTheta_zero_one_S}) by expanding the right-hand side over residues.
\end{dictionary}

\begin{proof}[Proof on paper]
Poisson summation: write $\theta_{n_0,r,t}(\tau)=\sum_m g(m)$ with $g(x)=e^{2\pi it(n_0+rx)^2\tau}$, whose Fourier
transform is again a Gaussian. Equivalently, apply Jacobi's transformation formula
$\vartheta(z/\tau,-1/\tau)=(-i\tau)^{1/2}e^{\pi iz^2/\tau}\vartheta(z,\tau)$ (in Mathlib) after completing the square.
\end{proof}

\begin{physics}
Exchanging space and time on the torus turns a sum over momenta into a sum over windings; this is the
$T$-duality-like Poisson resummation of the free boson partition function.
\end{physics}

\begin{status}
Proved, Lean's three axioms.
\end{status}

\subsection{Virasoro multiplicities and positivity of the theta weights}

\begin{definition}[The $c=1$ Virasoro characters; GR (3.4)--(3.7)]\label{l22:def:etavir}
At $c=1$ the irreducible Virasoro module $L(1,h)$ with $h\in\N$ has $\eta\chi_h=q^{h}-q^{(\sqrt h+1)^2}$ if $h$ is a
perfect square, and $\eta\chi_h=q^h$ otherwise (up to the common factor $q^{-1/24}$). For multiplicities
$A:\N\to\N$, $\sum_hA_h\,\eta\chi_h$ is the formal series whose $N$-th coefficient is the finite sum over $h\le N$.
\end{definition}
\noindent File \leanpath{VOA/\allowbreak Classification/\allowbreak COne/\allowbreak WeightPositivity.lean}.
\begin{leancode}
def etaVirasoroSum (A : ℕ → ℕ) : PowerSeries ℤ :=
  mk fun N => ∑ h ∈ range (N + 1), (A h : ℤ) * coeff N (etaVirasoroChar h)
\end{leancode}

\begin{dictionary}
\lean{etaVirasoroChar h} is $\eta\chi_h$ as an integer power series (defined just above in the same file);
\lean{coeff_etaVirasoroChar_of_lt} proves that its coefficients below $q^h$ vanish, which justifies truncating
the infinite sum at $h\le N$. Only integer $h$ occur because the weights of a VOA are integers; at $c=1$ the
degenerate weights are $h=\mu^2/4$, and for integer $h$ only integer $\mu/2$ occur.
\end{dictionary}

\begin{physics}
At $c=1$ the vacuum module of any chiral algebra decomposes into Virasoro representations; the degenerate ones
($h=j^2$) have a null vector at level $2j+1$, which is why their characters lose the term $q^{(j+1)^2}$. These are
the $SU(2)$ spin-$j$ multiplets of the free boson at the self-dual radius.
\end{physics}

\begin{status}
Definition.
\end{status}

\begin{theorem}[GR Lemma~3.11: the theta weights are non-negative]\label{l22:thm:wnonneg}
Suppose a power series $F$ equals $\sum_hA_h\eta\chi_h$ with $A_h\in\N$, $A_0=1$, and also
$F=\sum_{t\in T}\sum_{j\bmod r_t}a_t(j)\,\theta_{j,r_t,t}$ with all $t$ square-free. Then for each $t\in T$,
$w_t:=(r_t\sqrt{2t})^{-1}\sum_ja_t(j)$ is a non-negative real number.
\end{theorem}
\noindent File \leanpath{VOA/\allowbreak Classification/\allowbreak COne/\allowbreak WeightPositivity.lean}.
\begin{leancode}
theorem theta_weight_nonneg (F : PowerSeries ℂ) (A : ℕ → ℕ) (hA0 : A 0 = 1)
    (hVir : F = map (Int.castRingHom ℂ) (etaVirasoroSum A))
    (T : Finset ℕ) (hT : ∀ t ∈ T, Squarefree t) (r : ℕ → ℕ) [∀ t, NeZero (r t)]
    (a : (t : ℕ) → ZMod (r t) → ℂ)
    (hθ : F = ∑ t ∈ T, ∑ j : ZMod (r t), a t j • map (Int.castRingHom ℂ) (thetaSeries (r t) j t))
    {t : ℕ} (ht : t ∈ T) :
    ∃ x : ℝ, 0 ≤ x ∧
      (1 / ((r t : ℂ) * (Real.sqrt (2 * t) : ℂ))) * ∑ j : ZMod (r t), a t j = x :=
\end{leancode}

\begin{dictionary}
\lean{map (Int.castRingHom ℂ)} views an integer series as a complex series; \lean{thetaSeries (r t) j t} is the
formal $q$-expansion of $\theta_{j,r_t,t}$; \lean{a t j • …} is scalar multiplication. \lean{[∀ t, NeZero (r t)]}
says every modulus is nonzero (so \lean{ZMod (r t)} is finite). The conclusion \lean{∃ x : ℝ, 0 ≤ x ∧ … = x}
says the complex number $w_t$ is real and $\ge0$. Hypotheses are exactly GR's.
\end{dictionary}

\begin{proof}[Proof on paper]
Compare coefficients at $q^{tn^2}$ with $t$ square-free. By Definition~\ref{l22:def:etavir}, the coefficient of
$q^{(n+1)^2}$ in $\sum A_h\eta\chi_h$ is $A_{(n+1)^2}-A_{n^2}$, and the coefficient of a non-square $N$ is $A_N\ge0$.
On the theta side, the coefficient of $q^{tn^2}$ collects $a_t(j)$ for $j\equiv\pm n$. Summing over a full period
of $n$ and using $A\ge0$ (and $A_0=1$ for the telescoping square terms when $t=1$) gives
$\sum_ja_t(j)$ as half a sum of non-negative integers.
\end{proof}

\begin{physics}
The total weight $w_t$ of the momentum lattice $\sqrt{2t}\Z$ in the vacuum character is a count of Virasoro
primaries, hence non-negative.
\end{physics}

\begin{status}
Proved, Lean's three axioms.
\end{status}

\subsection{Uniqueness of the surviving theta weight}

\begin{theorem}[GR Prop.~3.12, arithmetic core]\label{l22:thm:unique}
Let $8\mid C$, $T$ a finite set of square-free $t$ with $4t\mid C$, $w_t\ge0$ with $\sum_tw_t=S_{00}>0$, and
suppose $\bigl|\sum_t\bigl(\tfrac{2t}{\ell}\bigr)w_t\bigr|\ge S_{00}$ for every $\ell$ coprime to $C$. Then exactly one
$w_t$ is nonzero.
\end{theorem}
\noindent File \leanpath{VOA/\allowbreak Classification/\allowbreak COne/\allowbreak ModularDataLemmas.lean}.
\begin{leancode}
theorem existsUnique_weight_ne_zero {C : ℕ} (hC0 : C ≠ 0) (hC8 : 8 ∣ C) (T : Finset ℕ)
    (hT : ∀ t ∈ T, Squarefree t ∧ 4 * t ∣ C) (w : ℕ → ℝ) (hw : ∀ t ∈ T, 0 ≤ w t) {S₀₀ : ℝ}
    (hS₀₀ : 0 < S₀₀) (hsum : ∑ t ∈ T, w t = S₀₀)
    (hbound : ∀ ℓ : ℕ, ℓ.Coprime C →
      S₀₀ ≤ |∑ t ∈ T, (jacobiSym (2 * (t : ℤ)) ℓ : ℝ) * w t|) :
    ∃! t, t ∈ T ∧ w t ≠ 0 :=
\end{leancode}

\begin{dictionary}
\lean{jacobiSym (2 * t) ℓ} is the Jacobi symbol $\bigl(\tfrac{2t}{\ell}\bigr)\in\{-1,0,1\}$; \lean{ℓ.Coprime C}
is $\gcd(\ell,C)=1$. \lean{∃! t, t ∈ T ∧ w t ≠ 0}: there is a unique such $t$. GR assume $24\mid C$; the theorem
needs only $8\mid C$, so it is slightly \emph{stronger}.
\end{dictionary}

\begin{proof}[Proof on paper]
Suppose $w_{t_1},w_{t_2}>0$ for $t_1\ne t_2$. Since they are square-free and distinct, $2t_1$ and $2t_2$ differ as
square classes, so by quadratic reciprocity and the Chinese remainder theorem there is $\ell$ coprime to $C$ with
$\bigl(\tfrac{2t_1}{\ell}\bigr)=1$ and $\bigl(\tfrac{2t_2}{\ell}\bigr)=-1$. Then
$|\sum_t(\tfrac{2t}{\ell})w_t|<\sum_tw_t=S_{00}$, since the two terms partially cancel, contradicting the bound.
No Dirichlet theorem is needed.
\end{proof}

\begin{physics}
The Galois symmetry of the modular data acts on the momentum sectors by Jacobi symbols. A mixture of two
different radii would be broken by some Galois conjugation, while unitarity-type bounds forbid that.
\end{physics}

\begin{status}
Proved, Lean's three axioms.
\end{status}

\subsection{The Tannakian squeeze}

\begin{theorem}[GR Theorem~2.18, matrix part]\label{l22:thm:squeeze}
Let $S$ be a finite matrix with vacuum label $0$, $d_a>0$, $d_0=1$, $S_{a0}=S_{00}d_a$, $|S_{ab}|\le S_{00}d_ad_b$, and
positive functions $\ltwotwoch_b(t)$ with $e^{-2\pi tc/24}\ltwotwoch_0(t)\to1$ and $e^{-2\pi tc/24}\ltwotwoch_b(t)\to0$ for
$b\ne0$. If $\sum_b\bigl(DS_{00}d_b-\sum_ad_aS_{ab}\bigr)\ltwotwoch_b(t)=0$ for all $t>0$, then $D=\sum_ad_a^2$ and
$S_{ab}=S_{00}d_ad_b$ for all $a,b$.
\end{theorem}
\noindent File \leanpath{VOA/\allowbreak Classification/\allowbreak COne/\allowbreak ModularDataLemmas.lean}.
\begin{leancode}
theorem tannakian_squeeze {ι : Type*} [Fintype ι] (i₀ : ι) (S : ι → ι → ℂ)
    (d : ι → ℝ) (S₀₀ c : ℝ) (D : ℂ) (ch : ι → ℝ → ℝ)
    (hd : ∀ a, 0 < d a) (hd₀ : d i₀ = 1) (hS₀₀ : 0 < S₀₀)
    (hSa0 : ∀ a, S a i₀ = S₀₀ * d a)
    (hbound : ∀ a b, ‖S a b‖ ≤ S₀₀ * d a * d b)
    (hch : ∀ b, ∀ t, 0 < t → 0 < ch b t)
    (hlead₀ : Tendsto (fun t => Real.exp (-(2 * Real.pi * t * c / 24)) * ch i₀ t) atTop (𝓝 1))
    (hlead : ∀ b, b ≠ i₀ →
      Tendsto (fun t => Real.exp (-(2 * Real.pi * t * c / 24)) * ch b t) atTop (𝓝 0))
    (hid : ∀ t, 0 < t →
      ∑ b, (D * S₀₀ * d b - ∑ a, d a * S a b) * (ch b t : ℂ) = 0) :
    D = ((∑ a, d a ^ 2 : ℝ) : ℂ) ∧ ∀ a b, S a b = S₀₀ * d a * d b :=
\end{leancode}

\begin{dictionary}
\lean{Tendsto f atTop (𝓝 x)}: $f(t)\to x$ as $t\to\infty$. The identity \lean{hid} is GR (2.33) evaluated at
$\tau=it$. The $S$-matrix need not be symmetric or unitary --- the theorem is slightly \emph{more general} than
GR's.
\end{dictionary}

\begin{proof}[Proof on paper]
Multiply the identity by $e^{-2\pi tc/24}$ and let $t\to\infty$: only the vacuum survives, so its coefficient
vanishes, $DS_{00}=\sum_ad_aS_{a0}=S_{00}\sum_ad_a^2$, i.e.\ $D=\sum d_a^2$. Then each coefficient
$DS_{00}d_b-\sum_ad_aS_{ab}\ge DS_{00}d_b-\sum_a d_a\cdot S_{00}d_ad_b=0$ by the bound, and a positive combination of
positive functions vanishing identically forces each coefficient to vanish, and equality in the bound for every
$a$: $S_{ab}=S_{00}d_ad_b$.
\end{proof}

\begin{physics}
$S_{ab}=S_{00}d_ad_b$ says the braiding among these anyons is trivial (they are bosons that braid trivially): the
subcategory is symmetric, so by Deligne it is $\operatorname{Rep}G$ and the anyons can be condensed (the
symmetry ``ungauged'').
\end{physics}

\begin{status}
Proved, Lean's three axioms.
\end{status}

\subsection{The theta Galois law}

\begin{definition}[The theta Galois law; GR Prop.~3.8, (3.18)]\label{l22:def:galoislaw}
For $24\mid M$, $4tr^2\mid M$, $\ell$ coprime to $M$, $\ell\ell'\equiv1$ and $\gamma\equiv\operatorname{diag}(\ell,\ell')
\pmod M$: $\theta_{n_0,r,t}(\gamma\tau)\,\eta(\tau)=\bigl(\tfrac{2t}{\ell}\bigr)\theta_{\ell n_0,r,t}(\tau)\,\eta(\gamma\tau)$.
\end{definition}
\noindent File \leanpath{VOA/\allowbreak Classification/\allowbreak COne/\allowbreak GapsPaper.lean}.
\begin{leancode}
def ThetaGaloisLaw : Prop :=
  ∀ (n₀ : ℤ) (r t M : ℕ), 0 < r → 0 < t → 4 * t * r ^ 2 ∣ M → 24 ∣ M →
    ∀ ℓ : ℕ, ℓ.Coprime M → ∀ ℓ' : ℤ, (ℓ : ℤ) * ℓ' ≡ 1 [ZMOD M] →
    ∀ γ : Matrix.SpecialLinearGroup (Fin 2) ℤ,
      γ 0 0 ≡ ℓ [ZMOD M] → γ 1 1 ≡ ℓ' [ZMOD M] → γ 0 1 ≡ 0 [ZMOD M] → γ 1 0 ≡ 0 [ZMOD M] →
      ∀ τ : ℂ, 0 < τ.im →
        rankOneTheta n₀ r t (mobius γ τ) * ModularForm.eta τ =
          (jacobiSym (2 * (t : ℤ)) ℓ : ℂ) * rankOneTheta ((ℓ : ℤ) * n₀) r t τ *
            ModularForm.eta (mobius γ τ)
\end{leancode}

\begin{dictionary}
Cross-multiplied, again, to avoid dividing by $\eta$. \lean{ModularForm.eta} is Mathlib's Dedekind $\eta$. For
$\ell=1$ this is the $\Gamma(M)$-invariance of $\theta/\eta$. The docstring records a numerical check and,
importantly, that the plausible variant with $\ell^{-1}n_0$ in place of $\ell n_0$ is \emph{false} for $r=7$.
\end{dictionary}

\begin{physics}
The lattice boson's characters transform under the congruence quotient by the Weil representation, in which
diagonal matrices act by Jacobi symbols and by scaling the momentum label.
\end{physics}

\begin{status}
A \lean{Prop}, \textbf{proved}: \lean{thetaGaloisLaw : ThetaGaloisLaw} in
\leanpath{VOA/\allowbreak Classification/\allowbreak COne/\allowbreak WeilKernel.lean}, through the Weil
representation and Gauss's sign theorem (\lean{WeilGalois.gaussG_one}).
\end{status}

%======================================================================
\section{The conditional theorem}\label{l22:sec:conditional}
%======================================================================

\subsection{Four gaps}

Besides the seven published inputs, the first conditional proof named four \emph{gaps}: steps of GR §3 or of
classical analysis that were true but not yet formalized. They are recorded here because they are hypotheses of
the conditional theorem; all four are now proved.

\begin{definition}[The four gaps]\label{l22:def:gaps}
(i) \textbf{GapQExpansion}: uniqueness of $q$-expansions on $\ltwotwoH$; multiplication of a $q$-series by $\eta$
corresponds to the formal product with Euler's series; the $q$-expansion of $\theta_{j,r,t}$.
(ii) \textbf{GapLemma37} (GR Lemma~3.7): for nice $c=1$ $V$ with $\Gamma(N)$-invariant characters,
$F_V=\eta\cdot\ltwotwoch_V$ is a holomorphic $\theta$-type weight-$\frac12$ form for $\Gamma_1(N)$.
(iii) \textbf{GapThetaReduction} (GR (3.13)): a combination of $\theta_{n_0,r,t}$ can be rewritten with square-free
$t$ and one modulus $R_t$ per $t$.
(iv) \textbf{GapGaloisBound} (GR (3.31)--(3.35)): under Galois symmetry,
$|\sum_t(\tfrac{2t}{\ell})w_t|\ge S_{00}$ for every $\ell$ coprime to $N$.
\end{definition}
\noindent File \leanpath{VOA/\allowbreak Classification/\allowbreak COne/\allowbreak Conditional.lean}.
\begin{leancode}
structure GapQExpansion : Prop where
  unique : ∀ (c c' : ℕ → ℂ) (f : ℂ → ℂ),
    (∀ τ : ℂ, 0 < τ.im → HasSum (fun n : ℕ => c n * cexp (2 * Real.pi * I * n * τ)) (f τ)) →
    (∀ τ : ℂ, 0 < τ.im → HasSum (fun n : ℕ => c' n * cexp (2 * Real.pi * I * n * τ)) (f τ)) →
    c = c'
  eta_mul : ∀ (c : ℕ → ℂ) (g : ℂ → ℂ),
    (∀ τ : ℂ, 0 < τ.im →
      HasSum (fun n : ℕ => c n * cexp (2 * Real.pi * I * ((n : ℂ) - 1 / 24) * τ)) (g τ)) →
    ∀ τ : ℂ, 0 < τ.im →
      HasSum (fun n : ℕ => PowerSeries.coeff n
        (PowerSeries.map (Int.castRingHom ℂ) eulerSeries * PowerSeries.mk c) *
          cexp (2 * Real.pi * I * n * τ)) (ModularForm.eta τ * g τ)
  theta : ∀ (r : ℕ) [NeZero r] (j : ZMod r) (t : ℕ), 0 < t → ∀ τ : ℂ, 0 < τ.im →
    HasSum (fun n : ℕ => (thetaCoeff r j t n : ℂ) * cexp (2 * Real.pi * I * n * τ))
      (rankOneTheta (j.val : ℕ) r t τ)
\end{leancode}
\noindent File \leanpath{VOA/\allowbreak Classification/\allowbreak COne/\allowbreak Conditional.lean}.
\begin{leancode}
def GapLemma37 : Prop :=
  ∀ (𝒱 : Type u) [AddCommGroup 𝒱] [Module ℂ 𝒱] (V : VertexOperatorAlgebra 𝒱), Nice V →
    V.centralCharge = 1 → ∀ D : ModularData V, D.IsModularData → ∀ N : ℕ, D.IsLevel N →
      (∀ γ ∈ CongruenceSubgroup.Gamma N, ∀ a (τ : ℂ), 0 < τ.im →
        D.ch a (mobius γ τ) = D.ch a τ) →
      IsThetaTypeForm N (fun τ => ModularForm.eta τ * D.ch D.vac τ)
\end{leancode}
\noindent File \leanpath{VOA/\allowbreak Classification/\allowbreak COne/\allowbreak Conditional.lean}.
\begin{leancode}
def GapThetaReduction : Prop :=
  ∀ (N : ℕ), 0 < N → ∀ (s : Finset (ℤ × ℕ × ℕ)) (coef : ℤ × ℕ × ℕ → ℂ),
    (∀ p ∈ s, 1 ≤ p.2.1 ∧ 1 ≤ p.2.2 ∧ 4 * p.2.2 * p.2.1 ^ 2 ∣ N) →
    ∃ (T : Finset ℕ) (R : ℕ → ℕ+) (a : (t : ℕ) → ZMod (R t) → ℂ),
      (∀ t ∈ T, Squarefree t ∧ 4 * t * (R t : ℕ) ^ 2 ∣ N) ∧
      ∀ τ : ℂ, 0 < τ.im → ∑ p ∈ s, coef p * rankOneTheta p.1 p.2.1 p.2.2 τ =
        ∑ t ∈ T, ∑ j : ZMod (R t), a t j * rankOneTheta (j.val : ℕ) (R t : ℕ) t τ
\end{leancode}
\noindent File \leanpath{VOA/\allowbreak Classification/\allowbreak COne/\allowbreak Conditional.lean}.
\begin{leancode}
def GapGaloisBound : Prop :=
  ∀ (𝒱 : Type u) [AddCommGroup 𝒱] [Module ℂ 𝒱] (V : VertexOperatorAlgebra 𝒱), Nice V →
    V.centralCharge = 1 → ∀ D : ModularData V, D.IsModularData → ∀ N : ℕ, D.IsLevel N →
      D.GaloisSymmetric N →
      ∀ (T : Finset ℕ) (R : ℕ → ℕ+) (a : (t : ℕ) → ZMod (R t) → ℂ),
        (∀ t ∈ T, Squarefree t ∧ 4 * t * (R t : ℕ) ^ 2 ∣ N) →
        (∀ τ : ℂ, 0 < τ.im → ModularForm.eta τ * D.ch D.vac τ =
          ∑ t ∈ T, ∑ j : ZMod (R t), a t j * rankOneTheta (j.val : ℕ) (R t : ℕ) t τ) →
        ∀ ℓ : ℕ, ℓ.Coprime N →
          D.S00 ≤ ‖∑ t ∈ T, (jacobiSym (2 * (t : ℤ)) ℓ : ℂ) *
            ((1 / ((R t : ℂ) * (Real.sqrt (2 * t) : ℂ))) * ∑ j : ZMod (R t), a t j)‖

\end{leancode}

\begin{dictionary}
\lean{ℕ+} is the type of positive natural numbers, so \lean{R t : ℕ+} is a positive modulus. \lean{Squarefree t}
is Mathlib's square-freeness. In \lean{GapGaloisBound} the quantity
\lean{(1 / ((R t : ℂ) * (Real.sqrt (2 * t) : ℂ))) * ∑ j : ZMod (R t), a t j} is the weight $w_t$ of
Theorem~\ref{l22:thm:wnonneg}, and the conclusion (the last two lines of the quotation) is the bound
$S_{00}\le\bigl|\sum_t\bigl(\tfrac{2t}{\ell}\bigr)w_t\bigr|$, exactly GR (3.35).
\end{dictionary}

\begin{physics}
(ii) is the statement that the vacuum character of a $c=1$ theory times $\eta$ behaves like a one-boson theta
function under the congruence subgroup; (iv) is the modular-invariance consistency of the Galois action with
positivity of quantum dimensions.
\end{physics}

\begin{status}
\lean{Prop}s. All proved: \lean{gapQExpansion}, \lean{gapThetaReduction} (in
\leanpath{GapsClassical.lean}), and \lean{gapLemma37_of_thetaGaloisLaw},
\lean{gapGaloisBound_of_thetaGaloisLaw} (in \leanpath{GapsPaper.lean}) together with \lean{thetaGaloisLaw}.
\end{status}

\subsection{The last step and the conditional theorem}

\begin{theorem}[Orbifold transport]\label{l22:thm:transport}
If $W\cong W'$ and $G\le\Aut W$, then the conjugate subgroup $G'\le\Aut W'$ satisfies $G\cong G'$ and
$W^G\cong W'^{G'}$.
\end{theorem}
\noindent File \leanpath{VOA/\allowbreak Classification/\allowbreak COne/\allowbreak Conditional.lean}.
\begin{leancode}
theorem orbifold_transport {W : VertexOperatorAlgebra 𝒱} {W' : VertexOperatorAlgebra 𝒲}
    (h : VOAIsomorphic W W') (G : Subgroup (Aut W)) :
    ∃ G' : Subgroup (Aut W'), Nonempty (G ≃* G') ∧
      VOAIsomorphic (orbifoldVOA G) (orbifoldVOA G') :=
\end{leancode}

\begin{dictionary}
The conjugate subgroup is the image of $G$ under $g\mapsto e g e^{-1}$, where $e:W\to W'$ is the isomorphism.
\end{dictionary}

\begin{proof}[Proof on paper]
$e$ maps $G$-fixed vectors bijectively to $eGe^{-1}$-fixed vectors and preserves products, vacuum and $\omega$.
\end{proof}

\begin{physics}
Gauging the same symmetry of the same theory, described in two different bases, gives the same gauged theory.
\end{physics}

\begin{status}
Proved, Lean's three axioms.
\end{status}

\begin{theorem}[The last step of GR's proof]\label{l22:thm:lastStep}
Assume \textup{IDongMason} and \textup{IEnumeration}. If $V\cong W^G$ with $W$ nice, $c=1$, $W_1\ne0$ and $G$
finite, then $V$ is in the list.
\end{theorem}
\noindent File \leanpath{VOA/\allowbreak Classification/\allowbreak COne/\allowbreak Conditional.lean}.
\begin{leancode}
theorem inCOneList_of_extension (hDM : IDongMason.{w}) (hEn : IEnumeration)
    {𝒱 : Type v} [AddCommGroup 𝒱] [Module ℂ 𝒱] {V : VertexOperatorAlgebra 𝒱}
    {𝒲 : Type w} [AddCommGroup 𝒲] [Module ℂ 𝒲] {W : VertexOperatorAlgebra 𝒲}
    (G : Subgroup (Aut W)) (hG : Finite G) (hW : Nice W) (hc : W.centralCharge = 1)
    (h1 : finrank ℂ (W.grading 1) ≠ 0) (hV : VOAIsomorphic V (orbifoldVOA G)) :
    InCOneList V :=
\end{leancode}

\begin{proof}[Proof on paper]
Dong--Mason gives $W\cong V_{2m}$; transport moves $G$ to $G'\le\Aut V_{2m}$ with $W^G\cong V_{2m}^{G'}$; the
enumeration puts $V_{2m}^{G'}$ in the list, and the list is closed under isomorphism (\lean{inCOneList_of_iso}).
\end{proof}

\begin{physics}
Every nice $c=1$ theory is a gauging of a free boson at a rational radius by a finite symmetry group.
\end{physics}

\begin{status}
Proved, conditional on the two displayed inputs; Lean's three axioms.
\end{status}

\begin{theorem}[Gannon--Rayhaun Theorem~3.1, forward, conditionally]\label{l22:thm:conditional}
The seven inputs and the four gaps imply the forward classification.
\end{theorem}
\noindent File \leanpath{VOA/\allowbreak Classification/\allowbreak COne/\allowbreak Conditional.lean}.
\begin{leancode}
theorem gannonRayhaun_of_inputs
    (hMod : IModData.{u}) (hCG : ICongruenceGalois.{u}) (hSS : ISerreStark24)
    (hVir : IVirasoro.{u}) (hTan : ITannakian.{u}) (hDM : IDongMason.{u})
    (hEn : IEnumeration)
    (g37 : GapLemma37.{u}) (gRed : GapThetaReduction) (gQ : GapQExpansion)
    (gGal : GapGaloisBound.{u}) :
    GannonRayhaunStatement.{u} := by
\end{leancode}

\begin{dictionary}
A theorem with eleven \lean{Prop} hypotheses and conclusion \lean{GannonRayhaunStatement.{u}}; all inputs are
taken in the same universe \lean{u}.
\end{dictionary}

\begin{proof}[Proof on paper, following GR §3]
Let $V$ be nice with $c=1$, with modular data $D$ (IModData) and a level $N$ divisible by $24$.
\begin{enumerate}
\item \emph{The character is a theta-type form.} The characters are $\Gamma(N)$-invariant (ICongruenceGalois), so
  $F_V=\eta\,\ltwotwoch_V$ is a holomorphic $\theta$-type form of weight $\frac12$ for $\Gamma_1(N)$ (GapLemma37):
  $\eta$ supplies the multiplier of $\theta$ on $\Gamma(24)$, $T$-invariance comes from integrality of the weights
  of $V$, and holomorphy at the cusps from the $S$-law and $h_a\ge0$.
\item \emph{Serre--Stark.} Hence $F_V=\sum_{t\in T}\sum_{j\bmod R_t}a_t(j)\,\theta_{j,R_t,t}$ with $t$ square-free
  (ISerreStark24, GapThetaReduction).
\item \emph{Lemma 3.10.} Apply $S$ and project onto integer powers of $q$ (an average over the
  translations $T^k$). By the $S$-transformation of the theta functions (Theorem~\ref{l22:thm:thetaS}) and the $S$-law, the projection of $\ltwotwoch_V(-1/\tau)$ is a combination of the
  characters of the integer-weight modules, with coefficients $S_{0a}=S_{00}d_a$, and on the theta side it is a
  combination of the $\theta_{0,1,t}$:
  \begin{equation*}
  \sum_{h_a\in\Z}d_a\,\ltwotwoch_a(\tau)=\frac{1}{S_{00}\,\eta(\tau)}\sum_{t\in T}w_t\,\theta_{0,1,t}(\tau),
  \end{equation*}
  with $w_t$ as in Theorem~\ref{l22:thm:wnonneg}.
\item \emph{Lemma 3.11 and the weight sum.} IVirasoro and the positivity of the theta weights (Theorem~\ref{l22:thm:wnonneg}) give $w_t\ge0$; comparing leading terms gives
  $\sum_tw_t=S_{00}$.
\item \emph{Prop.~3.12.} The Galois bound (GapGaloisBound) and the uniqueness of the surviving theta weight (Theorem~\ref{l22:thm:unique}) leave exactly one $t_\star$, so
  $\sum_{h_a\in\Z}d_a\ltwotwoch_a=\theta_{0,1,t_\star}/\eta$ --- the character of $V_{2t_\star}$.
\item \emph{Theorem 2.18.} This vacuum identity, the theta $S$-transformation and the Tannakian squeeze (Theorem~\ref{l22:thm:squeeze}) give $S_{ab}=S_{00}d_ad_b$ on integer-weight labels.
\item \emph{Tannakian step.} ITannakian produces a nice $W$ with $V\cong W^G$, $G$ finite, and
  $\dim W_n=\sum_{h_a\in\Z}d_a\dim(M_a)_n$, i.e.\ $\ltwotwoch_W=\theta_{0,1,t_\star}/\eta$, so
  $\dim W_1\ge1$.
\item \emph{Last step.} Theorem~\ref{l22:thm:lastStep}.
\end{enumerate}
\end{proof}

\begin{physics}
The argument reads: the $c=1$ theory's integer-spin anyons, weighted by quantum dimension, have exactly the partition
function of a single compact boson; Galois symmetry forbids a superposition of radii; so the integer-spin anyons
form a gaugeable $\operatorname{Rep}G$, condensing them gives a theory with a current, hence a free boson at a
rational radius, and the original theory is its $G$-gauging.
\end{physics}

\begin{status}
Proved, \lean{#print axioms} = \lean{[propext, Classical.choice, Quot.sound]} (machine-checked on 2026-10-13).
\textbf{Conditional} on the eleven displayed hypotheses.
\end{status}

%======================================================================
\section{Discharging the inputs: the unconditional forward theorem}\label{l22:sec:discharge}
%======================================================================

\subsection{How the discharge is organized}

The conditional theorem was followed by a chain of theorems, each replacing one hypothesis by a proof or by
finer hypotheses (files \leanpath{Main*.lean}, \leanpath{SerreStark*.lean}, \leanpath{QExp*.lean} in
\leanpath{VOA/\allowbreak Classification/\allowbreak COne/}, and \leanpath{HKL*.lean}, \leanpath{DLN*.lean},
\leanpath{NS*.lean} in \leanpath{VOA/\allowbreak Huang/}). The chain ends at a hypothesis-free theorem. We list
what discharges each hypothesis; the deep inputs are explained in Chapter~\ref{ch:L21}.

\begin{center}
\begin{longtable}{p{0.24\textwidth}p{0.70\textwidth}}
\toprule
Hypothesis & Discharged by \\
\midrule
\lean{GapQExpansion} & \lean{gapQExpansion} (Mathlib power-series uniqueness, Euler product, theta regrouping). \\
\lean{GapThetaReduction} & \lean{gapThetaReduction} (reindexing of theta series). \\
\lean{GapLemma37}, \lean{GapGaloisBound} & \lean{paperGaps_of_thetaGaloisLaw} with \lean{thetaGaloisLaw}
  (Weil representation, \lean{weilKernelLaw}, Gauss sums). \\
\lean{ISerreStark24} & \lean{SSHecke.iSerreStark24_of_boundedCoeff_of_levelBox} with
  \lean{SSHecke.serreStarkBoundedCoeff24_holds} and \lean{SSHecke.serreStarkLevelBox24}: a proof of Serre--Stark at
  levels $24\mid N$ via Hecke operators, a Fricke involution and the $q$-expansion principle in weight two
  (\lean{SSHecke.qExpPrinciple24Away_two}). \\
\lean{IVirasoro} & \lean{iVirasoro_of_composition} with \lean{virLocalCompositionSeries} (weight-by-weight
  Jordan--Hölder for positive-energy Virasoro modules) and \lean{virC1DegenerateCharacters_holds} (Feigin--Fuchs at
  $c=1$ via a radical-level bound). \\
\lean{IDongMason} & \lean{iDongMason_of_modData}: trivial Heisenberg commutant from the modular data, then
  \lean{heisLatticeReconstruction} (a VOA containing a rank-one Heisenberg with trivial commutant is a lattice VOA). \\
\lean{IEnumeration} & \lean{iEnumeration}: classification of finite subgroups of $SO(3)$
  (\lean{FiniteRotation.classification}), $\Aut V_2$ and $\Aut V_{2m}$, and the orbifold identifications. \\
\lean{IModData} & \lean{VOA.Huang.HKLRational.iModData_holds}: Zhu, Anderson--Moore, Huang's Verlinde theorem
  (Chapter~\ref{ch:L21}); uses \lean{VOA.pbw} (a theorem). \\
\lean{ICongruenceGalois} & \lean{iCongruenceGalois_of_split} with the Dong--Lin--Ng inputs, obtained from the
  Ng--Schauenburg congruence theorem \lean{nsProjectiveKernelAll_holds} and Beyl's Schur-multiplier facts for
  $\mathrm{SL}_2(\Z/p^k)$: \lean{threePowerSchur}, \lean{twoPowerInflationFromTwo_holds},
  \lean{largePrimeGaussLowerRelations} (Chapter~\ref{ch:L21}). \\
\lean{ITannakian} & \lean{iTannakian_of_rigid} with Deligne's theorem (\lean{symmetricFiberFunctorExists_holds},
  \lean{fiberFunctorReconstruction}) and the Huang--Kirillov--Lepowsky extension theory
  (\lean{hklRationalInput_holds}, \lean{hklLocalityInput_holds}; Chapter~\ref{ch:L21}). \\
\bottomrule
\end{longtable}
\end{center}

\begin{warning}[title={Not every input has a standalone theorem}]
For \lean{IVirasoro}, \lean{ISerreStark24}, \lean{ICongruenceGalois} and \lean{ITannakian} the library contains no
declaration of the literal form \lean{theorem … : IVirasoro}; the corresponding instances are built inside the
assembly chain (e.g.\ \lean{gannonRayhaun_of_iTannakian} builds the Serre--Stark instance with a \lean{let}). This
is harmless: the chain ends in a theorem with \emph{no} hypotheses, so every input it consumed has been proved. A
referee who wants a named statement can read off the composition from the table.
\end{warning}

\subsection{The discharging theorems quoted}

We quote the statements that close the analytic and arithmetic inputs; the Huang-side ones are quoted in
Chapter~\ref{ch:L21}.

\noindent File \leanpath{VOA/\allowbreak Classification/\allowbreak COne/\allowbreak GapsClassical.lean}.
\begin{leancode}
theorem gapQExpansion : GapQExpansion where
\end{leancode}
\noindent File \leanpath{VOA/\allowbreak Classification/\allowbreak COne/\allowbreak GapsClassical.lean}.
\begin{leancode}
theorem gapThetaReduction : GapThetaReduction := by
\end{leancode}
\noindent File \leanpath{VOA/\allowbreak Classification/\allowbreak COne/\allowbreak WeilKernel.lean}.
\begin{leancode}
theorem thetaGaloisLaw : ThetaGaloisLaw := thetaGaloisLaw_of_weilKernelLaw weilKernelLaw
\end{leancode}
\noindent File \leanpath{VOA/\allowbreak Classification/\allowbreak COne/\allowbreak GapsPaper.lean}.
\begin{leancode}
theorem paperGaps_of_thetaGaloisLaw (hL : ThetaGaloisLaw) :
    GapLemma37.{u} ∧ GapGaloisBound.{u} :=
  ⟨gapLemma37_of_thetaGaloisLaw hL, gapGaloisBound_of_thetaGaloisLaw hL⟩
\end{leancode}
\noindent File \leanpath{VOA/\allowbreak Classification/\allowbreak COne/\allowbreak SerreStarkLevelBox.lean}.
\begin{leancode}
theorem iSerreStark24_of_boundedCoeff_of_levelBox (hB : SerreStarkBoundedCoeff24)
    (hL : SerreStarkLevelBox24) : ISerreStark24 :=
\end{leancode}
\noindent File \leanpath{VOA/\allowbreak Classification/\allowbreak COne/\allowbreak SerreStarkFrickeBeta.lean}.
\begin{leancode}
theorem serreStarkLevelBox24 : SerreStarkLevelBox24 :=
\end{leancode}
\noindent File \leanpath{VOA/\allowbreak Classification/\allowbreak COne/\allowbreak QExpSpan2Headline.lean}.
\begin{leancode}
theorem SSHecke.serreStarkBoundedCoeff24_holds : SSHecke.SerreStarkBoundedCoeff24 :=
  SSHecke.boundedCoeff24_of_qExpPrinciple24Away_two SSHecke.qExpPrinciple24Away_two
\end{leancode}
\noindent File \leanpath{VOA/\allowbreak Classification/\allowbreak COne/\allowbreak VirasoroCOne.lean}.
\begin{leancode}
theorem iVirasoro_of_composition (hJH : VirLocalCompositionSeries.{u})
    (hdeg : VirC1DegenerateCharacters) : IVirasoro.{u} :=
\end{leancode}
\noindent File \leanpath{VOA/\allowbreak Classification/\allowbreak COne/\allowbreak VirComposition.lean}.
\begin{leancode}
theorem virLocalCompositionSeries : VirLocalCompositionSeries.{u} := by
\end{leancode}
\noindent File \leanpath{VOA/\allowbreak Classification/\allowbreak COne/\allowbreak FeiginFuchsLower.lean}.
\begin{leancode}
theorem virC1DegenerateCharacters_holds : VirC1DegenerateCharacters :=
  virC1DegenerateCharacters_of_radicalLevelBound virC1RadicalLevelBound_holds
\end{leancode}
\noindent File \leanpath{VOA/\allowbreak Classification/\allowbreak COne/\allowbreak CommutantTrivial.lean}.
\begin{leancode}
theorem iDongMason_of_modData (hMD : IModData.{u}) : IDongMason.{u} :=
  iDongMason_of_facts (heisCommutantTrivial_of_modData hMD) heisLatticeReconstruction
\end{leancode}
\noindent File \leanpath{VOA/\allowbreak Classification/\allowbreak COne/\allowbreak AutV2Bridge.lean}.
\begin{leancode}
theorem iEnumeration : IEnumeration := by
\end{leancode}
\noindent File \leanpath{VOA/\allowbreak Huang/\allowbreak HKLFrobDimMain.lean}.
\begin{leancode}
theorem iModData_holds : IModData.{u} :=
  VOA.Huang.HuangTierReduced.iModData_of_verlinde_pbw_inputs
    VOA.Huang.AndersonMoore.andersonMooreInput_holds
    (VOA.Huang.VerlindeFusion.verlindeInput_of_fusion huangVerlindeFusionInput_holds)
\end{leancode}

\begin{dictionary}
The statements are exactly ``the named \lean{Prop} holds'' (or, for the two reduction theorems, ``these proved
\lean{Prop}s imply that one''). The bodies shown in one line (\lean{:= by}, \lean{where}) are the start of proofs that
the book does not reproduce. \lean{iModData_holds} is shown with its proof term because it is informative: it is a
composition of Anderson--Moore (\lean{andersonMooreInput_holds}) and Huang's Verlinde theorem
(\lean{huangVerlindeFusionInput_holds}), through \lean{iModData_of_verlinde_pbw_inputs}, whose name records the
dependence on \lean{VOA.pbw}.
\end{dictionary}

\begin{physics}
Each line closes one of the classical pillars of rational CFT used by Gannon and Rayhaun: Serre--Stark (one-boson
partition functions), Feigin--Fuchs (the $c=1$ Virasoro spectrum), Dong--Mason (currents make a boson), Klein's
classification of rotation groups, and Huang's modular tensor category.
\end{physics}

\begin{status}
All proved. Lean's three axioms (re-probed: \lean{iModData_holds} included); \lean{iModData_holds} (and
\lean{iDongMason_of_modData} when instantiated with it) uses the theorem \lean{VOA.pbw}.
\end{status}

\subsection{The forward theorem}

\begin{theorem}[Gannon--Rayhaun Theorem~3.1, forward direction]\label{l22:thm:forward}
Every nice vertex operator algebra of central charge $1$ is isomorphic to some $V_{2m}$ ($m\ge1$), some
$V_{2m}^+$ ($m\ge2$), or a fixed-point algebra $V_2^G$ with $G\cong A_4$, $S_4$ or $A_5$.
\end{theorem}
\noindent File \leanpath{VOA/\allowbreak Classification/\allowbreak COne/\allowbreak MainUnconditional.lean}.
\begin{leancode}
theorem gannonRayhaun_forward : GannonRayhaunStatement.{u} :=
  gannonRayhaun_beyl_large VOA.Huang.DLNSchur.largePrimeGaussLowerRelations
\end{leancode}

\begin{dictionary}
No hypotheses. The proof term is a composition: \lean{gannonRayhaun_beyl_large} is the forward theorem conditional
on the single large-prime Schur residual \lean{LargePrimeGaussLowerRelations}, and
\lean{largePrimeGaussLowerRelations} proves it. The preceding layers are
\lean{gannonRayhaun_beyl_fine} (three Beyl residuals), \lean{gannonRayhaun_beyl}, \lean{gannonRayhaun_hkl_qexp},
\lean{gannonRayhaun_hkl_final}, \lean{gannonRayhaun_hkl_dim}, \ldots, back to \lean{gannonRayhaun_of_iTannakian} and
\lean{gannonRayhaun_current_gauss}, which instantiates the conditional theorem of Section~\ref{l22:sec:conditional}.
The theorem is universe-polymorphic.
\end{dictionary}

\begin{proof}[Proof on paper]
Theorem~\ref{l22:thm:conditional} with every hypothesis discharged as in the table.
\end{proof}

\begin{physics}
Any rational, positive-energy chiral algebra at $c=1$ is a free boson at a rational radius, its $\Z_2$ orbifold,
or one of the three exceptional gaugings of $SU(2)_1$.
\end{physics}

\begin{status}
\textbf{Proved with no hypotheses.} \lean{#print axioms gannonRayhaun_forward} (re-probed at an earlier commit, not re-run for this edition; on
2026-10-13, when \lean{VOA.pbw} was an axiom, the list also contained \lean{VOA.pbw}):
\begin{plaincode}
'VOA.Classification.COne.gannonRayhaun_forward' depends on axioms:
  [propext, Classical.choice, Quot.sound]
\end{plaincode}
So it is unconditional; it uses the PBW theorem \lean{VOA.pbw} (Chapter~\ref{ch:L02}). It is also recorded in the
maintained audit \leanpath{scripts/\allowbreak audit/\allowbreak ClassificationCOneAxiomAudit.lean} with expected axioms
\lean{PBW_AXIOMS} (since \lean{VOA.pbw} became a theorem, \lean{PBW_AXIOMS} is the standard three); the audit passed on 2026-10-13
(844 declarations).
\end{status}

%======================================================================
\section{Pairwise distinctness}\label{l22:sec:distinct}
%======================================================================

Most pairs of labels are separated by easy invariants: $\dim V_1$ ($=1$ for $V_{2m}$ with $m\ge2$, $3$ for $V_2$,
$0$ for $V_{2m}^+$ and the exceptional orbifolds), the weight of the first primary, the number of irreducible modules.
The two hard cases, isolated in \leanpath{VOA/\allowbreak Classification/\allowbreak COne/\allowbreak Distinct.lean},
are: $V_{2m}^+$ with $m\ge4$ against the exceptional orbifolds, and the three exceptional orbifolds among themselves.

\begin{definition}[The two residuals]\label{l22:def:residuals}
\textbf{PlusExceptionalResidual}: no VOA is both $V_{2m}^+$ ($m\ge4$) and an exceptional orbifold.
\textbf{ExceptionalResidual}: no VOA carries two different exceptional labels.
\end{definition}
\noindent File \leanpath{VOA/\allowbreak Classification/\allowbreak COne/\allowbreak Distinct.lean}.
\begin{leancode}
def PlusExceptionalResidual : Prop :=
  ∀ (𝒱 : Type u) [AddCommGroup 𝒱] [Module ℂ 𝒱] (V : VertexOperatorAlgebra 𝒱) (m : ℕ)
    (l : COneLabel), 4 ≤ m → l.IsExceptional → HasCOneLabel V (.plus m) → HasCOneLabel V l → False
\end{leancode}
\noindent File \leanpath{VOA/\allowbreak Classification/\allowbreak COne/\allowbreak Distinct.lean}.
\begin{leancode}
def ExceptionalResidual : Prop :=
  ∀ (𝒱 : Type u) [AddCommGroup 𝒱] [Module ℂ 𝒱] (V : VertexOperatorAlgebra 𝒱) (l l' : COneLabel),
    l.IsExceptional → l'.IsExceptional → HasCOneLabel V l → HasCOneLabel V l' → l = l'
\end{leancode}

\begin{dictionary}
\lean{l.IsExceptional} says $l\in\{\mathrm{tetra},\mathrm{octa},\mathrm{icosa}\}$. Concluding \lean{False} from two labels
is ``these cannot both hold''.
\end{dictionary}

\begin{physics}
The question is whether two differently constructed gaugings of $SU(2)_1$ could secretly be the same theory. Their
spectra of states differ, which decides it.
\end{physics}

\begin{status}
\lean{Prop}s, both proved below. The docstrings of both still say ``not proved here''; they are proved in
\leanpath{DistinctResiduals.lean} (Section~\ref{l22:sec:defects}).
\end{status}

\begin{theorem}[Molien's formula for orbifold dimensions]\label{l22:thm:molien}
For a VOA $W$ and a finite $G\le\Aut W$, $|G|\dim(W^G)_n=\sum_{g\in G}\operatorname{tr}(g|W_n)$.
\end{theorem}
\noindent File \leanpath{VOA/\allowbreak Classification/\allowbreak COne/\allowbreak DistinctResiduals.lean}.
\begin{leancode}
theorem card_mul_finrank_orbifold (G : Subgroup (Aut W)) [Fintype G] (n : ℤ)
    [FiniteDimensional ℂ (W.grading n)] :
    (Fintype.card G : ℂ) * finrank ℂ ((orbifoldVOA G).grading n)
      = ∑ g : G, LinearMap.trace ℂ _ (autRes W n g) :=
\end{leancode}

\begin{dictionary}
\lean{autRes W n g} is the restriction of $g$ to $W_n$; \lean{LinearMap.trace ℂ _} its trace. The left side is a
complex number; the formula is the averaging (Reynolds) projection.
\end{dictionary}

\begin{proof}[Proof on paper]
$P=|G|^{-1}\sum_g g$ is a projection of $W_n$ onto $(W_n)^G=(W^G)_n$; the trace of a projection is the dimension of its
image.
\end{proof}

\begin{physics}
The number of gauge-invariant states is the average over the group of the twisted traces --- the untwisted-sector
part of the orbifold partition function.
\end{physics}

\begin{status}
Proved, Lean's three axioms.
\end{status}

\begin{theorem}[Torus traces on $V_2$]\label{l22:thm:toruscharacter}
For the torus element $g_u$ ($u\in\C^\times$) acting on the charge-$k$ sector by $u^k$,
$\operatorname{tr}(g_u|(V_2)_n)=\sum_{w^2\le n}u^w p(n-w^2)$, where $p$ is the partition function.
\end{theorem}
\noindent File \leanpath{VOA/\allowbreak Classification/\allowbreak COne/\allowbreak DistinctResiduals.lean}.
\begin{leancode}
theorem trace_charU (n : ℕ) (u : ℂˣ) :
    haveI := (V1).finrank (n : ℤ)
    LinearMap.trace ℂ _ (autRes V1 n (charU 1 le_rfl u)) = torusTrace n u := by
\end{leancode}

\begin{dictionary}
\lean{V1} is $V_2$ (a local notation in \leanpath{DistinctResiduals.lean} for
\lean{latticeVOAGen (latP 1) (latP_even 1) (latP_posDef 1 le_rfl)}, the unfolding of \lean{V2m 1 le_rfl}); \lean{haveI := (V1).finrank (n : ℤ)}
only registers that $(V_2)_n$ is finite-dimensional, so that the trace is the usual one (Mathlib's
\lean{LinearMap.trace} is $0$ on infinite-dimensional spaces). \lean{charU 1 le_rfl u} is the automorphism of $V_2$
acting by $u^k$ on $e^{k\alpha}\otimes M(1)$; \lean{autRes V1 n g} its restriction to $(V_2)_n$;
\lean{torusTrace n u} is the right-hand side $\sum_{w^2\le n}u^wp(n-w^2)$. Every finite-order automorphism of $V_2$ is conjugate to some $g_u$
(\lean{exists_conj_charU}), so these traces suffice.
\end{dictionary}

\begin{proof}[Proof on paper]
$(V_2)_n=\bigoplus_w e^{w\alpha}\otimes M(1)_{n-w^2}$ and $\dim M(1)_k=p(k)$; $g_u$ is the scalar $u^w$ on the $w$-th
summand.
\end{proof}

\begin{physics}
A rotation about a fixed axis of $SU(2)$ acts on the states of momentum $w$ by the phase $u^w$; the trace is the
partition function with a $U(1)$ chemical potential.
\end{physics}

\begin{status}
Proved, Lean's three axioms.
\end{status}

\begin{theorem}[Pairwise distinctness]\label{l22:thm:distinct}
A VOA carries at most one label of the list.
\end{theorem}
\noindent File \leanpath{VOA/\allowbreak Classification/\allowbreak COne/\allowbreak DistinctResiduals.lean}.
\begin{leancode}
theorem plusExceptionalResidual : PlusExceptionalResidual.{u} := by
\end{leancode}
\noindent File \leanpath{VOA/\allowbreak Classification/\allowbreak COne/\allowbreak DistinctResiduals.lean}.
\begin{leancode}
theorem exceptionalResidual : ExceptionalResidual.{u} := by
\end{leancode}
\noindent File \leanpath{VOA/\allowbreak Classification/\allowbreak COne/\allowbreak DistinctResiduals.lean}.
\begin{leancode}
theorem gannonRayhaunDistinct : GannonRayhaunDistinct.{u} :=
  gannonRayhaunDistinct_of_residuals plusExceptionalResidual exceptionalResidual
\end{leancode}

\begin{dictionary}
Universe-polymorphic and hypothesis-free. Together with Theorem~\ref{l22:thm:forward} it gives GR's
``isomorphic to \emph{exactly} one''.
\end{dictionary}

\begin{proof}[Proof on paper]
With Molien and the torus traces, and the element counts by order of $A_4$, $S_4$, $A_5$ ($(1,3,8,0,0)$,
$(1,9,8,6,0)$, $(1,15,20,0,24)$ for orders $1,\dots,5$), one computes exactly
\begin{equation*}
\dim(V_2^G)_4=p(4)-p(3)=2,\qquad
\dim(V_2^G)_{16}=p(16)-p(15)+\begin{cases}p(7)&G\cong A_4\\ p(0)&G\cong S_4\\ 0&G\cong A_5\end{cases}
\end{equation*}
($70$, $56$, $55$). Since $p(7)=15\ge2>p(0)=1>0$, the three are pairwise different (\lean{exceptionalResidual}). For
$V_{2m}^+$ with $m\ge4$ one finds $\dim(V_{2m}^+)_4=e+[m=4]$ with $e$ the number of partitions of $4$ into an even
number of parts, while $p(4)-p(3)=e+o-3$ with $o\le2$ the odd count; equality is impossible
(\lean{plusExceptionalResidual}). The easy pairs are handled in \leanpath{Distinct.lean}. The weight-$16$ numbers were
also checked independently by a spin decomposition $\sum_j(p(n-j^2)-p(n-(j+1)^2))\dim(\mathrm{spin}\ j)^G$, which
is an independent consistency check, not part of the proof.
\end{proof}

\begin{physics}
The tetrahedral, octahedral and icosahedral gaugings keep different numbers of $SU(2)$-invariant combinations of
spin-$j$ states (the first difference is at spin $3$ and $4$); counting gauge-invariant states at energy $16$
distinguishes them.
\end{physics}

\begin{status}
\textbf{Proved with no hypotheses.} \lean{#print axioms gannonRayhaunDistinct} (machine-checked 2026-10-13):
\lean{[propext, Classical.choice, Quot.sound]} --- no \lean{VOA.pbw}.
\end{status}

%======================================================================
\section{The converse}\label{l22:sec:converse}
%======================================================================

The converse says every VOA in the list is nice and has $c=1$. For $V_{2m}$ and $V_{2m}^+$ this is proved directly; for
the three exceptional orbifolds the proof needs four facts (Definition~\ref{l22:def:Bfacts}): the
$C_2$-cofiniteness of $V_2^{A_4}$ and $V_2^{S_4}$, Xu's icosahedral relations (which give the $C_2$-cofiniteness of
$V_2^{A_5}$), and the positivity of the conformal weights of the exceptional orbifolds are explicit computations
(this section); the rationality of the exceptional orbifolds --- more precisely, of every $C_2$-cofinite orbifold of
$V_2$ by a finite group --- is a theorem of McRae whose formalization occupies Section~\ref{l22:sec:B3}. All four are proved, and the assembled converse is
Theorem~\ref{l22:thm:final}.

\subsection{The lattice and plus families}

\begin{theorem}[$V_{2m}$ is nice]\label{l22:thm:V2mnice}
For every $m\ge1$, $V_{2m}$ is strongly rational and pseudo-unitary.
\end{theorem}
\noindent File \leanpath{VOA/\allowbreak Classification/\allowbreak COne/\allowbreak LatticeStronglyRational.lean}.
\begin{leancode}
theorem V2m_stronglyRational (m : ℕ) (hm : 1 ≤ m) : StronglyRational (V2m m hm) :=
\end{leancode}
\noindent File \leanpath{VOA/\allowbreak Classification/\allowbreak COne/\allowbreak LatticeNice.lean}.
\begin{leancode}
theorem V2m_pseudoUnitary : PseudoUnitary (V2m m hm) :=
\end{leancode}
\noindent File \leanpath{VOA/\allowbreak Classification/\allowbreak COne/\allowbreak LatticeNice.lean}.
\begin{leancode}
theorem V2m_nice : Nice (V2m m hm) := ⟨V2m_stronglyRational m hm, V2m_pseudoUnitary m hm⟩
\end{leancode}

\begin{dictionary}
\lean{V2m_pseudoUnitary} and \lean{V2m_nice} are stated inside a section with variables \lean{(m : ℕ) (hm : 1 ≤ m)},
so they hold for every $m\ge1$.
\end{dictionary}

\begin{proof}[Proof on paper]
Simplicity and self-duality come from the nondegenerate invariant form of a positive-definite lattice VOA;
$C_2$-cofiniteness from the Fock-space description (the $C_2$ quotient is spanned by finitely many classes of
$e^{\beta}$ and $\alpha(-1)^k$); rationality (Dong; Dong--Li--Mason) from a Stone--von Neumann decomposition of every
admissible module into $M(1)\otimes\Omega$, semisimplicity of $h(0)$ on $\Omega$ and the classification of irreducibles
$V_{L+j/(2m)\alpha}$, $j\bmod 2m$. Pseudo-unitarity: the irreducible modules not isomorphic to $V_{2m}$ are
$V_{L+\lambda_j}$ with $0<|j|\le m$, of lowest weight $j^2/(4m)>0$.
\end{proof}

\begin{physics}
The compact boson at a rational radius has $2m$ primaries of weights $j^2/4m$, all positive except the vacuum.
\end{physics}

\begin{status}
Proved unconditionally, Lean's three axioms.
\end{status}

\begin{theorem}[$V_{2m}^+$ is nice]\label{l22:thm:plusnice}
For every $m\ge1$, $V_{2m}^+$ is nice.
\end{theorem}
\noindent File \leanpath{VOA/\allowbreak Classification/\allowbreak COne/\allowbreak MainConverse.lean}.
\begin{leancode}
theorem V2mPlus_nice (m : ℕ) (hm : 1 ≤ m) : Nice (V2mPlus m hm) :=
  nice_V2mPlus_of_rational_pseudoUnitary v2mPlusRationalFactGe2_holds
    v2mPlusPseudoUnitaryFactGe2_holds m hm
\end{leancode}

\begin{proof}[Proof on paper]
$C_2$-cofiniteness by an explicit strong-generation argument; pseudo-unitarity from the known list of irreducible
$V_{2m}^+$-modules (Dong--Nagatomo) and their weights; rationality ($m\ge2$) by a uniform version of Abe's
splitting argument. For $m=1$, $V_2^+\cong V_8$ reduces to Theorem~\ref{l22:thm:V2mnice}.
\end{proof}

\begin{physics}
The points of the orbifold line are unitary rational CFTs.
\end{physics}

\begin{status}
Proved unconditionally, Lean's three axioms (machine-checked 2026-10-13).
\end{status}

\subsection{The exceptional orbifolds: the four inputs}

\begin{definition}[The four exceptional inputs]\label{l22:def:Bfacts}
\textbf{(a) $C_2$-cofiniteness of $V_2^{A_4}$ and $V_2^{S_4}$}: $V_2^G$ is $C_2$-cofinite for $G\cong A_4,S_4$
(Miyamoto's solvable-orbifold theorem, specialized).
\textbf{(b) Xu's icosahedral relations}: for $G\cong A_5$ there are a weight-$36$ vector $X\in V_2^G$ and nonzero
couplings $c_0,c_6$ such that, in the $C_2$-algebra, the classes $D^n\mathbf 1$, $D^nX$ ($D$ = the action of $L_{-2}$)
span, and two explicit linear relations hold among $D^{37}\mathbf 1, D^{19}X$ and $D^{38}\mathbf 1, D^{20}X$ with Xu's
integer coefficients. \textbf{(c) Positivity of the conformal weights of the exceptional orbifolds}: if an exceptional
orbifold is strongly rational, it is pseudo-unitary.
\textbf{(d) Rationality of the exceptional orbifolds}: for every finite $G\le\Aut V_2$, if $V_2^G$ is $C_2$-cofinite then it is rational (McRae,
arXiv:1910.13226, Cor.~4.23).

The project's internal roadmap called these four facts \texttt{B2}, \texttt{B4}, \texttt{B5} and \texttt{B3} respectively; those labels survive in
some Lean names (for instance \lean{gannonRayhaunConverse_of_B3}, \leanpath{MainConverseB4.lean}), but are not used in the prose of
this book.
\end{definition}
\noindent File \leanpath{VOA/\allowbreak Classification/\allowbreak COne/\allowbreak ExceptionalOrbifoldNice.lean}.
\begin{leancode}
def SolvableExceptionalOrbifoldC2Fact : Prop :=
  ∀ G : Subgroup (Aut (V2m 1 le_rfl)),
    (Nonempty (G ≃* alternatingGroup (Fin 4)) ∨ Nonempty (G ≃* Equiv.Perm (Fin 4))) →
    C2Cofinite (orbifoldVOA G)
\end{leancode}
\noindent File \leanpath{VOA/\allowbreak Classification/\allowbreak COne/\allowbreak ExceptionalOrbifoldC2.lean}.
\begin{leancode}
def XuIcosahedralPoissonRelations : Prop :=
  ∀ (G : Subgroup (Aut (V2m 1 le_rfl))), Nonempty (G ≃* alternatingGroup (Fin 5)) →
    ∃ (X : (orbifold G).carrier) (c₀ c₆ : ℂ),
      X ∈ (orbifoldVOA G).grading 36 ∧ c₀ ≠ 0 ∧ c₆ ≠ 0 ∧
      let W := orbifoldVOA G
      let D := exceptionalConformalEnd W
      let v := mk₂ W W.vacuum
      let x := mk₂ W X
      Submodule.span ℂ (Set.range (Sum.elim (fun n : ℕ => (D ^ n) v)
        (fun n : ℕ => (D ^ n) x))) = ⊤ ∧
      (c₀ * (xuU37 : ℂ)) • ((D ^ 37) v) + (c₆ * (xuV19 : ℂ)) • ((D ^ 19) x) = 0 ∧
      (c₀ * (xuU38 : ℂ)) • ((D ^ 38) v) + (c₆ * (xuV20 : ℂ)) • ((D ^ 20) x) = 0
\end{leancode}
\noindent File \leanpath{VOA/\allowbreak Classification/\allowbreak COne/\allowbreak ExceptionalOrbifoldNice.lean}.
\begin{leancode}
def ExceptionalOrbifoldPositiveWeightsFact : Prop :=
  ∀ G : Subgroup (Aut (V2m 1 le_rfl)), IsExceptionalV2Group G →
    StronglyRational (orbifoldVOA G) → PseudoUnitary (orbifoldVOA G)
\end{leancode}
\noindent File \leanpath{VOA/\allowbreak Classification/\allowbreak COne/\allowbreak ExceptionalOrbifoldNice.lean}.
\begin{leancode}
def FiniteV2OrbifoldRationalOfC2Fact : Prop :=
  ∀ (G : Subgroup (Aut (V2m 1 le_rfl))) [Finite G],
    C2Cofinite (orbifoldVOA G) → AdmissiblyRational (orbifoldVOA G)
\end{leancode}

\begin{dictionary}
\lean{(orbifold G).carrier} is the subspace of fixed vectors; \lean{mk₂ W} the projection to the $C_2$-algebra
$R_W=W/C_2(W)$; \lean{exceptionalConformalEnd W} the endomorphism of $R_W$ induced by $L_{-2}$;
\lean{xuU37}, \lean{xuV19}, \lean{xuU38}, \lean{xuV20} are Xu's integers. The \lean{Submodule.span … = ⊤} clause says
$R_W$ is spanned by the two $L_{-2}$-strings, which with the two relations makes $R_W$ finite-dimensional.
\lean{IsExceptionalV2Group G} is the disjunction $G\cong A_4\vee G\cong S_4\vee G\cong A_5$. In
Theorem~\ref{l22:thm:B245} below, \lean{A₂ W} is the same $C_2$-algebra $R_W=W/C_2(W)$, so
\lean{Module.finrank ℂ (A₂ (orbifoldVOA G)) ≤ 58} bounds its dimension.
\lean{[Finite G]} in \lean{FiniteV2OrbifoldRationalOfC2Fact} quantifies over \emph{all} finite subgroups --- more than
the converse needs (only the three exceptional ones), but this is McRae's theorem at its natural generality, and it
is true. \textbf{Stronger than needed, not stronger than the literature.}
\end{dictionary}

\begin{physics}
Facts (a) and (b) say the exceptional gauged theories have finitely many ``classical'' states modulo derivatives
(lisse); fact (d) says lisse orbifolds of rational theories are rational (gauging a finite symmetry of a rational
theory gives a rational theory); fact (c) says gauging keeps energies positive (twisted sectors have positive weight).
\end{physics}

\begin{status}
(a) $C_2$-cofiniteness of $V_2^{A_4},V_2^{S_4}$: \textbf{proved} (\lean{solvableExceptionalOrbifoldC2Fact_holds}).
(b) Xu's relations: \textbf{proved} (\lean{xuIcosahedralPoissonRelations_holds}). (c) Positivity of the conformal
weights: \textbf{proved} (\lean{OrbIcoZ.exceptionalOrbifoldPositiveWeightsFact_holds}). (d) Rationality:
\textbf{proved}
(\lean{finiteV2OrbifoldRationalOfC2}, Theorem~\ref{l22:thm:B3}; it was the last open input of the converse until
2026-10-14).
\end{status}

\begin{theorem}[$C_2$-cofiniteness and positivity for the exceptional orbifolds]\label{l22:thm:B245}
The three facts (a), (b), (c) of Definition~\ref{l22:def:Bfacts} hold. In particular every $V_2^G$ with $G\cong A_5$ is $C_2$-cofinite with
$\dim R_{V_2^G}\le58$.
\end{theorem}
\noindent File \leanpath{VOA/\allowbreak Classification/\allowbreak COne/\allowbreak OrbOctaRCFinal.lean}.
\begin{leancode}
theorem solvableExceptionalOrbifoldC2Fact_holds : SolvableExceptionalOrbifoldC2Fact :=
\end{leancode}
\noindent File \leanpath{VOA/\allowbreak Classification/\allowbreak COne/\allowbreak MainConverseB4.lean}.
\begin{leancode}
theorem xuIcosahedralPoissonRelations_holds : XuIcosahedralPoissonRelations :=
\end{leancode}
\noindent File \leanpath{VOA/\allowbreak Classification/\allowbreak COne/\allowbreak MainConverseB4.lean}.
\begin{leancode}
theorem icosahedral_c2Cofinite_holds
    (G : Subgroup (Aut (V2m 1 le_rfl)))
    (hG : Nonempty (G ≃* alternatingGroup (Fin 5))) :
    C2Cofinite (orbifoldVOA G) ∧ Module.finrank ℂ (A₂ (orbifoldVOA G))≤58 :=
  OrbIco.icosa_c2_of_symbols OrbIcoV.icoSymbolRelations_holds G hG
\end{leancode}
\noindent File \leanpath{VOA/\allowbreak Classification/\allowbreak COne/\allowbreak OrbIcoZFinal.lean}.
\begin{leancode}
theorem exceptionalOrbifoldPositiveWeightsFact_holds :
    ExceptionalOrbifoldPositiveWeightsFact :=
\end{leancode}

\begin{proof}[Proof on paper]
\emph{(a) $C_2$-cofiniteness of $V_2^{A_4}$ and $V_2^{S_4}$.} For $A_4$ and $S_4$ the library proves the needed nilpotence of $L_{-2}$ on the $C_2$-algebra directly:
a spanning set of $V_2^G$ by Virasoro descendants of finitely many primaries (from invariant theory of the binary
polyhedral group acting on the charge sectors), and explicit relations (Dong--Jiang type) giving
$D^{25}=0$ in the octahedral case. A recorded trap: a constant printed in the source paper is wrong; the relation used
is the corrected one (Section~\ref{l22:sec:defects}).
\emph{(b) Xu's icosahedral relations.} Xu's relations are established in the actual VOA: a generic Ward-identity solution over a family of Fock
modules is specialized to $c=1$, resonant pivots are cancelled by exact jet arithmetic, and the four actual
coefficients are identified with nonzero ground couplings times Xu's values.
\emph{(c) Positivity of the conformal weights.} Positivity of the conformal weights of all irreducible modules of the exceptional orbifolds is reduced to
finite Zhu-algebra certificates (level-$13$ mode computations, exact determinant factorizations with strictly
positive roots).
\end{proof}

\begin{physics}
These are the hard, explicit computations behind the statement ``the exceptional $c=1$ theories are lisse and
unitary-like''. The icosahedral case is the hardest because $A_5$ is not solvable, so Miyamoto's theorem does not
apply.
\end{physics}

\begin{status}
Proved unconditionally. \lean{#print axioms}: Lean's three (maintained audit, 2026-10-13).
\end{status}

\begin{theorem}[The converse from orbifold rationality; historical reduction]\label{l22:thm:converseB3}
If the rationality of the exceptional orbifolds \textup{(}fact~(d) of Definition~\ref{l22:def:Bfacts}\textup{)} holds, every VOA in the list (with carrier in universe $0$) is nice and has central charge $1$.
\end{theorem}
\noindent File \leanpath{VOA/\allowbreak Classification/\allowbreak COne/\allowbreak MainConverseB4.lean}.
\begin{leancode}
theorem gannonRayhaunConverse_of_B3 (hRat : FiniteV2OrbifoldRationalOfC2Fact) :
    GannonRayhaunConverse.{0} :=
  gannonRayhaunConverse_of_B3_B4 xuIcosahedralPoissonRelations_holds hRat
\end{leancode}

\begin{dictionary}
One explicit hypothesis, \lean{hRat : FiniteV2OrbifoldRationalOfC2Fact}. Conclusion
\lean{GannonRayhaunConverse.{0}}: the converse for carriers in \lean{Type 0}. This was, until 2026-10-14, the form in
which the library stated the converse. The hypothesis is now a theorem (\lean{finiteV2OrbifoldRationalOfC2},
Theorem~\ref{l22:thm:B3}), so the reduction is kept as a step of the proof, not as a residual.
\end{dictionary}

\begin{remark}[Universe]\label{l22:rem:universe}
The converse is stated at universe \lean{0} only, while the forward theorem and distinctness hold in every universe.
Every VOA of the list has a model on a \lean{Type 0} carrier ($V_{2m}$ lives there), so nothing is lost
mathematically; but a listed VOA built on a carrier in a higher universe is not covered by the Lean statement, and
the library does not prove a universe-transfer lemma for \lean{Nice} (whose \lean{PseudoUnitary} clause quantifies
over modules in the universe of $V$). A referee should read the converse as ``every listed VOA built on a
\lean{Type 0} carrier is nice''. The restriction comes from the proof, not from the mathematics: a VOA in the list is compared, through an
isomorphism, with one of the concrete models $V_{2m}$, $V_{2m}^+$, $V_2^G$, all built on \lean{Type 0} carriers, and
niceness is transported along that isomorphism, which the library does only between carriers in the same universe.
The final theorem (Section~\ref{l22:sec:final}) therefore states all three parts at universe \lean{0}.
\end{remark}

\begin{proof}[Proof on paper]
Write $V$ with a label. Lattice: Theorem~\ref{l22:thm:V2mnice} and transport of niceness along isomorphisms. Plus:
Theorem~\ref{l22:thm:plusnice}. Exceptional: simplicity, self-duality, $c=1$ and $V_1=0$ come from the orbifold
structure; $C_2$-cofiniteness from fact (a) or from Xu's relations (b); rationality from $C_2$-cofiniteness by fact (d);
pseudo-unitarity by the positivity of the conformal weights (c).
\end{proof}

\begin{physics}
Once one knows that gauging a finite symmetry of $SU(2)_1$ yields a rational theory (Section~\ref{l22:sec:B3}), the
whole list consists of rational unitary-like chiral algebras: the only thing the reduction still needs from outside is
the absence of logarithmic sectors in the three exceptional gaugings.
\end{physics}

\begin{status}
Proved, \lean{#print axioms} = \lean{[propext, Classical.choice, Quot.sound]} (machine-checked 2026-10-13).
Conditional on \lean{FiniteV2OrbifoldRationalOfC2Fact}, \textbf{which is now proved} (Theorem~\ref{l22:thm:B3}).
Composing the two gives the unconditional converse \lean{gannonRayhaunConverse_holds}
(Theorem~\ref{l22:thm:final}); note that the composition uses the theorem \lean{VOA.pbw} through the proof of orbifold
rationality, although this reduction step alone does not.
\end{status}

%======================================================================
\section{Rationality of the orbifolds of \texorpdfstring{$V_2$}{V2}: McRae's orbifold theorem}\label{l22:sec:B3}
%======================================================================

This section proves the rationality of the exceptional orbifolds, fact~(d) of Definition~\ref{l22:def:Bfacts}, in
its general form: \emph{for every finite group $G$ of automorphisms of $V_2$
such that $V_2^G$ is $C_2$-cofinite, $V_2^G$ is rational}. It was the last open input of the converse. The route is
McRae's theorem on compact orbifolds (arXiv:1910.13226, Theorem~4.20 and Corollary~4.23), specialized to $V=V_2$:
\begin{quote}
\emph{If $V$ is strongly rational, $G$ is a finite group of automorphisms of $V$, and $V^G$ is $C_2$-cofinite, then
$V^G$ is strongly rational.}
\end{quote}
The library formalizes the argument in seven steps, listed with their titles in Section~\ref{l22:sec:B3plan} and
referred to below as Step~1, \ldots, Step~7. Several of them are proved for an arbitrary ambient VOA,
some only for $V_2$; the dependence on $V_2$ enters at exactly one point, Step~5 (twisted parts are modules for the
cyclic orbifolds), through the fact that every cyclic
orbifold $V_2^{\langle g\rangle}$ is a lattice VOA and hence rational. We first explain the shape of the argument, then
go through the steps.

\subsection{The plan, and the notation}\label{l22:sec:B3plan}

\paragraph{Notation.} Throughout, $V$ is the ambient VOA (for us $V=V_2$), $G\le\Aut V$ is a finite subgroup, and
$V^G$ is the fixed-point subalgebra (\lean{orbifoldVOA G}, Chapter~\ref{ch:04}). The reader should note a clash with
the Lean names: in the Lean files of this section the \emph{ambient VOA} is called \lean{W} and its \emph{carrier type}
is called \lean{V}, so \lean{orbifoldVOA G} is our $V^G$. We write $\mathcal C$ for the category of finitely generated
grading-restricted generalized (``logarithmic'') $V^G$-modules, \lean{LogModuleCat (orbifoldVOA G)}
(Chapter~\ref{ch:L21}, Definition~\ref{l21:def:logcat}); $\boxtimes$ for the Huang--Lepowsky--Zhang tensor product in
$\mathcal C$; $\mathbf 1=V^G$ for its unit; $\alpha,\lambda,\rho$ for associator and unitors; $c_{X,Y}:X\boxtimes Y\to
Y\boxtimes X$ for the braiding; and $M_{X,Y}=c_{Y,X}\circ c_{X,Y}$ for the \emph{monodromy} (double braiding).

\paragraph{The idea.} Rationality of $V^G$ means that every (admissible) $V^G$-module is a direct sum of
irreducibles. The danger, for a $C_2$-cofinite VOA, is \emph{logarithmic} modules: indecomposable modules on which
$L_0$ has Jordan blocks, which are extensions of irreducibles that do not split. McRae's insight is that every
$V^G$-module $Y$ can be ``lifted'' to a $V$-module in a generalized sense, namely the \emph{induced module}
$V\boxtimes Y$, and that such generalized $V$-modules are controlled by honest \emph{twisted} $V$-modules, which, for a
rational $V$, are well behaved. In more detail, the seven steps, with the titles by which we refer to them, are:
\begin{enumerate}
\item[Step 1.] \emph{Finitely many irreducibles.} Every $C_2$-cofinite VOA has finitely many irreducible modules,
  without assuming that its Zhu algebra is semisimple; this finiteness theorem is needed to construct $\boxtimes$ on
  finite-length modules.
\item[Step 2.] \emph{The braided tensor category of logarithmic modules.} Build the braided monoidal category
  $\mathcal C$ of logarithmic $V^G$-modules: the pentagon, the braiding and the hexagons.
\item[Step 3.] \emph{The ambient VOA as a commutative algebra.} Show that $V$ itself, viewed as a $V^G$-module, is an object of $\mathcal C$ and a \emph{commutative
  algebra} in it: the vertex operator of $V$ is a $V^G$-intertwining operator of type $\binom{V}{V\,V}$, hence a morphism
  $\mu:V\boxtimes V\to V$, which is associative and commutative. Left modules for this algebra are ``generalized
  $V$-modules'' (they may be twisted).
\item[Step 4.] \emph{Projections onto twisted parts.} Show that every module $X$ for the algebra $V$ splits, by explicit idempotents $\pi_g$ ($g\in G$), into
  parts $\pi_gX$ on which the monodromy of $V$ around $X$ acts by $g$: the $g$-\emph{twisted parts} (McRae's
  Theorem~3.3). The idempotents are built from a \emph{coevaluation} $\tilde\imath:\mathbf 1\to V\boxtimes V$, which
  exists because $V$ is a rigid self-dual object of dimension $|G|$, proved here by Galois descent. We split this
  step in three: \emph{Step 4a}, the Dong--Li--Mason Schur--Weyl duality, which supplies the finiteness and density
  inputs; \emph{Step 4b}, rigidity and self-duality of $V$ by Galois descent; \emph{Step 4c}, McRae's projections
  $\pi_g$ onto the twisted parts.
\item[Step 5.] \emph{Twisted parts are modules for the cyclic orbifolds.} On a $g$-twisted part, the elements of $V$
  fixed by $g$ have trivial monodromy, so their vertex operators
  have integral powers and no logarithms (\emph{Step 5a}); they satisfy the Jacobi identity (\emph{Step 5b}, the
  ``Jacobi bridge''), so $\pi_gX$ is an honest module for
  the cyclic orbifold $V^{\langle g\rangle}$. For $V=V_2$, $V_2^{\langle g\rangle}$ is a lattice VOA $V_{2n^2}$, hence
  rational (\emph{Step 5c}). So every surjection of $V$-modules splits $V^G$-linearly on each twisted part, and hence
  on the whole module (\emph{Step 5d}).
\item[Step 6.] \emph{Every surjection splits (McRae's Lemma~4.21).} From Step~5 applied to induced modules
  $V\boxtimes Y$, every surjection in $\mathcal C$ splits.
\item[Step 7.] \emph{Split surjections imply rationality.} A $C_2$-cofinite VOA all of whose surjections in
  $\mathcal C$ split is rational.
\end{enumerate}
The Lean files and some Lean names carry the project's internal labels for these steps, \texttt{O1}, \ldots,
\texttt{O7} (with sub-steps such as \texttt{O4b}, \texttt{O5b}); for instance \leanpath{McRaeO1FinitePhi.lean} belongs
to Step~1 and the namespace \lean{McRaeO5b} to Step~5b. The label \texttt{O}$n$ always means Step~$n$, and the
label \texttt{M2} in \leanpath{McRaeM2LogBraided.lean} is left over from the internal labels of the earlier,
abandoned route (Section~\ref{l22:sec:oldroute}); that file belongs to Step~2. We use the step numbers only.

The argument never needs the rigidity of $\mathcal C$ itself, modular invariance of logarithmic characters,
factorizability, or the semisimplicity of Zhu's algebra of $V^G$ --- the inputs of the earlier, abandoned route
(Section~\ref{l22:sec:oldroute}). That is why it could be completed.

\paragraph{Where \texttt{VOA.pbw} enters.} The construction of $\boxtimes$ uses the finiteness theorem of Step~1 in its
\lean{_pbw} form, so every statement below that mentions the tensor category uses the theorem \lean{VOA.pbw}. The purely
algebraic or analytic steps (Schur--Weyl duality, the abstract projection theorem, the residue identity, the
rationality of cyclic orbifolds, splitting $\Rightarrow$ rationality) do not. The status boxes record which is which,
from the project's audit logs. For this revision the \lean{#print axioms} of every headline
quoted below was re-run against the built tree at an earlier commit (with \texttt{lake env lean
-DmaxSynthPendingDepth=3}, the audit-safe invocation; not repeated for this edition), and agrees with those logs; the final theorems were in addition
machine-checked by the coordinator on 2026-10-14 (Section~\ref{l22:sec:final}).

\paragraph{The physical picture of the whole argument.} Gauging a finite symmetry $G$ of a rational CFT produces a
theory whose sectors are labelled by a twist $g\in G$ (the boundary condition of the fields around the insertion
point) and by a representation of the centralizer. The untwisted sector decomposes $V$ under $G$; the $g$-twisted
sectors are the $g$-twisted modules of $V$. McRae's argument is the precise mathematical version of the statement
``every sector of the gauged theory comes from a twisted sector of the parent theory''. The parent theory $SU(2)_1$
($V_2$) is rational, and each of its twisted sectors, seen through the cyclic subgroup $\langle g\rangle$, is a sector of a
compact boson at a rational radius; so no logarithmic sector can appear in the gauged theory.

%----------------------------------------------------------------------
\subsection{Step 1: finitely many irreducibles, without a semisimple Zhu algebra}

\begin{theorem}[Step 1: finitely many irreducibles]\label{l22:thm:O1}
Every $C_2$-cofinite VOA has finitely many isomorphism classes of irreducible ordinary modules.
\end{theorem}
\noindent File \leanpath{VOA/\allowbreak Classification/\allowbreak COne/\allowbreak McRaeO1FinitePhi.lean}.
\begin{leancode}
theorem finite_Φ_of_C2Cofinite_pbw (hC2 : C2Cofinite W) : Finite (Φ W) :=
  finite_Φ_of_C2Cofinite' (voaPBWBottomProjection (parabolicRepOfAMod W (regularAMod W))) hC2
\end{leancode}

\begin{dictionary}
\lean{Φ W} is the type of isomorphism classes of irreducible ordinary $W$-modules, \lean{Finite} says it is finite.
The only hypothesis is \lean{C2Cofinite W}. The proof term shows how \lean{VOA.pbw} enters: the bottom projection of
Zhu's induced module is supplied by \lean{voaPBWBottomProjection}. The full fidelity check, the \lean{Pr}-hypothesis primed form
and the consumer \lean{finiteLength_maximal_of_c2} are in Chapter~\ref{ch:L21}, Theorem~\ref{l21:thm:finitephi}.
\end{dictionary}

\begin{proof}[Proof on paper]
Zhu's algebra $A(W)$ is finite-dimensional, hence Artinian; $A/J(A)$ is semisimple with finitely many simple summands;
every simple $A$-module is killed by $J(A)$, so it is one of these; bottom levels of irreducible modules are simple
$A$-modules and determine the module. (Chapter~\ref{ch:L21}.)
\end{proof}

\begin{physics}
A lisse chiral algebra has finitely many primary fields, logarithmic or not; this is what makes the fusion of
finite-length modules land in finite-length modules.
\end{physics}

\begin{status}
Proved; Lean's three axioms (it uses the theorem \lean{VOA.pbw}; while \lean{VOA.pbw} was an axiom the report also listed \lean{VOA.pbw}) (audit log of Step~1). It removed the hypothesis \lean{hss} (semisimple Zhu
algebra) from the whole logarithmic tensor chain.
\end{status}

%----------------------------------------------------------------------
\subsection{Step 2: the braided tensor category of logarithmic modules}

\begin{theorem}[Step 2: braided tensor category; Huang--Lepowsky--Zhang]\label{l22:thm:O2}
For every $C_2$-cofinite VOA $W$, the category of finitely generated logarithmic $W$-modules, with the
Huang--Lepowsky--Zhang tensor product, the vacuum module as unit, the analytically constructed associator, the unitors
and the braiding induced by skew-symmetry, is a braided monoidal category.
\end{theorem}
\noindent File \leanpath{VOA/\allowbreak Classification/\allowbreak COne/\allowbreak McRaeM2LogBraided.lean}.
\begin{leancode}
theorem nonempty_logarithmicBraidedCategory (hC2 : C2Cofinite W) :
    Nonempty (@BraidedCategory (LogModuleCat W) _ (logarithmicMonoidalCategory hC2)) :=
  ⟨logarithmicBraidedCategory hC2⟩
\end{leancode}

\begin{dictionary}
The two structures \lean{logarithmicMonoidalCategory hC2} and \lean{logarithmicBraidedCategory hC2} are quoted in full,
with their fidelity checks, in Chapter~\ref{ch:L21}, Theorems~\ref{l21:thm:logmon} and~\ref{l21:thm:logbraid}. The
point a referee should check is that the bare type \lean{BraidedCategory (LogModuleCat W)} is not the content: any
category admits trivial monoidal structures. The content is that the \emph{specific} structures are the intended
ones, which the library pins down by definitional lemmas proved by \lean{rfl}: the tensor object is
\lean{tensorObject hC2 N P} (the HLZ $P(z)$-tensor product), the associator is \lean{associatorIso}, the unit is the
adjoint module, and the braiding is \lean{swapIso hC2 X Y 0}, induced by the skew-symmetry operator $\Omega_0$. Every
later step uses these structures through the local instances \lean{letI := logarithmicMonoidalCategory hC2} and
\lean{letI := logarithmicBraidedCategory hC2}.
\end{dictionary}

\begin{proof}[Proof on paper]
See Chapter~\ref{ch:L21}, Section~\ref{l21:sec:log}: the pentagon is the equality of the two ways of associating a
product of four intertwining operators, proved by showing that all five bracketings converge on a common domain,
agree there by repeated use of the associativity of two intertwiners, and that a two-variable power-log series which
vanishes on an open set has zero coefficients. The hexagons combine associativity with the skew-symmetry operator
$\Omega_0(\mathcal Y)(w_2,z)w_1=e^{zL_{-1}}\mathcal Y(w_1,e^{\pi i}z)w_2$, which turns an intertwining operator of type
$\binom{W_3}{W_1\,W_2}$ into one of type $\binom{W_3}{W_2\,W_1}$ and induces the braiding, using analytic continuation of products to a half-plane
and the gluing of four local expansions around a common point.
\end{proof}

\begin{physics}
Fusion of sectors of a logarithmic CFT is associative up to $F$-moves and commutative up to $R$-moves; the pentagon and
hexagons are the consistency of four- and five-point conformal blocks. Here they are proved for logarithmic blocks,
which carry powers of $\log z$.
\end{physics}

\begin{status}
Proved; the only hypothesis is \lean{C2Cofinite W}; Lean's three axioms (it uses the theorem \lean{VOA.pbw}; while \lean{VOA.pbw} was an axiom the report also listed \lean{VOA.pbw}) (audit logs of the two halves of Step~2, monoidal and braided).
\end{status}

%----------------------------------------------------------------------
\subsection{Step 3: the ambient VOA as a commutative algebra}

\begin{definition}[The object $V$ of $\mathcal C$]\label{l22:def:vObj}
Let $V$ be a VOA, $G\le\Aut V$, and assume $V$ is finitely generated as a weak $V^G$-module. Then $V$, with the
$V^G$-action by restriction of its own vertex operator, is an object of $\mathcal C$.
\end{definition}
\noindent File \leanpath{VOA/\allowbreak Classification/\allowbreak COne/\allowbreak McRaeO3Object.lean}.
\begin{leancode}
def vObj (hfg : FinGen (orbifold G).adjointRestrict) : LogModuleCat (orbifoldVOA G) :=
  LogModuleCat.of (ambientLog G) hfg
\end{leancode}
\noindent File \leanpath{VOA/\allowbreak Classification/\allowbreak COne/\allowbreak McRaeO3Object.lean}.
\begin{leancode}
abbrev ambientLog : LogVModule (orbifoldVOA G) V := (orbifold G).adjointRestrictOrd.toLog
\end{leancode}

\begin{dictionary}
\lean{(orbifold G).adjointRestrict} is the adjoint module of the ambient VOA restricted to the subalgebra $V^G$;
\lean{FinGen} says it is finitely generated (finitely many vectors whose $V^G$-submodule is everything).
\lean{ambientLog G} views this restriction as a logarithmic module (Chapter~\ref{ch:L12}, Definition~\ref{l12:def:log});
its grading is the ambient $L_0$-grading, so $L_0$ acts semisimply. \lean{LogModuleCat.of} bundles a logarithmic module
with a proof of finite generation into an object of $\mathcal C$. The hypothesis \lean{hfg} is discharged for $V_2$ in
Step~4a (\lean{V2m_orbifold_finGen}, Section~\ref{l22:sec:B3plan}).
\end{dictionary}

\begin{physics}
The parent theory's chiral algebra, decomposed into sectors of the gauged theory: the untwisted sector of the
orbifold, before projection onto invariants.
\end{physics}

\begin{status}
Definition. Lean's three axioms.
\end{status}

\begin{theorem}[Step 3: $V$ is a haploid commutative algebra in $\mathcal C$]\label{l22:thm:O3}
Let $G\le\Aut V$ be finite, with $V^G$ $C_2$-cofinite and $V$ finitely generated over $V^G$, and $V$ nonzero. There
are morphisms $\mu:V\boxtimes V\to V$ (induced by the vertex operator of $V$), $\iota:\mathbf 1\to V$ (the inclusion
$V^G\subset V$), $\varepsilon:V\to\mathbf 1$ (the average $\frac1{|G|}\sum_g g$) and, for $g\in G$, $g:V\to V$, such
that $(V,\mu,\iota)$ is a commutative monoid object of the braided category $\mathcal C$; $\dim\operatorname{Hom}(\mathbf
1,V)=1$ (``haploid''); each $g$ is an algebra automorphism; and $\varepsilon\circ\iota=\mathrm{id}$, while
$\iota\circ\varepsilon$ is the average over $G$:
\begin{equation*}
  \mu\circ c_{V,V}=\mu,\qquad \varepsilon\circ\iota=\mathrm{id}_{\mathbf 1},\qquad \iota\circ\varepsilon=\frac1{|G|}\sum_{g\in G}g .
\end{equation*}
\end{theorem}
\noindent File \leanpath{VOA/\allowbreak Classification/\allowbreak COne/\allowbreak McRaeO3MonObj.lean}.
\begin{leancode}
theorem haploidCommAlgebra [Nontrivial V] [Fintype G] :
    letI := logarithmicMonoidalCategory hC2
    letI := logarithmicBraidedCategory hC2
    letI := monObj G hfg hC2
    IsCommMonObj (vObj G hfg) ∧
      Module.finrank ℂ (𝟙_ (LogModuleCat (orbifoldVOA G)) ⟶ vObj G hfg) = 1 ∧
      (∀ g : G, IsMonHom (autHom G hfg g)) ∧
      MonObj.one ≫ counitHom G hfg = 𝟙 _ ∧
      counitHom G hfg ≫ MonObj.one =
        (Fintype.card G : ℂ)⁻¹ • ∑ g : G, autHom G hfg g :=
\end{leancode}

\begin{dictionary}
The section variables are \lean{(G : Subgroup (Aut W)) (hfg : FinGen (orbifold G).adjointRestrict)
(hC2 : C2Cofinite (orbifoldVOA G))}. \lean{monObj G hfg hC2} is a Mathlib \lean{MonObj} (monoid object: \lean{one},
\lean{mul} with unit and associativity laws) on \lean{vObj G hfg}, whose \lean{one} is \lean{unitHom} ($\iota$) and whose
\lean{mul} is \lean{mulHom} ($\mu$), both by \lean{rfl} (\lean{monObj_one}, \lean{monObj_mul}). \lean{IsCommMonObj} says
\lean{(β_ V V).hom ≫ mul = mul} for the braiding of \lean{logarithmicBraidedCategory}. \lean{𝟙_ C} is the monoidal unit,
\lean{⟶} the morphism type, so the second clause is $\dim_{\C}\operatorname{Hom}_{\mathcal C}(\mathbf 1,V)=1$.
\lean{autHom G hfg g} is $g$ acting on $V$ as a morphism of $\mathcal C$ (it is $V^G$-linear because $V^G$ is fixed by
$g$), and \lean{IsMonHom} says it preserves unit and multiplication. \lean{counitHom} is the Reynolds average
$\varepsilon$; in the last clause \lean{•} is scalar multiplication of morphisms. \lean{[Nontrivial V]} excludes the
zero carrier (needed for haploidy); \lean{[Fintype G]} gives $|G|$. The statement is the commutative-algebra part of
McRae's Assumption~3.1. Could it hold for a wrong $\mu$? No: \lean{mulHom} is defined (\lean{mulHom}, \lean{McRaeO3Algebra.lean})
as the morphism out of $V\boxtimes V$ corresponding, by the universal property of $\boxtimes$, to the vertex operator
of $V$ itself (\lean{vertexIntertwiner G}).
\end{dictionary}

\begin{proof}[Proof on paper]
\emph{The product.} For $u\in V^G$ and $a,b\in V$, the Jacobi identity of $V$ for the triple $(u,a,b)$ says precisely
that $Y_V(\cdot,z)$ is an intertwining operator of type $\binom{V}{V\,V}$ for the algebra $V^G$; it has integral powers
and no logarithms. The universal property of $V\boxtimes V$ turns it into $\mu$. \emph{Commutativity.} The braiding
is induced by $\Omega_0$: $\Omega_0(\mathcal Y)(a,z)b=e^{zL_{-1}}\mathcal Y(b,e^{\pi i}z)a$. For $\mathcal Y=Y_V$ the
skew-symmetry of a VOA, $Y(a,z)b=e^{zL_{-1}}Y(b,-z)a$, and the fact that $e^{\pi i(n+1)}=(-1)^{n+1}$ for integral $n$,
give $\Omega_0(Y_V)=Y_V$, i.e.\ $\mu\circ c_{V,V}=\mu$. \emph{Associativity.} Testing $\mu\circ(\mu\boxtimes1)$ and
$\mu\circ(1\boxtimes\mu)\circ\alpha$ on the canonical elements of the triple tensor product turns them into the
iterate $Y_V(Y_V(a,z_1-z_2)b,z_2)c$ and the product $Y_V(a,z_1)Y_V(b,z_2)c$; these agree on the common domain by the
associativity of $V$ (the Jacobi identity of $V$, summed into a convergent series at a complex point). \emph{Unit laws}
are creation and the vacuum property. \emph{Haploid.} A $V^G$-module map $f:V^G\to V$ is determined by $f(\mathbf 1)$,
which has weight $0$, and $V_0=\C\mathbf 1$ (the library's VOAs are of CFT type); so $f$ is a multiple of $\iota$.
\emph{Counit.} The average $\varepsilon(a)=\frac1{|G|}\sum_g g a$ lands in $V^G$ and is $V^G$-linear; it fixes $V^G$.
\end{proof}

\begin{physics}
In the gauged theory, the parent chiral algebra $V$ becomes a ``condensable algebra'' of the orbifold's sectors: its
fields close under fusion with trivial self-braiding (locality of the parent theory), and the vacuum appears exactly
once. Undoing the gauging is condensing this algebra.
\end{physics}

\begin{status}
Proved, conditional on the displayed hypotheses \lean{hfg}, \lean{hC2} (both discharged for $V_2$: \lean{hfg} by Step~4a,
\lean{hC2} in the application by the $C_2$-cofiniteness facts (a) and (b) of Definition~\ref{l22:def:Bfacts}). Lean's three axioms (it uses the theorem \lean{VOA.pbw}; while \lean{VOA.pbw} was an axiom the report also listed \lean{VOA.pbw}), inherited from the tensor category (audit log of Step~3). Right exactness of $V\boxtimes-$ is proved in the concrete form \lean{tensorMorphism_epi} (if $q$ is
surjective, $1\boxtimes q$ is an epimorphism), which is what Step~6 uses.
\end{status}

%----------------------------------------------------------------------
\subsection{Step 4a: Schur--Weyl duality and finite generation}

\begin{definition}[Dong--Li--Mason Schur--Weyl duality]\label{l22:def:dlm}
For a simple VOA $V$ and finite $G\le\Aut V$, and each irreducible representation $E$ of $G$, let
$M_E=\operatorname{Hom}_G(E,V)$ (the multiplicity space). DLM duality says: $M_E\ne0$; $M_E$ is an irreducible
$V^G$-module; and $M_E$, $M_{E'}$ are inequivalent for $E\not\cong E'$. Hence $V=\bigoplus_E E\otimes M_E$ with the
$M_E$ pairwise distinct irreducible $V^G$-modules.
\end{definition}
\noindent File \leanpath{VOA/\allowbreak Huang/\allowbreak SchurWeylDLM.lean}.
\begin{leancode}
def DLMSchurWeyl : Prop :=
  ∀ (𝒱 : Type u) [AddCommGroup 𝒱] [Module ℂ 𝒱] (W : VertexOperatorAlgebra 𝒱)
    (G : Subgroup (Aut W)), W.IsSimple → Finite G →
    ∀ (E : Type) [AddCommGroup E] [Module ℂ E] [FiniteDimensional ℂ E]
      (σ : Representation ℂ G E), σ.IsIrreducible →
      multSpace G σ ≠ ⊥ ∧
      (∀ P : Submodule ℂ (E →ₗ[ℂ] 𝒱), P ≤ multSpace G σ → ModeStable G P →
        P = ⊥ ∨ P = multSpace G σ) ∧
      (∀ (E' : Type) [AddCommGroup E'] [Module ℂ E'] [FiniteDimensional ℂ E']
        (σ' : Representation ℂ G E'), σ'.IsIrreducible → IsEmpty (σ.Equiv σ') →
        ∀ T : multSpace G σ →ₗ[ℂ] multSpace G σ',
          (∀ u (hu : u ∈ (orbifold G).carrier) (n : ℤ) (φ : multSpace G σ),
            (T ⟨(modeOp W u n).comp φ, modeOp_comp_mem G σ hu n φ.2⟩ : E' →ₗ[ℂ] 𝒱) =
              (modeOp W u n).comp (T φ : E' →ₗ[ℂ] 𝒱)) → T = 0)
\end{leancode}
\noindent File \leanpath{VOA/\allowbreak Huang/\allowbreak SchurWeylDLMProof.lean}.
\begin{leancode}
theorem dlmSchurWeyl_holds : DLMSchurWeyl.{u} := by
\end{leancode}
\noindent File \leanpath{VOA/\allowbreak Huang/\allowbreak SchurWeylDLMProofFinGen.lean}.
\begin{leancode}
theorem V2m_orbifold_finGen
    (G : Subgroup (Aut (VOA.Classification.COne.V2m 1 le_rfl))) [Finite G] :
    VOA.LogTensor.FinGen (orbifold G).adjointRestrict :=
  orbifold_finGen G (VOA.Classification.COne.V2m_isSimple 1 le_rfl)
\end{leancode}

\begin{dictionary}
\lean{multSpace G σ} is the subspace of linear maps $\varphi:E\to V$ with $g(\varphi(x))=\varphi(\sigma(g)x)$, i.e.\
$\operatorname{Hom}_G(E,V)$ (with $V=\bigoplus_nV_n$ the algebraic direct sum, so this is DLM's
$\bigoplus_n\operatorname{Hom}_G(E,V_n)$). $V^G$ acts on it by post-composition with modes, $\varphi\mapsto
u_{(n)}\circ\varphi$ (\lean{modeOp W u n}). Clause one: $M_E\ne0$. Clause two: every subspace stable under all these
modes (\lean{ModeStable}) is $0$ or all, i.e.\ $M_E$ is irreducible. Clause three: every $V^G$-mode-commuting linear map
$M_E\to M_{E'}$ with $E\not\cong E'$ (\lean{IsEmpty (σ.Equiv σ')}) is zero (Schur form of ``inequivalent''). Irreducible
representations are quantified over all finite-dimensional $E$ in \lean{Type}, which covers every isomorphism class.
\lean{FinGen (orbifold G).adjointRestrict}: $V$ is finitely generated as a weak $V^G$-module, the field that
\lean{vObj} requires. The decomposition $V=\bigoplus_EE\otimes M_E$ as a direct sum is \emph{not} formalized (it would
need a finite list of representatives of the irreducibles of $G$); only the three clauses and finite generation are.
\end{dictionary}

\begin{proof}[Proof on paper]
\emph{Density.} $V$ is a simple module over the algebra generated by its modes; by Schur and $V_0=\C\mathbf 1$ its
commutant is $\C$, so by the Jacobson density theorem the modes act densely. \emph{Dedekind's lemma for automorphisms:}
if $\sum_g (y_g)_{(n)}\,g(w)=0$ for all $n$ and all $w$, then all $y_g=0$; this uses an annihilator-stability lemma (the
vanishing persists after inserting further modes, by the Jacobi identity) and simplicity. \emph{Averaging.} If a sum
$\sum_i f_i\circ\chi_i$ of mode-algebra elements composed with $G$-equivariant maps is itself $G$-equivariant, averaging
over $G$ rewrites it with modes of $G$-invariant vectors, i.e.\ of $V^G$. These three facts give non-vanishing,
irreducibility (any nonzero $\varphi\in M_E$ generates $M_E$ under $V^G$) and inequivalence (a nonzero intertwiner
would have a graph which is a $V^G$-stable subspace of $M_{E\oplus E'}$ not of product form). \emph{Finite generation.}
Apply the same density and Dedekind arguments to the regular representation $\C^G$: finitely many orbit maps
$g\mapsto g(b_j)$ already separate, so finitely many vectors $b_j$ generate $V$ over $V^G$.
\end{proof}

\begin{physics}
The untwisted sector of the gauged theory decomposes the parent chiral algebra by $G$-charge: each irreducible charge
$E$ comes with its own primary family $M_E$, every charge occurs, and different charges give different sectors. This is
the ``Schur--Weyl'' or ``charge--sector'' duality familiar from orbifold CFT.
\end{physics}

\begin{status}
Proved: \lean{dlmSchurWeyl_holds} and \lean{V2m_orbifold_finGen} use Lean's three axioms only, no \lean{VOA.pbw} (Step~4a
audit log).
\end{status}

%----------------------------------------------------------------------
\subsection{Step 4b: \texorpdfstring{$V$}{V} is rigid and self-dual of dimension \texorpdfstring{$|G|$}{|G|}: Galois descent}

\begin{theorem}[Step 4b: rigidity and self-duality; McRae Assumption~3.1, rigidity and dimension]\label{l22:thm:O4b}
In the setting of Theorem~\ref{l22:thm:O3}, with $V$ simple, there is a morphism $\tilde\imath:\mathbf 1\to V\boxtimes V$
such that $(V,\ \mathrm{ev}=\varepsilon\circ\mu,\ \tilde\imath)$ satisfies both zigzag identities (so $V$ is a rigid,
self-dual object of $\mathcal C$), and $\mu\circ\tilde\imath=|G|\,\iota$; in particular $\varepsilon\circ\mu\circ\tilde\imath
=|G|\cdot\mathrm{id}_{\mathbf 1}$, i.e.\ $\dim V=|G|$.
\end{theorem}
\noindent File \leanpath{VOA/\allowbreak Classification/\allowbreak COne/\allowbreak McRaeO4bRigid.lean}.
\begin{leancode}
theorem rigidity_dimension :
    letI := logarithmicMonoidalCategory hC2
    ∃ coev : 𝟙_ (LogModuleCat (orbifoldVOA G)) ⟶ vObj G hfg ⊗ vObj G hfg,
      coev ▷ vObj G hfg ≫ (α_ _ _ _).hom ≫ vObj G hfg ◁ (mulHom G hfg hC2 ≫ counitHom G hfg) =
          (λ_ _).hom ≫ (ρ_ _).inv ∧
        vObj G hfg ◁ coev ≫ (α_ _ _ _).inv ≫ (mulHom G hfg hC2 ≫ counitHom G hfg) ▷ vObj G hfg =
          (ρ_ _).hom ≫ (λ_ _).inv ∧
        coev ≫ mulHom G hfg hC2 = (Fintype.card G : ℂ) • unitHom G hfg := by
\end{leancode}

\begin{dictionary}
Section variables: \lean{G}, \lean{hfg}, \lean{hC2} as before, plus \lean{[Fintype G] [DecidableEq G] (hW : W.IsSimple)}.
Mathlib's notation: \lean{f ▷ X} is $f\boxtimes\mathrm{id}_X$, \lean{X ◁ f} is $\mathrm{id}_X\boxtimes f$, \lean{α_},
\lean{λ_}, \lean{ρ_} the associator and unitors, \lean{≫} composition in diagrammatic order (first left, then right).
The first conjunct is the zigzag $(\mathrm{id}_V\boxtimes\mathrm{ev})\circ\alpha\circ(\tilde\imath\boxtimes\mathrm{id}_V)
=\rho^{-1}\circ\lambda$, the second the other zigzag, with $\mathrm{ev}=\varepsilon\circ\mu$ (\lean{mulHom ≫ counitHom});
together they say $(V,\mathrm{ev},\tilde\imath)$ is an exact pairing, Mathlib's \lean{ExactPairing V V}
(\lean{vObjExactPairing}). The third is the dimension condition. This is exactly the rigidity-and-dimension part of
McRae's Assumption~3.1. It does \emph{not} say $\mathcal C$ is rigid --- only the object $V$ is --- and McRae's argument
needs no more. The $V_2$ instance \lean{V2m_rigidity_dimension} (same file) discharges \lean{hfg} and \lean{hW}.
\end{dictionary}

\begin{proof}[Proof on paper: Galois descent]
McRae proves this with Kirillov's functor $\operatorname{Rep}G\to\mathcal C$; the library uses instead the
Chase--Harrison--Rosenberg argument for Galois extensions of commutative rings, which is shorter. For $g\in G$ put
\begin{equation*}
  \psi_g=\mu\circ(\mathrm{id}\boxtimes g):V\boxtimes V\to V,\qquad a\boxtimes b\mapsto a\cdot g(b).
\end{equation*}
\emph{Part (i) (abstract, in any $\C$-linear monoidal category).} Suppose $\tilde\imath$ satisfies the \emph{Galois
condition} $\psi_g\circ\tilde\imath=|G|\,\delta_{g,1}\,\iota$. Write symbolically $\tilde\imath=\sum_ix_i\boxtimes y_i$.
Then for $a\in V$,
\begin{equation*}
  \sum_i x_i\,\varepsilon(y_ia)=\frac1{|G|}\sum_i\sum_{g}x_i\,g(y_i)\,g(a)
  =\frac1{|G|}\sum_g\Big(\sum_ix_ig(y_i)\Big)g(a)=\frac1{|G|}\,|G|\,a=a,
\end{equation*}
using $\iota\varepsilon=\frac1{|G|}\sum_gg$, that each $g$ is an algebra map, and associativity and commutativity.
This is the first zigzag; the second is symmetric; and $\mu\circ\tilde\imath=\sum_ix_iy_i=|G|\iota$ is the case $g=1$.
In the library this computation is carried out with morphisms (\lean{GaloisData.exactPairing}), the step ``$\sum_h
a\,h(bc)\otimes d=\sum_y a\otimes y(b\,\psi_{y^{-1}}(c\otimes d))$'' being the one identity that needs care.
\emph{Part (ii) (surjectivity).} Write $a\boxtimes_nb\in V\boxtimes V$ for the coefficient of $z^{-n-1}$ in the canonical
intertwining operator $\mathcal Y_\boxtimes(a,z)b$ (no logarithms occur here, because the vertex operator of $V$ has none).
There is $w'\in V\boxtimes V$ with $\psi_g(w')=\delta_{g,1}\mathbf 1$: by the density and
Dedekind arguments of Step~4a, every function $\C^G\to V$ is a combination of $a_{(n)}\circ(\text{orbit map of }b)$, and
$\psi_g(a\boxtimes_nb)=a_{(n)}g(b)$. \emph{Part (iii) (injectivity).} The joint map $\Psi=(\psi_g)_g$ is injective on
$V\boxtimes V$ (an identity of part (i) evaluated on canonical elements). \emph{Part (iv).} Define $\tilde\imath(u)=u_{(-1)}
(|G|w')$. That this is a $V^G$-module map is checked after composing with each $\psi_g$ and transported back by
injectivity; then $\psi_g\circ\tilde\imath=|G|\delta_{g,1}\iota$, and part (i) applies.
\end{proof}

\begin{physics}
$\dim V=|G|$ is the statement that the parent theory, seen as a sector of the gauged theory, has quantum dimension
$|G|$: gauging $G$ divides the ``number of states'' by $|G|$, and the parent algebra is self-conjugate. The Galois
condition says that the fusion $V\times V$ contains the vacuum once for each group element, which is the classical
statement $V\otimes_{V^G}V\cong\bigoplus_{g\in G}V$ of Galois theory.
\end{physics}

\begin{status}
Proved, conditional on the displayed hypotheses (for $V_2$ only \lean{hC2} remains). Lean's three axioms (it uses the theorem \lean{VOA.pbw}; while \lean{VOA.pbw} was an axiom the report also listed \lean{VOA.pbw}), inherited from the tensor category; the abstract Galois-descent lemmas use Lean's three only (Step~4b audit
log). McRae's Lemma~3.8 (haploid + rigid + $\dim=|G|$ $\Rightarrow$ Galois condition) is not formalized and not needed:
the Galois condition is proved directly (\lean{coevHom_psi}).
\end{status}

%----------------------------------------------------------------------
\subsection{Step 4c: McRae's Theorem 3.3, projections onto twisted parts}

\begin{definition}[Modules, monodromy, twisted]\label{l22:def:twisted}
In a braided monoidal category, let $(A,m,u)$ be an algebra. A \emph{left $A$-module} is $(X,a)$ with $a:A\otimes X\to X$
satisfying $a\circ(u\otimes1)=\lambda$ and $a\circ(m\otimes1)=a\circ(1\otimes a)\circ\alpha$. The \emph{monodromy} is
$M_{X,Y}=c_{Y,X}\circ c_{X,Y}$. For an automorphism $\varphi$ of $A$, the module is \emph{$\varphi$-twisted} if
$a\circ(\varphi\otimes1)\circ M_{A,X}=a$. Finally, \emph{McRae's Assumption~3.1} on a commutative algebra $A$ with an action
of a finite group $G$ by algebra automorphisms $g:A\to A$ is the list: $A$ is associative, unital and commutative; a
counit $e:A\to\mathbf 1$ with $u\circ e=\frac1{|G|}\sum_gg$; and a coevaluation $\tilde\imath:\mathbf 1\to A\otimes A$
with $m\circ(1\otimes g)\circ\tilde\imath=|G|\,\delta_{g,1}\,u$ for every $g$ (the Galois condition, established in the proof of Theorem~\ref{l22:thm:O4b} as \lean{coevHom_psi}).
In Lean these are the structures \lean{GaloisData} and \lean{McRaeData} quoted below.
\end{definition}
\noindent File \leanpath{VOA/\allowbreak Classification/\allowbreak COne/\allowbreak McRaeO4cAbstract.lean}.
\begin{leancode}
def monodromy (X Y : C) : X ⊗ Y ⟶ X ⊗ Y := (β_ X Y).hom ≫ (β_ Y X).hom
\end{leancode}
\noindent File \leanpath{VOA/\allowbreak Classification/\allowbreak COne/\allowbreak McRaeO4cAbstract.lean}.
\begin{leancode}
structure IsLeftModule (m : A ⊗ A ⟶ A) (u : 𝟙_ C ⟶ A) {X : C} (a : A ⊗ X ⟶ X) : Prop where
  one_smul : u ▷ X ≫ a = (λ_ X).hom
  mul_smul : m ▷ X ≫ a = (α_ A A X).hom ≫ A ◁ a ≫ a
\end{leancode}
\noindent File \leanpath{VOA/\allowbreak Classification/\allowbreak COne/\allowbreak McRaeO4cAbstract.lean}.
\begin{leancode}
def IsTwisted {X : C} (a : A ⊗ X ⟶ X) (φ : A ⟶ A) : Prop :=
  monodromy A X ≫ φ ▷ X ≫ a = a
\end{leancode}
\noindent File \leanpath{VOA/\allowbreak Classification/\allowbreak COne/\allowbreak McRaeO4cAbstract.lean}.
\begin{leancode}
structure McRaeData (A : C) (m : A ⊗ A ⟶ A) (u : 𝟙_ C ⟶ A) (e : A ⟶ 𝟙_ C)
    (G : Type*) [Group G] [Fintype G] [DecidableEq G] (act : G → (A ⟶ A))
    (coev : 𝟙_ C ⟶ A ⊗ A) : Prop extends GaloisData A m u e G act where
  mul_comm : (β_ A A).hom ≫ m = m
  act_comp : ∀ g h, act (g * h) = act h ≫ act g
  coev_psi : ∀ g, coev ≫ A ◁ act g ≫ m = if g = 1 then (Fintype.card G : ℂ) • u else 0
\end{leancode}
\noindent File \leanpath{VOA/\allowbreak Classification/\allowbreak COne/\allowbreak McRaeO4bAbstract.lean}.
\begin{leancode}
structure GaloisData (A : C) (m : A ⊗ A ⟶ A) (u : 𝟙_ C ⟶ A) (e : A ⟶ 𝟙_ C)
    (G : Type*) [Group G] [Fintype G] (act : G → (A ⟶ A)) : Prop where
  one_mul : u ▷ A ≫ m = (λ_ A).hom
  mul_one : A ◁ u ≫ m = (ρ_ A).hom
  mul_assoc : m ▷ A ≫ m = (α_ A A A).hom ≫ A ◁ m ≫ m
  act_mul : ∀ g, (act g ⊗ₘ act g) ≫ m = m ≫ act g
  unit_act : ∀ g, u ≫ act g = u
  act_inv : ∀ g, act g⁻¹ ≫ act g = 𝟙 A
  act_one : act 1 = 𝟙 A
  counit : e ≫ u = (Fintype.card G : ℂ)⁻¹ • ∑ g, act g
\end{leancode}

\begin{dictionary}
\lean{C} is any category with \lean{[MonoidalCategory C] [BraidedCategory C]} (and, for \lean{McRaeData}, $\C$-linear
structure); \lean{β_} is the braiding. The module axioms are Mathlib's \lean{ModObj} axioms
(\lean{isLeftModule_of_modObj}). Note the direction of composition: \lean{monodromy A X ≫ φ ▷ X ≫ a = a} reads
$a\circ(\varphi\otimes1)\circ M_{A,X}=a$, McRae's Definition~2.3. \lean{GaloisData} lists the algebra axioms, that $G$
acts by algebra automorphisms (\lean{act_mul}: $m\circ(g\otimes g)=g\circ m$; \lean{unit_act}; invertibility), and that
\lean{e ≫ u} is the average. \lean{McRaeData} adds commutativity, that \lean{act} is a group action (\lean{act_comp}: \lean{act (g * h) = act h ≫ act g},
i.e.\ $\mathrm{act}(gh)=\mathrm{act}(g)\circ\mathrm{act}(h)$, an ordinary left action written in diagrammatic order), and
the Galois condition \lean{coev_psi}: $m\circ(1\otimes g)\circ\tilde\imath=|G|\delta_{g,1}u$. For $\mathcal C$ and $A=V$
all fields are supplied by Steps~3 and~4b (\lean{mcraeData}, \leanpath{McRaeO4cInstance.lean}). Every field is a true
property of the intended data, so the structure is satisfiable and not vacuous.
\end{dictionary}

\begin{physics}
A module of the condensed algebra $V$ is a sector of the parent theory possibly with a twisted boundary condition: the
monodromy $M_{A,X}$ is the effect of carrying a parent field once around the sector's insertion point, and
``$\varphi$-twisted'' says that this effect is the symmetry $\varphi$ --- exactly a $\varphi$-twisted sector.
\end{physics}

\begin{status}
Definitions. Lean's three axioms.
\end{status}

\begin{theorem}[Step 4c: projections onto twisted parts; McRae Theorem~3.3, projection form]\label{l22:thm:O4c}
Under McRae's Assumption~3.1 (\lean{McRaeData}), every left $A$-module $(X,a)$ carries module endomorphisms
$\pi_g$ ($g\in G$) that are orthogonal idempotents summing to the identity, and the image of $\pi_g$ is
$g$-twisted:
\begin{equation*}
  a\circ(1\otimes\pi_g)=\pi_g\circ a,\quad a\circ(g\otimes1)\circ M_{A,X}\circ(1\otimes\pi_g)=a\circ(1\otimes\pi_g),\quad
  \pi_g\pi_h=\delta_{g,h}\pi_g,\quad \sum_g\pi_g=\mathrm{id}_X .
\end{equation*}
In particular this holds for $A=V_2$ in $\mathcal C$, for every finite $G\le\Aut V_2$ with $V_2^G$ $C_2$-cofinite.
\end{theorem}
\noindent File \leanpath{VOA/\allowbreak Classification/\allowbreak COne/\allowbreak McRaeO4cAbstract.lean}.
\begin{leancode}
theorem theorem_3_3 :
    ∃ π : G → (X ⟶ X),
      (∀ g, a ≫ π g = A ◁ π g ≫ a) ∧
      (∀ g, A ◁ π g ≫ monodromy A X ≫ act g ▷ X ≫ a = A ◁ π g ≫ a) ∧
      (∀ g h, π g ≫ π h = if g = h then π g else 0) ∧
      ∑ g, π g = 𝟙 X :=
\end{leancode}
\noindent File \leanpath{VOA/\allowbreak Classification/\allowbreak COne/\allowbreak McRaeO4cInstance.lean}.
\begin{leancode}
theorem V2m_theorem_3_3 (G : Subgroup (Aut (V2m 1 le_rfl))) [Fintype G] [DecidableEq G]
    (hC2 : C2Cofinite (orbifoldVOA G)) :
    letI := logarithmicMonoidalCategory hC2
    letI := logarithmicBraidedCategory hC2
    ∀ (X : LogModuleCat (orbifoldVOA G)) (a : V2mObj G ⊗ X ⟶ X),
      IsLeftModule (mulHom G (VOA.Huang.SchurWeyl.V2m_orbifold_finGen G) hC2)
        (unitHom G (VOA.Huang.SchurWeyl.V2m_orbifold_finGen G)) a →
      ∃ π : G → (X ⟶ X),
        (∀ g, a ≫ π g = V2mObj G ◁ π g ≫ a) ∧
        (∀ g, V2mObj G ◁ π g ≫ monodromy (V2mObj G) X ≫
            autHom G (VOA.Huang.SchurWeyl.V2m_orbifold_finGen G) g ▷ X ≫ a =
          V2mObj G ◁ π g ≫ a) ∧
        (∀ g h, π g ≫ π h = if g = h then π g else 0) ∧
        ∑ g, π g = 𝟙 X := by
\end{leancode}

\begin{dictionary}
In \lean{theorem_3_3} the data are section variables: \lean{(D : McRaeData A m u e G act coev)} and a module
\lean{(hX : IsLeftModule m u a)}. Clause one: $\pi_g$ is a module map. Clause two: on the image of $\pi_g$ the action is
\lean{act g}-twisted (the twisted condition precomposed with $1\otimes\pi_g$, which is how one says ``the image is
twisted'' without forming the image). Clauses three and four: a complete family of orthogonal idempotents.
\textbf{Weaker in form than McRae's statement, same in content:} McRae states $X\cong\bigoplus_gX_g$ with $X_g$
$g$-twisted; the library gives the idempotents, because $\mathcal C$ has no images or biproducts as Lean instances.
Given the idempotents, the decomposition is formal; and Steps~5--6 consume $\pi_g$ directly. Also proved: naturality
of $\pi_g$ in module maps (\lean{proj_natural}) and that there are no nonzero module maps between a $g_1$-twisted and a
$g_2$-twisted module for $g_1\ne g_2$ (\lean{twisted_hom_eq_zero_orbifold}). In \lean{V2m_theorem_3_3},
\lean{V2mObj G} is \lean{vObj} for $V_2$ and the only hypothesis left is \lean{hC2}.
\end{dictionary}

\begin{proof}[Proof on paper]
Write $\theta_g=a\circ(g\otimes1)\circ M_{A,X}:A\otimes X\to X$ (``act after twisting by $g$ and carrying around $X$'').
McRae's projection is
\begin{equation*}
  \pi_g=\frac1{|G|}\;a\circ(1\otimes\theta_g)\circ\alpha\circ(\tilde\imath\otimes1)\circ\lambda^{-1},
  \qquad\text{symbolically}\quad \pi_g(x)=\frac1{|G|}\sum_i x_i\cdot\theta_g(y_i\otimes x).
\end{equation*}
\emph{Sum.} $\sum_g\theta_g(y\otimes x)$ involves $\sum_g g(y)=|G|\,\iota\varepsilon(y)$, and the monodromy of the unit is
trivial, so $\sum_g\pi_g(x)=\sum_ix_i\varepsilon(y_i)\cdot x=x$ by the zigzag identities of Step~4b. \emph{Module map.} The Casimir
element $\tilde\imath$ commutes with the action up to the braiding: $\sum_i ax_i\otimes y_i$ and $\sum_ix_i\otimes y_ia$ agree
after multiplication (McRae's Lemma~3.7; proved from the zigzag identities of Step~4b), which, with the associativity of the monodromy
action $M\circ a$ (a consequence of the pure-braid relation $\sigma_2^2\sigma_1\sigma_2^2\sigma_1=\sigma_2\sigma_1^2\sigma_2\sigma_1^2$
and commutativity), gives that $\pi_g$ commutes with $a$ and that its image is $g$-twisted. \emph{Orthogonality.}
On the image of $\pi_h$ one has $(1\otimes\pi_h)\circ\theta_g$ expressed through $\theta_{h^{-1}g}$; a trace lemma using the
Galois condition $\sum_ix_ig(y_i)=|G|\delta_{g,1}$ (McRae's Lemmas~3.6 and~3.8) then gives $\pi_h\pi_g=\delta_{g,h}\pi_h$.
\end{proof}

\begin{physics}
Every sector of the parent theory that is compatible with the condensed algebra decomposes uniquely into twisted sectors,
one for each group element: the projection $\pi_g$ is the ``twist-$g$ projector'' of orbifold CFT, built from the
operator that inserts the parent algebra and averages over the symmetry.
\end{physics}

\begin{status}
\lean{theorem_3_3}: proved, Lean's three axioms (pure category theory). \lean{V2m_theorem_3_3}: proved, conditional on
\lean{hC2}; Lean's three axioms (it uses the theorem \lean{VOA.pbw}; while \lean{VOA.pbw} was an axiom the report also listed \lean{VOA.pbw}) (audit log of Step~4c).
\end{status}

%----------------------------------------------------------------------
\subsection{Step 5: a twisted part is an untwisted module of the cyclic orbifold}

This step is where the analysis of logarithmic intertwining operators meets the categorical structure, and where the
specific VOA $V_2$ enters. Fix a left $V$-module $(X,a)$ in $\mathcal C$ and $g\in G$, and let $p$ be a module
endomorphism of $X$ whose image is $g$-twisted (for instance $p=\pi_g$). The action $a:V\boxtimes X\to X$ corresponds,
by the universal property of $\boxtimes$, to a logarithmic intertwining operator of type $\binom{X}{V\,X}$ over $V^G$,
\begin{equation*}
  \mathcal Y_a(v,z)x=\sum_{n\in\C}\sum_{k\ge0}v_{(n;k)}x\;z^{-n-1}(\log z)^k ,
\end{equation*}
written \lean{actY hC2 a} in Lean (\lean{actY a = a ∘ 𝒴_⊠}, \leanpath{McRaeO5Twisted.lean}; the mode
\lean{(actY hC2 a).1 n k v x} is $v_{(n;k)}x$, the coefficient of $z^{-n-1}(\log z)^k$).

\paragraph{Step 5a: monodromy kills logarithms and non-integral powers.} The braiding is $\Omega_0$, and $\Omega_0\circ\Omega_0$ is the
operator that continues $\mathcal Y$ once around $z=0$: $z\mapsto e^{2\pi i}z$, which multiplies $z^{-n-1}$ by
$e^{-2\pi i(n+1)}$ and replaces $\log z$ by $\log z+2\pi i$. On modes,
\begin{equation*}
  \big(\Omega_0\Omega_0\mathcal Y\big)_{(n;k)}=e^{-2\pi i(n+1)}\sum_{j\ge k}\binom jk(2\pi i)^{j-k}\mathcal Y_{(n;j)} .
\end{equation*}
If $v\in V$ is fixed by $g$, the twisted condition on the image of $p$ says that this monodromy acts trivially on
$\mathcal Y_a(v,z)$ restricted to $p(X)$. A family of modes fixed by this unipotent-times-phase operator must have
$e^{-2\pi i(n+1)}=1$ (so $n\in\Z$) and no component with $k\ge1$ (the top logarithmic power would be multiplied by
$1$ but the next one would pick up $2\pi i$ times it). Hence, on $p(X)$, \emph{the vertex operators of $g$-fixed vectors
have integral powers and no logarithms} (\lean{twisted_modes}).

\paragraph{Step 5b: the Jacobi identity (the ``Jacobi bridge'').} To make $p(X)$ a module for the cyclic orbifold $V^{\langle g\rangle}$
(the $g$-fixed vectors of $V$), one needs the Jacobi identity for these integral, log-free modes. This is a named
\lean{Prop}.

\begin{definition}[The twisted Jacobi bridge]\label{l22:def:bridge}
For all left $V$-modules $(X,a)$ in $\mathcal C$, all $g\in G$, all module endomorphisms $p$ with $g$-twisted image, all
$g$-fixed $u,v\in V$, all $x\in X$ and all $\ell,m,n\in\Z$, the modes $w_{(k)}:=w_{(k;0)}$ of $\mathcal Y_a$ satisfy on
$p(x)$ the Jacobi (Borcherds) identity
\begin{equation*}
  \sum_{i\ge0}(-1)^i\binom{\ell}{i}\Big(u_{(\ell+m-i)}v_{(n+i)}-(-1)^{\ell}\,v_{(\ell+n-i)}u_{(m+i)}\Big)p(x)
  =\sum_{i\ge0}\binom mi\big(u_{(\ell+i)}v\big)_{(m+n-i)}p(x).
\end{equation*}
\end{definition}
\noindent File \leanpath{VOA/\allowbreak Classification/\allowbreak COne/\allowbreak McRaeO5Main.lean}.
\begin{leancode}
def TwistedJacobiBridge : Prop :=
  letI := logarithmicMonoidalCategory hC2
  letI := logarithmicBraidedCategory hC2
  ∀ (X : LogModuleCat (orbifoldVOA G)) (a : vObj G hfg ⊗ X ⟶ X),
    IsLeftModule (mulHom G hfg hC2) (unitHom G hfg) a →
    ∀ (g : G) (p : X ⟶ X), a ≫ p = vObj G hfg ◁ p ≫ a →
      vObj G hfg ◁ p ≫ monodromy (vObj G hfg) X ≫ autHom G hfg g ▷ X ≫ a =
        vObj G hfg ◁ p ≫ a →
      ∀ u v : V, (g : Aut W) u = u → (g : Aut W) v = v → ∀ (x : X) (ℓ m n : ℤ),
        (∑ᶠ i : ℤ, (negOnePow i * cbz ℓ i) •
            ((actY hC2 a).1 ((ℓ + m - i : ℤ) : ℂ) 0 u
                ((actY hC2 a).1 ((n + i : ℤ) : ℂ) 0 v (p.1 x))
              - negOnePow ℓ • (actY hC2 a).1 ((ℓ + n - i : ℤ) : ℂ) 0 v
                  ((actY hC2 a).1 ((m + i : ℤ) : ℂ) 0 u (p.1 x))))
          = ∑ᶠ i : ℤ, cbz m i •
              (actY hC2 a).1 ((m + n - i : ℤ) : ℂ) 0 (W.smul u (ℓ + i) v) (p.1 x)
\end{leancode}

\begin{dictionary}
Section variables \lean{G}, \lean{hfg}, \lean{hC2} as in Step~3. The two hypotheses on \lean{p} are ``module map''
(\lean{a ≫ p = V ◁ p ≫ a}) and ``$g$-twisted image'' (the twisted condition of Definition~\ref{l22:def:twisted}
precomposed with $1\boxtimes p$). \lean{p.1} is the underlying linear map, \lean{(g : Aut W) u = u} says $u$ is
$g$-fixed. \lean{negOnePow i} is $(-1)^i$; \lean{cbz ℓ i} is the binomial coefficient $\binom{\ell}{i}$ for $i\ge0$ (any
integer $\ell$, the generalized binomial coefficient) and $0$ for $i<0$, so the sums over $i\in\Z$ are the sums over
$i\ge0$; \lean{∑ᶠ} is a sum with finitely many nonzero terms (the truncation property makes the sums finite).
\lean{(actY hC2 a).1 ((k : ℤ) : ℂ) 0} is the mode $(\cdot)_{(k;0)}$ at an \emph{integral} index with \emph{no}
logarithm; \lean{W.smul u (ℓ + i) v} is $u_{(\ell+i)}v$ computed in the ambient VOA. So the display is exactly the
Borcherds identity in the form of Chapter~\ref{ch:L03}. The statement quantifies over \emph{all} such data and is
therefore strong; a referee should note that it is about the integral log-free modes only, which by Step~5a are all the
modes that act on $p(X)$ for $g$-fixed vectors. (Its docstring still says ``open''; it is proved,
Theorem~\ref{l22:thm:bridge}; Section~\ref{l22:sec:defects}.)
\end{dictionary}

\begin{physics}
In a $g$-twisted sector, the parent fields invariant under $g$ are single-valued around the insertion point, and their
operator product expansions close; the Jacobi identity is the statement that these OPEs are associative and local, so
that the twisted sector is an ordinary sector of the smaller gauged theory $V^{\langle g\rangle}$.
\end{physics}

\begin{status}
Definition. It was the last open \lean{Prop} of the converse until 2026-10-14; now proved (Theorem~\ref{l22:thm:bridge}).
\end{status}

\begin{theorem}[Step 5b: the Jacobi bridge holds]\label{l22:thm:bridge}
For every VOA $V$, every subgroup $G\le\Aut V$ (finite or not), every proof of finite generation and of
$C_2$-cofiniteness of $V^G$, the twisted Jacobi bridge holds.
\end{theorem}
\noindent File \leanpath{VOA/\allowbreak Classification/\allowbreak COne/\allowbreak McRaeO5bMain.lean}.
\begin{leancode}
theorem twistedJacobiBridge_holds (G : Subgroup (Aut W)) (hfg : FinGen (orbifold G).adjointRestrict)
    (hC2 : C2Cofinite (orbifoldVOA G)) : TwistedJacobiBridge G hfg hC2 := by
\end{leancode}
The analytic heart is a statement about one holomorphic function, independent of vertex algebras:
\noindent File \leanpath{VOA/\allowbreak Classification/\allowbreak COne/\allowbreak McRaeO5bLaurent.lean}.
\begin{leancode}
theorem residue_of_glue {s t q : ℤ → ℂ} {n₀ n₁ m₁ : ℤ}
    (hs0 : ∀ n, n₀ < n → s n = 0) (ht0 : ∀ n, n < n₁ → t n = 0) (hq0 : ∀ m, m < m₁ → q m = 0)
    {Ψ : ℂ → ℂ} (hΨ : ∀ ζ : ℂ, ζ ≠ 0 → ζ ≠ 1 → DifferentiableAt ℂ Ψ ζ)
    (hS : ∀ ζ : ℂ, 1 < ‖ζ‖ → HasSum (fun n : ℤ => s n * ζ ^ n) (Ψ ζ))
    (hT : ∀ ζ : ℂ, ζ ≠ 0 → ‖ζ‖ < 1 → HasSum (fun n : ℤ => t n * ζ ^ n) (Ψ ζ))
    (hQ : ∀ r : ℂ, r ≠ 0 → ‖r‖ < 1 → HasSum (fun m : ℤ => q m * r ^ m) (Ψ (1 + r)))
    (a ℓ : ℤ) : resInf a ℓ s - resZero a ℓ t = resOne a ℓ q := by
\end{leancode}

\begin{dictionary}
\lean{twistedJacobiBridge_holds} has no hypotheses beyond the parameters of the \lean{Prop}; in particular neither
simplicity of $V$ nor finiteness of $G$ is assumed. In \lean{residue_of_glue}: $\Psi$ is holomorphic on
$\C\setminus\{0,1\}$ (\lean{DifferentiableAt ℂ}); it has three Laurent expansions, $\sum_ns_n\zeta^n$ for $|\zeta|>1$
with $s_n=0$ for $n>n_0$ (a pole at $\infty$ of finite order), $\sum_nt_n\zeta^n$ for $0<|\zeta|<1$ with $t_n=0$ for
$n<n_1$, and $\sum_mq_m r^m$ in $r=\zeta-1$ for $0<|r|<1$ with $q_m=0$ for $m<m_1$. \lean{resInf a ℓ s},
\lean{resZero a ℓ t}, \lean{resOne a ℓ q} are the coefficients of $\zeta^{-1}$ (resp.\ $r^{-1}$) of
$\zeta^a(\zeta-1)^{\ell}\Psi$ in the three expansions (\leanpath{McRaeO5bResidue.lean}; e.g.\ \lean{resOne a ℓ q}
$=\sum_j\binom aj q_{-1-\ell-j}$). The conclusion is the residue theorem on the Riemann sphere,
$\operatorname{Res}_\infty+\operatorname{Res}_0+\operatorname{Res}_1=0$, in the sign conventions of the Jacobi identity.
\end{dictionary}

\begin{proof}[Proof on paper]
Fix $(X,a)$, $g$, $p$, $g$-fixed homogeneous $u,v$, a vector $x=p(x')$ in the image of $p$ with $x'$ of generalized
weight $c$, a linear functional $f$ on $X$, and an integer $D$, the \emph{index line}: in every term below the two mode
indices add up to $D$ (for instance $(-n-1)+(D+n+1)=D$), so each channel is a single sequence indexed by one integer.
(For the Jacobi identity with parameters $\ell,m,n$ one takes $D=\ell+m+n$; the modes are the integral, log-free modes
$w_{(k)}=w_{(k;0)}$ of $\mathcal Y_a$, and in the iterate $u_{(-m-1)}v$ is the product in $V$.) Consider three
\emph{channels}, all functions of one complex variable $\zeta$ ($u$ inserted at $\zeta$, $v$ at $1$, $x$ at $0$), each
single-valued by Step~5a:
\begin{center}\small
\begin{tabular}{lll}
channel & series & domain\\\hline
product & $\sum_n s_n\zeta^n$, \ $s_n=f\big(u_{(-n-1)}v_{(D+n+1)}x\big)$ & $|\zeta|>1$\\
reversed product & $\sum_n t_n\zeta^n$, \ $t_n=f\big(v_{(D+n+1)}u_{(-n-1)}x\big)$ & $0<|\zeta|<1$\\
iterate & $\sum_m q_m(\zeta-1)^m$, \ $q_m=f\big((u_{(-m-1)}v)_{(D+m+1)}x\big)$ & $0<|\zeta-1|<1$
\end{tabular}
\end{center}
(The intervals below are intervals of real $\zeta$.) \emph{Product $=$ iterate} on $(1,3/2)$: this is the module axiom $a\circ(\mu\boxtimes1)=a\circ(1\boxtimes a)\circ\alpha$
tested on canonical elements; it is the converse direction of the associativity proof of Step~3. \emph{Reversed product
$=$ iterate} on $(1/2,1)$: here the commutativity of the algebra enters --- skew-symmetry $\Omega_0(Y_V)=Y_V$, an inner
translation theorem for the iterate, and the module axiom at scale $\zeta$. \emph{Gluing.} The product row satisfies a
rational Fuchsian differential equation with its only finite singularity at $1$, so it continues analytically to
$\C\setminus[1,\infty)$ in $1/\zeta$; with the other two channels and the identity theorem on convex pieces this produces
one function $\Psi$ holomorphic on $\C\setminus\{0,1\}$ with the three expansions. \emph{Residues.} Since the pole at
$1$ has finite order $N$, $(\zeta-1)^N\Psi$ is bounded near $1$, and a Liouville argument shows $\Psi$ is a Laurent
polynomial divided by $(\zeta-1)^N$; uniqueness of Laurent coefficients identifies $s,t,q$ with the binomial
expansions of that rational function at $\infty$, $0$, $1$, and the residue theorem (a Vandermonde-convolution
identity for binomial coefficients) gives \lean{resInf - resZero = resOne}. Choosing the exponents $a=m$ and $\ell$
(the $a$ of \lean{residue_of_glue}, not the action) and $D=\ell+m+n$ turns this into $f$ applied to the Jacobi
identity for homogeneous data; since $f$ is arbitrary, the identity itself holds; additivity (in $x$ along generalized
$L_0$-eigenspaces, in $u,v$ along the grading of $V^{\langle g\rangle}$) gives the general case. This is the classical
derivation of the Jacobi identity from rationality, commutativity and associativity of correlation functions (Frenkel--Lepowsky--Meurman,
Frenkel--Huang--Lepowsky), carried out on a twisted part.
\end{proof}

\begin{physics}
The three channels are the three operator-product limits of one chiral three-point function
$\langle f,\,u(\zeta)\,v(1)\,x(0)\rangle$: $u$ near infinity, near $0$, near $v$. Crossing symmetry says they are one
meromorphic function on the sphere, and the Jacobi identity is the statement that the residues of
$\zeta^a(\zeta-1)^\ell$ times this function sum to zero.
\end{physics}

\begin{status}
\lean{twistedJacobiBridge_holds}: proved, no hypotheses; Lean's three axioms (re-probed at an earlier commit, not re-run for this edition; while
\lean{VOA.pbw} was an axiom the report also listed \lean{VOA.pbw}; every VOA-side lemma uses the theorem \lean{VOA.pbw} through
the tensor category). It is in the maintained audit
(\leanpath{scripts/\allowbreak audit/\allowbreak ClassificationCOneAxiomAudit.lean}) with expected axioms \lean{PBW_AXIOMS}
(now the standard three).
\lean{residue_of_glue}: proved, Lean's three axioms (audit log of Step~5b).
\end{status}

\paragraph{Step 5c: the cyclic orbifolds of $V_2$ are rational.}

\begin{theorem}[Step 5c: the cyclic orbifolds of $V_2$ are rational]\label{l22:thm:cyclic}
For every automorphism $g$ of $V_2$ of finite order, $V_2^{\langle g\rangle}$ is (admissibly) rational.
\end{theorem}
\noindent File \leanpath{VOA/\allowbreak Classification/\allowbreak COne/\allowbreak McRaeO5Cyclic.lean}.
\begin{leancode}
theorem V2m_cyclic_admissiblyRational (g : Aut (V2m 1 le_rfl)) (hg : IsOfFinOrder g) :
    AdmissiblyRational (orbifoldVOA (Subgroup.zpowers g)) :=
\end{leancode}

\begin{dictionary}
\lean{IsOfFinOrder g}: $g^n=1$ for some $n\ge1$. \lean{Subgroup.zpowers g} is the cyclic group $\langle g\rangle$.
\lean{AdmissiblyRational} (Chapter~\ref{ch:L12}) says that in every admissible module every weak submodule has a weak
complement. This is the library's rationality convention, whose relation to Dong--Li--Mason rationality is not settled
in Lean (Remark~\ref{l22:rem:rational}); so the theorem is ``$V_2^{\langle g\rangle}$ is admissibly rational'', not
literally ``DLM-rational''.
\end{dictionary}

\begin{proof}[Proof on paper]
$\Aut V_2\cong\mathrm{PSL}_2(\C)$; an element of finite order is semisimple, hence conjugate to an element of the
maximal torus $\{e^{2\pi i\,t\,\beta(0)}\}$ of $\Aut V_{\sqrt2\Z}$ ($\beta$ the generator of the lattice, $\beta^2=2$).
A torus element of order $n$ acts on the charge-$k$ sector $M(1)\otimes e^{k\beta}$ by $\zeta^{k}$ for a primitive
$n$-th root of unity $\zeta$, and trivially on the Heisenberg part. Conjugate subgroups have isomorphic fixed-point algebras, and the fixed points of the torus element are the
sectors with $n\mid k$, i.e.\ the lattice VOA of $n\sqrt2\Z$:
\begin{equation*}
  V_2^{\langle g\rangle}\;\cong\;V_{\sqrt{2n^2}\,\Z}=V_{2n^2}.
\end{equation*}
Rank-one positive-definite even lattice VOAs are rational (Theorem~\ref{l22:thm:V2mnice}), and rationality is
transported along isomorphisms.
\end{proof}

\begin{physics}
Gauging a cyclic $\Z_n$ subgroup of the $SU(2)$ symmetry of the self-dual boson moves it along the circle line to the
rational radius corresponding to $V_{2n^2}$: a compact boson, again rational. This is the single place where the
specific parent theory matters.
\end{physics}

\begin{status}
Proved; Lean's three axioms (audit log of Step~5).
\end{status}

\paragraph{Step 5d: sections on twisted parts.} Combining Steps~5a, 5b and~5c: on each $\pi_g$-part of a $V_2$-module, the modes
of $V_2^{\langle g\rangle}$ satisfy the Jacobi identity, the module is admissible (its generalized $L_0$-eigenspaces span
and its spectrum is bounded below by that of $X$), so by Step~5c it is completely reducible. Hence a surjective $V_2$-module
map restricted to the $g$-parts has a $V_2^{\langle g\rangle}$-linear section, which is in particular $V_2^G$-linear.

\begin{theorem}[Step 5d: surjections of $V_2$-modules split $V_2^G$-linearly]\label{l22:thm:O5}
Let $G\le\Aut V_2$ be finite with $V_2^G$ $C_2$-cofinite, and assume the bridge. If $f:X_1\to X_2$ is a surjective map of
left $V_2$-modules in $\mathcal C$, there is a morphism $s:X_2\to X_1$ of $\mathcal C$ (a $V_2^G$-module map) with
$f\circ s=\mathrm{id}$.
\end{theorem}
\noindent File \leanpath{VOA/\allowbreak Classification/\allowbreak COne/\allowbreak McRaeO5Main.lean}.
\begin{leancode}
theorem V2m_exists_section (G : Subgroup (Aut (V2m 1 le_rfl))) [Fintype G] [DecidableEq G]
    (hC2 : C2Cofinite (orbifoldVOA G))
    (hB : TwistedJacobiBridge G (VOA.Huang.SchurWeyl.V2m_orbifold_finGen G) hC2) :
    letI := logarithmicMonoidalCategory hC2
    letI := logarithmicBraidedCategory hC2
    ∀ (X₁ X₂ : LogModuleCat (orbifoldVOA G)) (a₁ : V2mObj G ⊗ X₁ ⟶ X₁)
      (a₂ : V2mObj G ⊗ X₂ ⟶ X₂),
      IsLeftModule (mulHom G (VOA.Huang.SchurWeyl.V2m_orbifold_finGen G) hC2)
        (unitHom G (VOA.Huang.SchurWeyl.V2m_orbifold_finGen G)) a₁ →
      IsLeftModule (mulHom G (VOA.Huang.SchurWeyl.V2m_orbifold_finGen G) hC2)
        (unitHom G (VOA.Huang.SchurWeyl.V2m_orbifold_finGen G)) a₂ →
      ∀ f : X₁ ⟶ X₂, IsModuleHom a₁ a₂ f → Function.Surjective f.1 →
        ∃ s : X₂ ⟶ X₁, s ≫ f = 𝟙 X₂ :=
\end{leancode}

\begin{dictionary}
\lean{IsModuleHom a₁ a₂ f} is $f\circ a_1=a_2\circ(1\boxtimes f)$; \lean{Function.Surjective f.1}: the underlying linear
map is onto. The conclusion \lean{s ≫ f = 𝟙 X₂} is $f\circ s=\mathrm{id}$. \textbf{Note:} $s$ is only a morphism of
$\mathcal C$, not a $V_2$-module map; that is all Step~6 needs. The hypothesis \lean{hB} is now discharged by
Theorem~\ref{l22:thm:bridge}. The section is $s=\sum_gs_g\circ\pi_g$ with $s_g$ the section on the $g$-part
(\lean{V2m_exists_twisted_section}).
\end{dictionary}

\begin{physics}
A short exact sequence of sectors of the parent theory (with all twists allowed) always splits as a sequence of sectors
of the gauged theory: no logarithmic gluing survives on any twisted sector, because each is a sector of a rational
compact boson.
\end{physics}

\begin{status}
Proved, conditional on \lean{hC2} and \lean{hB}; \lean{hB} is now proved. Lean's three axioms (it uses the theorem \lean{VOA.pbw}; while \lean{VOA.pbw} was an axiom the report also listed \lean{VOA.pbw}) (audit
log of Step~5).
\end{status}

%----------------------------------------------------------------------
\subsection{Step 6: McRae's Lemma 4.21, every surjection in \texorpdfstring{$\mathcal C$}{C} splits}

\begin{theorem}[Step 6: every surjection splits; McRae Lemma~4.21 for $V_2$]\label{l22:thm:O6}
Let $G\le\Aut V_2$ be finite with $V_2^G$ $C_2$-cofinite. Every morphism $f:Y_1\to Y_2$ of $\mathcal C$ whose underlying
map is surjective has a section $\sigma$, $f\circ\sigma=\mathrm{id}_{Y_2}$.
\end{theorem}
\noindent File \leanpath{VOA/\allowbreak Classification/\allowbreak COne/\allowbreak McRaeO6Main.lean}.
\begin{leancode}
theorem V2m_lemma_4_21 (G : Subgroup (Aut (V2m 1 le_rfl))) [Fintype G] [DecidableEq G]
    (hC2 : C2Cofinite (orbifoldVOA G))
    (hB : TwistedJacobiBridge G (VOA.Huang.SchurWeyl.V2m_orbifold_finGen G) hC2) :
    ∀ {Y₁ Y₂ : LogModuleCat (orbifoldVOA G)} (f : Y₁ ⟶ Y₂), Function.Surjective f.1 →
      ∃ σ : Y₂ ⟶ Y₁, σ ≫ f = 𝟙 Y₂ :=
\end{leancode}

\begin{dictionary}
Now $Y_1,Y_2$ are arbitrary objects of $\mathcal C$ (no $V_2$-action), and $\sigma$ is a morphism of $\mathcal C$. The
hypothesis \lean{hB} is discharged by Theorem~\ref{l22:thm:bridge}. Every epimorphism of $\mathcal C$ has surjective
underlying map (\lean{surjective_of_epi}, which needed quotient objects in $\mathcal C$, \leanpath{McRaeO6Quotient.lean}),
so the hypothesis ``surjective'' is the categorical ``epi''.
\end{dictionary}

\begin{proof}[Proof on paper]
Induce: $V\boxtimes f:V\boxtimes Y_1\to V\boxtimes Y_2$ is a map of left $V$-modules (action
$(\mu\boxtimes1)\circ\alpha^{-1}$), and it is an epimorphism (right exactness, Step~3), hence surjective. By Step~5 there is a
$\mathcal C$-morphism $s$ with $(V\boxtimes f)\circ s=\mathrm{id}$. Put
\begin{equation*}
  \sigma=\lambda\circ(\varepsilon\boxtimes1)\circ s\circ(\iota\boxtimes1)\circ\lambda^{-1}:Y_2\to Y_1 .
\end{equation*}
Then, by naturality of $\lambda$ and the interchange law,
$f\circ\sigma=\lambda\circ(\varepsilon\boxtimes1)\circ(V\boxtimes f)\circ s\circ(\iota\boxtimes1)\circ\lambda^{-1}
=\lambda\circ((\varepsilon\circ\iota)\boxtimes1)\circ\lambda^{-1}=\mathrm{id}$, since $\varepsilon\circ\iota=\mathrm{id}$.
This works in any monoidal category (\lean{section_of_whisker}).
\end{proof}

\begin{physics}
Every sector of the gauged theory is a summand of an induced sector (``fuse with the parent algebra and project back by
averaging''), and induced sectors split by Step~5; so the gauged theory has no logarithmic extensions among its sectors.
\end{physics}

\begin{status}
Proved, conditional on \lean{hC2}, \lean{hB} (now proved). Lean's three axioms (it uses the theorem \lean{VOA.pbw}; while \lean{VOA.pbw} was an axiom the report also listed \lean{VOA.pbw}) (O67 audit log).
\end{status}

%----------------------------------------------------------------------
\subsection{Step 7: split surjections imply rationality, and the rationality of the orbifolds}

\begin{theorem}[Step 7: split surjections imply rationality, general form]\label{l22:thm:O7}
Let $W$ be a $C_2$-cofinite VOA such that every surjective morphism of $\mathcal C_W=$ \lean{LogModuleCat W} has a
section. Then $W$ is admissibly rational.
\end{theorem}
\noindent File \leanpath{VOA/\allowbreak Classification/\allowbreak COne/\allowbreak McRaeO7Main.lean}.
\begin{leancode}
def SurjectionsSplit (W : VertexOperatorAlgebra V) : Prop :=
  ∀ {Y₁ Y₂ : LogModuleCat W} (f : Y₁ ⟶ Y₂), Function.Surjective f.1 →
    ∃ σ : Y₂ ⟶ Y₁, σ ≫ f = 𝟙 Y₂
\end{leancode}
\noindent File \leanpath{VOA/\allowbreak Classification/\allowbreak COne/\allowbreak McRaeO7Main.lean}.
\begin{leancode}
theorem admissiblyRational_of_surjectionsSplit (hC2 : C2Cofinite W) (hs : SurjectionsSplit W) :
    AdmissiblyRational W := by
\end{leancode}

\begin{dictionary}
\lean{SurjectionsSplit W} quantifies over objects of $\mathcal C_W$ only, i.e.\ finitely generated logarithmic modules;
the conclusion \lean{AdmissiblyRational W} is about \emph{all} admissible modules (Chapter~\ref{ch:L12}), so the theorem
is a genuine passage from a finite-length statement to complete reducibility in general. No semisimplicity of Zhu's
algebra and no positive-energy hypothesis is used.
\end{dictionary}

\begin{proof}[Proof on paper]
Let $X$ be a weak module with a simple submodule $P$ and simple quotient $X/P$. Then $X$ has length $2$, hence (by
$C_2$-cofiniteness) is a finitely generated logarithmic module, an object $N$ of $\mathcal C_W$; the quotient
$N\to N/P$ is surjective, so it has a section $\sigma$, and $\sigma(N/P)$ is a complement of $P$. So every extension of a
simple module by a simple module splits. The rest is a lattice-theoretic induction that the library already had
from the proof of rationality of $V_{2m}^+$: for a $C_2$-cofinite VOA, every two-step interval $a<b<c$ in the submodule
lattice of a weak module can be reduced, through its ``factor bottoms'' (lowest-weight pieces of the simple factors,
which exist by $C_2$-cofiniteness), to a module of length two, which splits; so every two-step interval has a relative
complement. A module in which all two-step intervals of all cyclic submodules split is semisimple
(\lean{semisimple_of_generalized_cyclic_splitting}). Hence every admissible module is completely reducible.
\end{proof}

\begin{physics}
A lisse chiral algebra is rational as soon as no two primaries can be glued into a logarithmic pair: logarithmic
partners always appear first in extensions of length two.
\end{physics}

\begin{status}
Proved; Lean's three axioms, no \lean{VOA.pbw} (O67 audit log).
\end{status}

\begin{theorem}[Step 7 for $V_2$: orbifold rationality from the Jacobi bridge]\label{l22:thm:O7V2}
If the bridge holds for every finite $G\le\Aut V_2$ with $V_2^G$ $C_2$-cofinite, then the rationality of the exceptional orbifolds (fact~(d)) holds, and the converse of the
classification holds.
\end{theorem}
\noindent File \leanpath{VOA/\allowbreak Classification/\allowbreak COne/\allowbreak McRaeO7Main.lean}.
\begin{leancode}
def TwistedJacobiBridgeAll : Prop :=
  ∀ (G : Subgroup (Aut (V2m 1 le_rfl))) [Finite G] (hC2 : C2Cofinite (orbifoldVOA G)),
    McRaeO5.TwistedJacobiBridge G (VOA.Huang.SchurWeyl.V2m_orbifold_finGen G) hC2
\end{leancode}
\noindent File \leanpath{VOA/\allowbreak Classification/\allowbreak COne/\allowbreak McRaeO7Main.lean}.
\begin{leancode}
theorem finiteV2OrbifoldRationalOfC2_of_bridge (hB : TwistedJacobiBridgeAll) :
    FiniteV2OrbifoldRationalOfC2Fact := by
\end{leancode}
\noindent File \leanpath{VOA/\allowbreak Classification/\allowbreak COne/\allowbreak McRaeO7Main.lean}.
\begin{leancode}
theorem gannonRayhaunConverse_of_bridge (hB : TwistedJacobiBridgeAll) :
    GannonRayhaunConverse.{0} :=
  gannonRayhaunConverse_of_B3 (finiteV2OrbifoldRationalOfC2_of_bridge hB)
\end{leancode}

\begin{dictionary}
\lean{TwistedJacobiBridgeAll} is the bridge of Definition~\ref{l22:def:bridge} for every finite $G\le\Aut V_2$ with
$C_2$-cofinite $V_2^G$. The second theorem is Theorem~\ref{l22:thm:O6} followed by Theorem~\ref{l22:thm:O7}; the third
is the historical reduction Theorem~\ref{l22:thm:converseB3} applied to it. (The docstrings of
\lean{TwistedJacobiBridgeAll} and of the file still call it ``open''; it is proved, Theorem~\ref{l22:thm:B3}.)
\end{dictionary}

\begin{physics}
Gauging any finite symmetry of $SU(2)_1$ with a lisse result gives a rational theory --- provided twisted sectors behave
as sectors of the cyclic orbifolds, which is the content of the bridge.
\end{physics}

\begin{status}
Proved, conditional on \lean{hB}, which is now proved. Lean's three axioms (it uses the theorem \lean{VOA.pbw}; while \lean{VOA.pbw} was an axiom the report also listed \lean{VOA.pbw}) (O67 audit log).
\end{status}

\begin{theorem}[Rationality of the orbifolds of $V_2$: McRae's Corollary 4.23]\label{l22:thm:B3}
For every finite subgroup $G\le\Aut V_2$ with $V_2^G$ $C_2$-cofinite, $V_2^G$ is admissibly rational.
\end{theorem}
\noindent File \leanpath{VOA/\allowbreak Classification/\allowbreak COne/\allowbreak MainConverseFinal.lean}.
\begin{leancode}
theorem twistedJacobiBridgeAll_holds : McRaeO7.TwistedJacobiBridgeAll := by
  intro G _ hC2
  exact McRaeO5b.twistedJacobiBridge_holds G _ hC2
\end{leancode}
\noindent File \leanpath{VOA/\allowbreak Classification/\allowbreak COne/\allowbreak MainConverseFinal.lean}.
\begin{leancode}
theorem finiteV2OrbifoldRationalOfC2 : FiniteV2OrbifoldRationalOfC2Fact :=
  McRaeO7.finiteV2OrbifoldRationalOfC2_of_bridge twistedJacobiBridgeAll_holds
\end{leancode}

\begin{dictionary}
\lean{FiniteV2OrbifoldRationalOfC2Fact} is Definition~\ref{l22:def:Bfacts}: \lean{∀ (G : Subgroup (Aut (V2m 1 le_rfl)))
[Finite G], C2Cofinite (orbifoldVOA G) → AdmissiblyRational (orbifoldVOA G)}. The theorem has no hypotheses. It is the
statement at its natural generality (all finite subgroups), not only the three exceptional ones the converse consumes.
Its conclusion is admissible rationality, which is the rationality clause of \lean{StronglyRational}; the other clauses of
niceness of the exceptional orbifolds are supplied by the converse reduction (Theorem~\ref{l22:thm:converseB3}).
\end{dictionary}

\begin{proof}[Proof on paper]
Steps~1--7 above: the Jacobi bridge (Step~5b) for every such $G$, then Step~7 for $V_2$.
\end{proof}

\begin{physics}
Every lisse gauging of a finite symmetry of the $SU(2)_1$ WZW model (the self-dual compact boson) is a rational CFT.
\end{physics}

\begin{status}
\textbf{Proved with no hypotheses.} \lean{#print axioms}: Lean's three (re-probed at an earlier commit, not re-run for this edition); on 2026-10-14, when \lean{VOA.pbw} was an
axiom, the coordinator machine-checked Lean's three plus \lean{VOA.pbw} (aggregate \texttt{lake build} then green, 13,306
jobs; green again on 2026-10-28, 18{,}901 jobs, and later 18{,}907, in pre-release builds; for the release build see Chapter~\ref{ch:15}, Section~\ref{ch15:sec:record}), recorded in the maintained audit with expected axioms \lean{PBW_AXIOMS} (now the standard three) (\texttt{run\_axiom\_audits.py} passes, 5,890 reports, of which 854 for
\texttt{ClassificationCOne}).
\end{status}

%----------------------------------------------------------------------
\subsection{Narrowing orbifold rationality to the exceptional groups (historical)}

Before the proof of orbifold rationality was found, the library isolated the part of it that the converse consumes.
These statements are proved and remain in the library; with orbifold rationality proved, their hypotheses are all
discharged.

\begin{theorem}[The converse needs orbifold rationality only at the exceptional groups]\label{l22:thm:narrowed}
Orbifold rationality \textup{(}fact~(d) of Definition~\ref{l22:def:Bfacts}\textup{)} implies its restriction to
$G\cong A_4,S_4,A_5$; and facts (a), (b), (c) together with that restriction imply the converse.
\end{theorem}
\noindent File \leanpath{VOA/\allowbreak Classification/\allowbreak COne/\allowbreak OrbRat2Reduction.lean}.
\begin{leancode}
def ExceptionalV2OrbifoldRationalOfC2 : Prop :=
  ∀ G : Subgroup (Aut (V2m 1 le_rfl)), IsExceptionalV2Group G →
    C2Cofinite (orbifoldVOA G) → AdmissiblyRational (orbifoldVOA G)
\end{leancode}
\noindent File \leanpath{VOA/\allowbreak Classification/\allowbreak COne/\allowbreak OrbRat2Reduction.lean}.
\begin{leancode}
theorem exceptionalV2OrbifoldRationalOfC2_of_finite (h : FiniteV2OrbifoldRationalOfC2Fact) :
    ExceptionalV2OrbifoldRationalOfC2 := by
\end{leancode}
\noindent File \leanpath{VOA/\allowbreak Classification/\allowbreak COne/\allowbreak OrbRat2Reduction.lean}.
\begin{leancode}
theorem gannonRayhaunConverse_of_narrowed (hSolv : SolvableExceptionalOrbifoldC2Fact)
    (hXu : XuIcosahedralPoissonRelations) (hRat : ExceptionalV2OrbifoldRationalOfC2)
    (hPU : ExceptionalOrbifoldPositiveWeightsFact) : GannonRayhaunConverse.{0} := by
\end{leancode}

\begin{dictionary}
\lean{ExceptionalV2OrbifoldRationalOfC2} is \lean{FiniteV2OrbifoldRationalOfC2Fact} with the quantifier over finite
subgroups replaced by one over the exceptional subgroups (\lean{IsExceptionalV2Group G}, Definition~\ref{l22:def:Bfacts};
finiteness is then automatic). The first theorem shows it is \emph{no stronger} than fact~(d). The second takes four
\lean{Prop}s: facts (a) and (b), the narrowed fact~(d), and fact~(c). All four are now proved (Theorems~\ref{l22:thm:B245} and~\ref{l22:thm:B3},
the narrowed form by \lean{exceptionalV2OrbifoldRationalOfC2_of_bridge}). The same file splits the narrowed statement
further into single concrete models, $(V_8^+)^{\Z/3}$ and $(V_2^{A_4})^{\Z/2}$, and an icosahedral statement
(\lean{gannonRayhaunConverse_of_models}); that split is no longer needed.
\end{dictionary}

\begin{physics}
Only the three exceptional gaugings of $SU(2)_1$ appear in the $c=1$ list; the general theorem is more than the
classification needs.
\end{physics}

\begin{status}
Proved; \lean{#print axioms} of both theorems is Lean's three (re-checked for this revision; not listed in the
maintained audit). All hypotheses of \lean{gannonRayhaunConverse_of_narrowed} are now proved (the rationality input uses the theorem
\lean{VOA.pbw}).
\end{status}

%----------------------------------------------------------------------
\subsection{The earlier route, abandoned}\label{l22:sec:oldroute}

The first plan for proving orbifold rationality followed McRae's rationality criterion (arXiv:2108.01898): for a $C_2$-cofinite $W$ with
semisimple Zhu algebra, rationality is equivalent to every simple module $S$ having a projective cover whose socle is
$S$ itself. Its formal pieces remain in the library. They are recorded here because they are proved statements that a
reader will meet, and because their open input is \emph{not} an input of the classification.

\begin{definition}[McRae covers]\label{l22:def:mcrae}
A \emph{McRae cover} of a simple weak module $S$ is a finite-length projective weak module $P$ with a surjection
$P\to S$ with superfluous kernel and an embedding $S\hookrightarrow P$ which is the essential socle.
\lean{McRaeCoverFact W}: every simple weak $W$-module has one.
\end{definition}
\noindent File \leanpath{VOA/\allowbreak Classification/\allowbreak COne/\allowbreak OrbMcRaeCover.lean}.
\begin{leancode}
def McRaeCoverFact (W : VertexOperatorAlgebra 𝒱) : Prop :=
  ∀ {X : Type u} [AddCommGroup X] [Module ℂ X] (B : WeakModule W X),
    IsSimpleModule (modeAlg B) X → Nonempty (McRaeCover B)
\end{leancode}
\noindent File \leanpath{VOA/\allowbreak Classification/\allowbreak COne/\allowbreak OrbMcRaeCover.lean}.
\begin{leancode}
theorem mcRaeCoverFact_iff_admissiblyRational (hC2 : C2Cofinite W)
    (hss : letI := instRingA W; IsSemisimpleRing (VOA.A W)) :
    McRaeCoverFact W ↔ AdmissiblyRational W :=
\end{leancode}

\begin{dictionary}
\lean{WeakModule W X}, \lean{modeAlg B} (the associative algebra generated by the modes acting on $X$, whose
submodules are the $W$-submodules) and \lean{IsSimpleModule} are Chapter~\ref{ch:L12}'s weak modules in ring-theoretic
form. \lean{letI := instRingA W} installs the ring structure of Zhu's algebra \lean{VOA.A W}, and
\lean{IsSemisimpleRing (VOA.A W)} says it is semisimple. The equivalence says that, for $C_2$-cofinite $W$ with
semisimple Zhu algebra, \emph{rationality is exactly the existence of McRae covers}: nothing is lost by the reduction.
At the time the route was pursued, the hypothesis \lean{hss} was \emph{not} available for $V_2^{S_4}$ and $V_2^{A_5}$ (for
$V_2^{A_4}$ the library proves $A(V_2^{A_4})\cong\C^{21}$), which is one reason the route was abandoned. Today \lean{hss}
for these orbifolds follows by composing two proved theorems: orbifold rationality \lean{finiteV2OrbifoldRationalOfC2}
(Theorem~\ref{l22:thm:B3}, applied to their proved $C_2$-cofiniteness) and Zhu's semisimplicity theorem
\lean{zhuAdmissibleSemisimple_holds} (Chapter~\ref{ch:L12}, Theorem~\ref{l12:thm:litproved}). The composition is not a named
declaration of the library, and it is useless for this route, which was meant to \emph{prove} rationality: \lean{hss} obtained
from rationality makes the criterion circular.
\end{dictionary}

\begin{status}
The equivalence is proved (Lean's three axioms; maintained audit). \lean{McRaeCoverFact} for the exceptional orbifolds is
not proved in the library, and \textbf{not needed}: orbifold rationality is proved along the orbifold route (Theorem~\ref{l22:thm:B3}).
\end{status}

\begin{physics}
In a logarithmic (non-semisimple) theory, simple modules have projective covers that glue several sectors together
(Jordan blocks of $L_0$). A McRae cover whose socle is the simple module itself can only be the module itself when Zhu's
algebra is semisimple, which is how the criterion detects that no logarithmic sector exists.
\end{physics}

\begin{theorem}[Huang's projective covers exist; first step of McRae's criterion]\label{l22:thm:M1}
For every $C_2$-cofinite $W$, every simple weak module has a projective cover in the category of finite-length weak
modules (Huang 2009).
\end{theorem}
\noindent File \leanpath{VOA/\allowbreak Classification/\allowbreak COne/\allowbreak OrbMcRaeMain.lean}.
\begin{leancode}
def HuangCoverFact (W : VertexOperatorAlgebra 𝒱) : Prop :=
  ∀ {X : Type u} [AddCommGroup X] [Module ℂ X] (B : WeakModule W X),
    IsSimpleModule (modeAlg B) X → Nonempty (HuangCover B)
\end{leancode}
\noindent File \leanpath{VOA/\allowbreak Classification/\allowbreak COne/\allowbreak McRaeM1Main.lean}.
\begin{leancode}
theorem huangCoverFact_of_c2 (hC2 : C2Cofinite W) : HuangCoverFact W :=
\end{leancode}

\begin{dictionary}
\lean{HuangCover B} packages a finite-length projective weak module with a surjection onto $B$ with superfluous kernel
(the projective cover in Huang's sense); \lean{HuangCoverFact W} says every simple weak module has one. The only
hypothesis of the theorem is $C_2$-cofiniteness.
\end{dictionary}

\begin{physics}
A lisse chiral algebra, even a logarithmic one, has for each primary a ``largest indecomposable sector built on it'' (its
projective cover), the object in which all logarithmic partners of that primary are glued together.
\end{physics}

\begin{status}
Proved, Lean's three axioms (maintained audit).
\end{status}

\begin{definition}[The categorical input of the earlier route]\label{l22:def:coverdual}
\lean{CoverDualLitFact}: the contragredient of the projective cover of the contragredient of $S$ is the projective
cover of $S$ ($P'_{S'}\cong P_S$), McRae (2108.01898) Prop.~5.14 via Etingof--Nikshych--Ostrik unimodularity of a
factorizable finite tensor category.
\end{definition}
\noindent File \leanpath{VOA/\allowbreak Classification/\allowbreak COne/\allowbreak McRaeM2Main.lean}.
\begin{leancode}
def CoverDualLitFact (hC2 : C2Cofinite W) : Prop :=
  ∀ {X : Type u} [AddCommGroup X] [Module ℂ X] (B : WeakModule W X)
    (hB : IsSimpleModule (modeAlg B) X), ∀ C : HuangCover B,
    ∃ D : HuangCover (contra (gwOfFiniteLength hC2 B (finiteLength_of_simple hB))).toWeakModule,
      ∃ e : (contra (gwOfFiniteLength hC2 D.BP D.finiteLength)).toWeakModule.Hom C.BP,
        Function.Bijective e.toLinearMap
\end{leancode}

\begin{dictionary}
\lean{contra} is the contragredient module, \lean{gwOfFiniteLength} views a finite-length weak module as a generalized
weight module so the contragredient is defined; \lean{D.BP} is the projective object of the cover \lean{D}; the bijective
module map \lean{e} is the isomorphism $P'_{S'}\cong P_S$. \lean{unimodularSocleFact_of_coverDualLit} turns this into the
socle statement, and with Theorem~\ref{l22:thm:M1} into McRae covers.
\end{dictionary}

\begin{physics}
Conjugating (charge conjugation) a logarithmic projective sector and taking the conjugate of its top gives back the same
projective sector: a symmetry of logarithmic CFTs whose chiral algebra has a non-degenerate (factorizable) braiding.
\end{physics}

\begin{status}
\textbf{Open}, and \textbf{not an input of the classification} (orbifold rationality is proved without it,
Theorem~\ref{l22:thm:B3}). It is a named \lean{Prop}, not proved and not an axiom. Proving it would require the rigid, factorizable logarithmic tensor category of $W$ with unimodularity, plus
semisimplicity of the Zhu algebras of $V_2^{S_4}$, $V_2^{A_5}$ to conclude.
\end{status}

%======================================================================
\section{The classification, complete}\label{l22:sec:final}
%======================================================================

\begin{theorem}[Gannon--Rayhaun Theorem~3.1 and Proposition~3.2, with the library's rationality convention]\label{l22:thm:final}
A vertex operator algebra (with carrier in universe $0$) is nice (``rational'' read as \lean{AdmissiblyRational},
Remark~\ref{l22:rem:rational}) and has central charge $1$ if and only if it is
isomorphic to one of $V_{2m}$ ($m\ge1$), $V_{2m}^+$ ($m\ge2$), $V_2^{A_4}$, $V_2^{S_4}$, $V_2^{A_5}$; and it is
isomorphic to exactly one of them.
\end{theorem}
\noindent File \leanpath{VOA/\allowbreak Classification/\allowbreak COne/\allowbreak MainConverseFinal.lean}.
\begin{leancode}
theorem gannonRayhaunConverse_holds : GannonRayhaunConverse.{0} :=
  McRaeO7.gannonRayhaunConverse_of_bridge twistedJacobiBridgeAll_holds
\end{leancode}
\noindent File \leanpath{VOA/\allowbreak Classification/\allowbreak COne/\allowbreak MainConverseFinal.lean}.
\begin{leancode}
theorem gannonRayhaun_classification :
    GannonRayhaunStatement.{0} ∧ GannonRayhaunDistinct.{0} ∧ GannonRayhaunConverse.{0} :=
  ⟨gannonRayhaun_forward, gannonRayhaunDistinct, gannonRayhaunConverse_holds⟩
\end{leancode}

\begin{dictionary}
\textbf{No hypotheses.} The three conjuncts are the three \lean{Prop}s of Section~\ref{l22:sec:statement}:
\lean{GannonRayhaunStatement} (Definition~\ref{l22:def:GRS}: nice and $c=1$ $\Rightarrow$ \lean{InCOneList}),
\lean{GannonRayhaunDistinct} and \lean{GannonRayhaunConverse} (Definition~\ref{l22:def:labels}). Together: ``nice with
$c=1$'' $\iff$ ``in the list'', and the label is unique, which is GR's ``isomorphic to exactly one of''. The proof term is
the triple of Theorems~\ref{l22:thm:forward}, \ref{l22:thm:distinct} and \lean{gannonRayhaunConverse_holds}; the last
is the reduction of Theorem~\ref{l22:thm:O7V2} applied to the proved bridge (Theorem~\ref{l22:thm:B3}). The statement
definitions are unchanged since the snapshot (they were not modified to make the proof go through).
\textbf{Universes.} The suffix \lean{.{0}} instantiates the universe parameter of each \lean{Prop} at \lean{0}: the
quantifier ``for every VOA'' ranges over VOAs whose carrier is in \lean{Type 0}. The forward theorem
\lean{gannonRayhaun_forward : GannonRayhaunStatement.{u}} and distinctness
\lean{gannonRayhaunDistinct : GannonRayhaunDistinct.{u}} hold separately in \emph{every} universe; only the converse is
proved at universe \lean{0} (Remark~\ref{l22:rem:universe}), so the conjunction is stated there. Mathematically nothing
is lost: every VOA of the list has a model on a \lean{Type 0} carrier.
\textbf{Rationality.} ``Nice'' is the library's \lean{Nice}, whose rationality conjunct is \lean{AdmissiblyRational}; its
relation to Dong--Li--Mason rationality is not settled in Lean (Remark~\ref{l22:rem:rational}). So this is GR's theorem
exactly for Lean-niceness; for GR's own niceness the forward half follows by the on-paper implication DLM-rational
$\Rightarrow$ \lean{AdmissiblyRational}, and the converse half is proved in the library's sense only. The library's
rationality and pseudo-unitarity also quantify modules only in the universe of $V$: harmless for the forward direction
(it weakens a hypothesis), formally weaker for the converse.
\end{dictionary}

\begin{proof}[Proof on paper]
Forward: Section~\ref{l22:sec:discharge}. Distinctness: Section~\ref{l22:sec:distinct}. Converse: $V_{2m}$ and
$V_{2m}^+$ are nice (Theorems~\ref{l22:thm:V2mnice}, \ref{l22:thm:plusnice}); the exceptional orbifolds are simple,
self-dual, of CFT type with $c=1$ and $V_1=0$ by their orbifold structure, $C_2$-cofinite by facts (a) and (b) of Definition~\ref{l22:def:Bfacts},
rational by orbifold rationality (Section~\ref{l22:sec:B3}), and pseudo-unitary by the positivity of the conformal
weights (Theorem~\ref{l22:thm:converseB3}).
\end{proof}

\begin{physics}
The rational, positive-energy chiral algebras at $c=1$ are exactly: the compact boson at the rational radii, its
$\Z_2$-orbifold, and the three gaugings of the $SU(2)_1$ WZW model by the tetrahedral, octahedral and icosahedral groups
$A_4$, $S_4$, $A_5\subset\mathrm{SO}(3)$ (the images of the binary polyhedral subgroups of $SU(2)$) --- the chiral version of Ginsparg's picture of the $c=1$ moduli space, with its three isolated
points --- and these are pairwise different theories.
\end{physics}

\begin{status}
\textbf{Proved with no hypotheses.} Machine-checked by the coordinator on 2026-10-14: aggregate \texttt{lake build}
then green (13,306 jobs; green again on 2026-10-28, 18{,}901 jobs, and later 18{,}907, in pre-release builds; for the release build see Chapter~\ref{ch:15}, Section~\ref{ch15:sec:record}), with axiom list \lean{[propext, Classical.choice, Quot.sound, VOA.pbw]} while \lean{VOA.pbw} was an
axiom; re-probed at an earlier commit, not re-run for this edition:
\begin{plaincode}
'VOA.Classification.COne.gannonRayhaun_classification' depends on axioms:
  [propext, Classical.choice, Quot.sound]
\end{plaincode}
and the maintained audit (\texttt{run\_axiom\_audits.py}: 5,890 reports, 854 for \texttt{ClassificationCOne}) passes,
listing \lean{gannonRayhaun_classification}, \lean{gannonRayhaunConverse_holds}, \lean{finiteV2OrbifoldRationalOfC2},
\lean{twistedJacobiBridgeAll_holds} and \lean{McRaeO5b.twistedJacobiBridge_holds} with expected axioms
\lean{PBW_AXIOMS} (now the standard three). So the classification is \textbf{unconditional}. Distinctness alone does
not use \lean{VOA.pbw}; the forward theorem uses the theorem \lean{VOA.pbw} through the modular-data package, the
converse through the logarithmic tensor category.
\end{status}

%======================================================================
\section{Recorded traps, dead ends and Lean defects}\label{l22:sec:defects}
%======================================================================

\begin{warning}[title={Traps and dead ends of the $c=1$ campaign}]
\begin{enumerate}
\item \textbf{Characters do not pin the $S$-matrix} (Warning~\ref{l22:warn:spin}): the first modular-data package
  admitted a continuum of $S$-matrices (\lean{isModularDataWeak_perturbS}); repaired by \lean{SPinned}.
\item \textbf{The theta Galois law with $\ell^{-1}n_0$} instead of $\ell n_0$ is false (fails for $r=7$); the
  correct form is Definition~\ref{l22:def:galoislaw}.
\item \textbf{A printed constant in Dong--Jiang is wrong.} The weight-$20$ constant of Dong--Jiang's Lemma~4.4, transcribed as
  \lean{OrbC2f.CubicSeedRemainders}, is \emph{false} (\lean{OrbC2h.not_cubicSeedRemainders}; the true coefficient is
  $-447232/169744575$). The existential relation \lean{TetrahedralDongJiangRelations} is nevertheless true, by rescaling
  (\lean{OrbC2h.tetrahedralDongJiangRelations}), and the $C_2$ argument survives.
\item \textbf{Library-wide defects that touch the campaign} (Chapter~\ref{ch:L12}): the old \lean{fusionCoeff} is
  identically zero (replaced by \lean{fusionCoeffFHL}), the old bilinear \lean{Unitary} is unsatisfiable in dimension
  $\ge2$ (replaced by \lean{UnitaryVOADL}), \lean{SelfContragredient} is vacuous. None of these is used by the
  classification statements, which use \lean{SelfDual} and the repaired notions.
\item \textbf{A clean axiom audit is not a proof of a hypothesis.} \lean{gannonRayhaunConverse_of_B3} has a clean
  \lean{#print axioms} (Lean's three) but takes orbifold rationality as an argument; until 2026-10-14 that was the only form of the
  converse, and it was correctly reported as conditional. The hypothesis is now proved (Theorem~\ref{l22:thm:B3}).
\item \textbf{Proving the hypothesis changed the axiom footprint.} The conditional converse was free of \lean{VOA.pbw};
  the proved converse \lean{gannonRayhaunConverse_holds} is not. Every categorical step of the proof of orbifold rationality (Steps~2--6)
  works in the logarithmic tensor category, whose construction uses \lean{VOA.pbw} (through the finite-length theorem
  that rests on Step~1); only the algebraic and analytic steps (Schur--Weyl duality, the abstract projection theorem, the
  residue identity, cyclic-orbifold rationality, Step~7 in general) are free of it. So the converse moved from
  ``\lean{VOA.pbw}-free but conditional'' to ``unconditional relative to \lean{VOA.pbw}'', like the forward theorem.
  Since then \lean{VOA.pbw} has become a theorem, so the footprint is again the standard three axioms and
  the converse is unconditional outright.
\item \textbf{The earlier route to orbifold rationality stalled on inputs the final route avoids} (Section~\ref{l22:sec:oldroute}): McRae's
  rationality criterion needs semisimplicity of the Zhu algebras of $V_2^{S_4}$, $V_2^{A_5}$ (not proved independently of
  rationality; it now follows \emph{from} orbifold rationality by Zhu's semisimplicity theorem, which is circular for this
  route; Definition~\ref{l22:def:mcrae}) and
  \lean{CoverDualLitFact}, which needs rigidity, factorizability and unimodularity of the logarithmic category (not
  proved). The orbifold route needs rigidity of the single object $V_2$ only, which Galois descent gives directly.
\end{enumerate}
\end{warning}

\paragraph{Lean defects found while writing this chapter.} All are documentation defects only (wrong or stale
docstrings; no statement is affected):
\begin{enumerate}
\item \leanpath{VOA/\allowbreak Classification/\allowbreak COne/\allowbreak Distinct.lean}: the docstrings of
  \lean{PlusExceptionalResidual} and \lean{ExceptionalResidual} say ``true by GR Thm.~3.1; not proved here''. Both are
  proved (\lean{plusExceptionalResidual}, \lean{exceptionalResidual}, \leanpath{DistinctResiduals.lean}).
\item \leanpath{VOA/\allowbreak Classification/\allowbreak COne/\allowbreak Statement.lean}: the docstrings of
  \lean{GannonRayhaunDistinct}, \lean{GannonRayhaunConverse} and \lean{GannonRayhaunStatement} say ``Not proved
  here''. True of that file, but a reader may take it as a status; all three are now proved
  (\lean{gannonRayhaunDistinct}, \lean{gannonRayhaun_forward}, \lean{gannonRayhaunConverse_holds}).
\item \leanpath{VOA/\allowbreak Classification/\allowbreak COne/\allowbreak McRaeO5Main.lean}: the docstring of
  \lean{TwistedJacobiBridge} is headed ``O5 (b), the analytic bridge (open)'' (``O5 (b)'' is the internal label of Step~5b) and says the conversion to formal identities
  ``is not done''. It is proved (\lean{McRaeO5b.twistedJacobiBridge_holds}).
\item \leanpath{VOA/\allowbreak Classification/\allowbreak COne/\allowbreak McRaeO7Main.lean}: the module docstring and the
  docstring of \lean{TwistedJacobiBridgeAll} call it ``(open)''; it is proved (\lean{twistedJacobiBridgeAll_holds},
  \leanpath{MainConverseFinal.lean}).
\item \leanpath{VOA/\allowbreak Classification/\allowbreak COne/\allowbreak ExceptionalOrbifoldNice.lean}: the docstring
  before \lean{SolvableExceptionalOrbifoldC2Fact} says ``the analytic orbifold theorem [is] not proved here''; true of that
  file, but the theorem (orbifold rationality) is now proved (\lean{finiteV2OrbifoldRationalOfC2}).
\item \leanpath{VOA/\allowbreak Classification/\allowbreak COne/\allowbreak McRaeO5Split.lean}: the module docstring
  describes part (b), \lean{TwistedJacobiBridge}, as ``isolated, not proved''; it is proved
  (\lean{McRaeO5b.twistedJacobiBridge_holds}). (Found by the referee.)
\item \leanpath{VOA/\allowbreak Classification/\allowbreak COne/\allowbreak OrbRat2Reduction.lean}: the module docstring
  says ``all three Props remain OPEN'', and the docstrings of \lean{CubicModelRationalOfC2},
  \lean{InvolutionModelRationalOfC2}, \lean{IcosahedralV2OrbifoldRationalOfC2} begin ``OPEN''. The same file proves
  that orbifold rationality implies each of them, and it is now proved, so all three are consequences of proved theorems (no declaration
  packages them as theorems; they are no longer needed). (Found by the referee.)
\end{enumerate}

%======================================================================
\section{Referee's checklist}\label{l22:sec:checklist}
%======================================================================

\begin{center}\footnotesize
\begin{longtable}{>{\raggedright\arraybackslash}p{0.36\textwidth}>{\raggedright\arraybackslash}p{0.21\textwidth}>{\raggedright\arraybackslash}p{0.21\textwidth}p{0.09\textwidth}}
\toprule
Declaration & Intended meaning & Hypotheses / status & Fidelity \\
\midrule
\lean{GannonRayhaunStatement} & nice, $c=1$ $\Rightarrow$ in the list & definition & Def.~\ref{l22:def:GRS} \\
\lean{Nice}, \lean{PseudoUnitary} & GR Defs.~2.4, 2.6, with rationality read as \lean{AdmissiblyRational} (relation to DLM not settled in Lean) & definitions & Defs.~\ref{l22:def:pu}, \ref{l22:def:nice}, Rem.~\ref{l22:rem:rational} \\
\lean{InCOneList} & the five families & definition; inhabited (Thm.~\ref{l22:thm:subgroups}) & Def.~\ref{l22:def:list} \\
\lean{IsModularData} & modular data of a nice VOA & definition; $S$ pinned & Def.~\ref{l22:def:ismoddata} \\
\lean{gannonRayhaun_of_inputs} & GR Thm.~3.1 from inputs & 7 inputs + 4 gaps; Lean's 3 & Thm.~\ref{l22:thm:conditional} \\
\lean{existsUnique_weight_ne_zero} & GR Prop.~3.12 core & none; Lean's 3 & Thm.~\ref{l22:thm:unique} \\
\lean{tannakian_squeeze} & GR Thm.~2.18 matrix part & none; Lean's 3 & Thm.~\ref{l22:thm:squeeze} \\
\lean{theta_weight_nonneg} & GR Lemma~3.11 & none; Lean's 3 & Thm.~\ref{l22:thm:wnonneg} \\
\lean{gannonRayhaun_forward} & GR Thm.~3.1 forward & \textbf{none}; Lean's 3, via thm \lean{VOA.pbw} & Thm.~\ref{l22:thm:forward} \\
\lean{gannonRayhaunDistinct} & ``exactly one'' & \textbf{none}; Lean's 3 & Thm.~\ref{l22:thm:distinct} \\
\lean{V2m_nice}, \lean{V2mPlus_nice} & lattice and plus VOAs are nice & none; Lean's 3 & Thms.~\ref{l22:thm:V2mnice}, \ref{l22:thm:plusnice} \\
\lean{xuIcosahedralPoissonRelations_holds} & Xu's icosahedral relations & none; Lean's 3 & Thm.~\ref{l22:thm:B245} \\
\lean{gannonRayhaunConverse_of_B3} & GR Prop.~3.2 from orbifold rationality (historical reduction) & conditional on orbifold rationality, \textbf{now proved}; Lean's 3; universe 0 & Thm.~\ref{l22:thm:converseB3} \\
\lean{FiniteV2OrbifoldRationalOfC2Fact} & McRae Cor.~4.23 for $V_2^G$ (orbifold rationality) & definition; \textbf{proved} by \lean{finiteV2Orbifold}\-\lean{RationalOfC2} & Def.~\ref{l22:def:Bfacts} \\
\lean{gannonRayhaunConverse_of_narrowed} & converse from orbifold rationality at $A_4,S_4,A_5$ & all four hypotheses now proved; Lean's 3 & Thm.~\ref{l22:thm:narrowed} \\
\lean{finite_Φ_of_C2Cofinite_pbw} & Step 1: finitely many irreducibles & \lean{hC2}; Lean's 3, via thm \lean{VOA.pbw} & Thm.~\ref{l22:thm:O1} \\
\lean{nonempty_logarithmicBraidedCategory} & Step 2: HLZ braided category & \lean{hC2}; Lean's 3, via thm \lean{VOA.pbw} & Thm.~\ref{l22:thm:O2}; Ch.~\ref{ch:L21} \\
\lean{McRaeO3.haploidCommAlgebra} & Step 3: $V$ haploid commutative algebra & \lean{hfg}, \lean{hC2}, \lean{[Nontrivial V]}, \lean{[Fintype G]}; Lean's 3, via thm \lean{VOA.pbw} & Thm.~\ref{l22:thm:O3} \\
\lean{dlmSchurWeyl_holds}, \lean{V2m_orbifold_finGen} & Step 4a: DLM Schur--Weyl; $V_2$ fin.\ gen.\ over $V_2^G$ & none / \lean{[Finite G]}; Lean's 3 & Def.~\ref{l22:def:dlm} \\
\lean{McRaeO4b.rigidity_dimension} & Step 4b: $V$ rigid self-dual, $\dim=|G|$ & \lean{hfg}, \lean{hC2}, \lean{hW}; Lean's 3, via thm \lean{VOA.pbw} & Thm.~\ref{l22:thm:O4b} \\
\lean{McRaeData.theorem_3_3}, \lean{V2m_theorem_3_3} & Step 4c: McRae Thm.~3.3 (projections) & \lean{McRaeData} / \lean{hC2}; Lean's 3 / Lean's 3, via thm \lean{VOA.pbw} & Thm.~\ref{l22:thm:O4c} \\
\lean{TwistedJacobiBridge} & Jacobi identity on twisted parts & definition; \textbf{proved} & Def.~\ref{l22:def:bridge} \\
\lean{McRaeO5b.twistedJacobiBridge_holds} & Step 5b: the Jacobi bridge & \textbf{none}; Lean's 3, via thm \lean{VOA.pbw} & Thm.~\ref{l22:thm:bridge} \\
\lean{V2m_cyclic_admissiblyRational} & Step 5c: $V_2^{\langle g\rangle}\cong V_{2n^2}$ rational & finite order; Lean's 3 & Thm.~\ref{l22:thm:cyclic} \\
\lean{McRaeO5.V2m_exists_section} & Step 5d: $V_2^G$-linear sections & \lean{hC2}, bridge (proved); Lean's 3, via thm \lean{VOA.pbw} & Thm.~\ref{l22:thm:O5} \\
\lean{McRaeO6.V2m_lemma_4_21} & Step 6: McRae Lemma~4.21 & \lean{hC2}, bridge (proved); Lean's 3, via thm \lean{VOA.pbw} & Thm.~\ref{l22:thm:O6} \\
\lean{McRaeO7.admissiblyRational_of}\-\lean{_surjectionsSplit} & Step 7: split surjections $\Rightarrow$ rational & \lean{hC2}, \lean{hs}; Lean's 3 & Thm.~\ref{l22:thm:O7} \\
\lean{finiteV2OrbifoldRationalOfC2} & rationality of $C_2$-cofinite $V_2^G$ (McRae Cor.~4.23) & \textbf{none}; Lean's 3, via thm \lean{VOA.pbw} & Thm.~\ref{l22:thm:B3} \\
\lean{gannonRayhaunConverse_holds} & GR Prop.~3.2 & \textbf{none}; Lean's 3, via thm \lean{VOA.pbw}; universe 0 & Thm.~\ref{l22:thm:final} \\
\lean{gannonRayhaun_classification} & GR Thm.~3.1 + Prop.~3.2, ``exactly one'', for Lean-niceness (\lean{AdmissiblyRational}) & \textbf{none}; Lean's 3, via thm \lean{VOA.pbw}; stated at universe 0 (\lean{Type 0} carriers) & Thm.~\ref{l22:thm:final}, Rem.~\ref{l22:rem:rational} \\
\bottomrule
\end{longtable}
\end{center}

%======================================================================
\section{Notes and references}\label{l22:sec:notes}
%======================================================================

The theorem is T.~Gannon and B.~C.~Rayhaun, \emph{Classification of rational $c=1$ vertex operator algebras and vertex
operator superalgebras}, arXiv:2609.00122 (2026); the library formalizes the bosonic Theorem~3.1 only (the super case,
their Theorem~4.2, is out of scope by the maintainers' decision). The physical moduli space is in P.~Ginsparg,
\emph{Curiosities at $c=1$}, Nucl.\ Phys.\ B295 (1988) 153--170, and R.~Dijkgraaf, E.~Verlinde and H.~Verlinde,
\emph{$c=1$ conformal field theories on Riemann surfaces}, Comm.\ Math.\ Phys.\ 115 (1988) 649--690. Inputs:
Y.~Zhu, \emph{Modular invariance of characters of vertex operator algebras}, J.\ AMS 9 (1996); C.~Dong, H.~Li and
G.~Mason, \emph{Modular-invariance of trace functions in orbifold theory}, Comm.\ Math.\ Phys.\ 214 (2000);
Y.-Z.~Huang, \emph{Vertex operator algebras, the Verlinde conjecture, and modular tensor categories}, PNAS 102 (2005),
and \emph{Rigidity and modularity of vertex tensor categories}, Commun.\ Contemp.\ Math.\ 10 (2008); C.~Dong, X.~Lin
and S.-H.~Ng, \emph{Congruence property in conformal field theory}, Algebra Number Theory 9 (2015); S.-H.~Ng and
P.~Schauenburg, \emph{Congruence subgroups and generalized Frobenius--Schur indicators}, Comm.\ Math.\ Phys.\ 300 (2010);
A.~Coste and T.~Gannon, \emph{Remarks on Galois symmetry in rational conformal field theories}, Phys.\ Lett.\ B323
(1994); F.~R.~Beyl, \emph{The Schur multiplicator of $\mathrm{SL}(2,\Z/m\Z)$ and the congruence subgroup property},
Math.\ Z.\ 191 (1986); J.-P.~Serre and H.~M.~Stark, \emph{Modular forms of weight 1/2}, LNM 627 (1977);
B.~L.~Feigin and D.~B.~Fuchs, \emph{Verma modules over the Virasoro algebra}, LNM 1060 (1984); M.~Müger, \emph{On the
structure of modular categories}, Proc.\ LMS 87 (2003); P.~Deligne, \emph{Catégories tensorielles}, Mosc.\ Math.\ J.\ 2
(2002); A.~Kirillov Jr.\ and V.~Ostrik, \emph{On a $q$-analog of the McKay correspondence and the ADE classification of
$\mathfrak{sl}_2$ conformal field theories}, Adv.\ Math.\ 171 (2002); Y.-Z.~Huang, A.~Kirillov Jr.\ and J.~Lepowsky,
\emph{Braided tensor categories and extensions of vertex operator algebras}, Comm.\ Math.\ Phys.\ 337 (2015); C.~Dong
and G.~Mason, \emph{Rational vertex operator algebras and the effective central charge}, IMRN 2004; C.~Dong and
K.~Nagatomo, \emph{Automorphism groups and twisted modules for lattice vertex operator algebras} (1999); M.~Miyamoto,
\emph{$C_2$-cofiniteness of cyclic-orbifold models}, Comm.\ Math.\ Phys.\ 335 (2015) (arXiv:1306.5031); S.~Xu,
arXiv:2608.13872 (2026), for the icosahedral relations; R.~McRae, \emph{On the tensor structure of modules for compact
orbifold vertex operator algebras}, Math.\ Z.\ 296 (2020) (arXiv:1910.13226), and \emph{On rationality for
$C_2$-cofinite vertex operator algebras}, arXiv:2108.01898; S.~Carnahan and M.~Miyamoto, \emph{Regularity of fixed-point vertex operator subalgebras},
arXiv:1603.05645; Y.-Z.~Huang, \emph{Cofiniteness conditions, projective
covers and the logarithmic tensor product theory}, J.\ Pure Appl.\ Algebra 213 (2009). For the proof of orbifold rationality
(Section~\ref{l22:sec:B3}): Y.-Z.~Huang, J.~Lepowsky and L.~Zhang, \emph{Logarithmic tensor category theory for generalized
modules for a conformal vertex algebra}, I--VIII, arXiv:1012.4193 and sequels (2010--2011); C.~Dong, H.~Li and G.~Mason,
\emph{Compact automorphism groups of vertex operator algebras}, Internat.\ Math.\ Res.\ Notices 1996, no.~18; S.~U.~Chase,
D.~K.~Harrison and A.~Rosenberg, \emph{Galois theory and Galois cohomology of commutative rings}, Mem.\ AMS 52 (1965);
A.~Kirillov Jr., \emph{Modular categories and orbifold models}, Comm.\ Math.\ Phys.\ 229 (2002); R.~McRae, \emph{Twisted
modules and $G$-equivariantization in logarithmic conformal field theory}, Comm.\ Math.\ Phys.\ 383 (2021)
(arXiv:1810.00747); T.~Creutzig, S.~Kanade and R.~McRae, \emph{Tensor categories for vertex operator superalgebra
extensions}, Mem.\ AMS (2024); I.~Frenkel, J.~Lepowsky and A.~Meurman, \emph{Vertex operator algebras and the Monster},
Academic Press (1988), and I.~Frenkel, Y.-Z.~Huang and J.~Lepowsky, \emph{On axiomatic approaches to vertex operator
algebras and modules}, Mem.\ AMS 494 (1993), for the derivation of the Jacobi identity from rationality, commutativity
and associativity used in Step~5b. The project's working notes are
\leanpath{Notes/\allowbreak C1\_CLASSIFICATION\_ROADMAP.md} and \leanpath{Notes/\allowbreak C1\_DEPENDENCY\_TREE.md} (last
section: the orbifold route to rationality).

\paragraph{Verification record.} Quotations checked verbatim by \texttt{scripts/textbook/sync.py}; aggregate build and
axiom audit as recorded in the status of Theorem~\ref{l22:thm:final} (2026-10-14; re-probed at an earlier commit, not re-run for this edition).
Release audit 2026-10-27: see Notes/Textbook/referee/RELEASE\_AUDIT\_L22.md. (Verdict: announceable with caveats --- the
rationality convention, \lean{Type 0}, and the aggregate build to be re-confirmed on the release commit; the three
headlines have no hypotheses and report Lean's three standard axioms; \lean{VOA.pbw} is a theorem.)

\endgroup
